\documentclass[11pt]{article}
\usepackage[T1]{fontenc}
\usepackage[utf8]{inputenc}
\usepackage{lmodern,amsmath,amssymb,amsthm,mathtools}
\usepackage[margin=1in]{geometry}
\usepackage{xurl}
\usepackage[hidelinks]{hyperref}
\usepackage{microtype}
\hypersetup{pdftitle={Worst-case transport and light cones for quasi-periodic operators},pdfauthor={Seth Douglas and Nidhal Mghirbi}}
\newcommand{\C}{\mathbb C}
\newcommand{\N}{\mathbb N}
\newcommand{\Z}{\mathbb Z}
\newcommand{\Q}{\mathbb Q}
\newcommand{\R}{\mathbb R}
\newcommand{\T}{\mathbb T}
\newcommand{\un}{\mathbf 1}
\newcommand{\tr}{\operatorname{tr}}
\newcommand{\ad}{\operatorname{ad}}
\newcommand{\dist}{\operatorname{dist}}
\newcommand{\supp}{\operatorname{supp}}
\newcommand{\Lip}{\operatorname{Lip}}
\newcommand{\Ga}{G_\alpha}
\newcommand{\au}{\alpha_u^+}
\newcommand{\aunif}{\alpha_u^{\mathrm{unif}}}
\newcommand{\eps}{\varepsilon}
\newcommand{\om}{\omega}
\newcommand{\omt}{\tilde\omega}
\newcommand{\omh}{\hat\omega}
\newcommand{\FU}{F_{\mathrm U}}
\newtheorem{theorem}{Theorem}[section]
\newtheorem{lemma}[theorem]{Lemma}
\newtheorem{proposition}[theorem]{Proposition}
\newtheorem{corollary}[theorem]{Corollary}
\newtheorem{conjecture}[theorem]{Conjecture}
\newtheorem{maintheorem}{Theorem}
\renewcommand{\themaintheorem}{\Alph{maintheorem}}
\newtheorem{maincorollary}[maintheorem]{Corollary}
\theoremstyle{definition}
\newtheorem{definition}[theorem]{Definition}
\newtheorem{remark}[theorem]{Remark}
\newtheorem{example}[theorem]{Example}
\setlength{\emergencystretch}{2em}
\raggedbottom
\title{Worst-case transport and light cones\\ for quasi-periodic operators\thanks{The authors developed this work with AI models from more than one developer and take responsibility for every claim. How the work was done is described at the end of the paper.}}
\author{Seth Douglas \and Nidhal Mghirbi}
\date{Version 1.0.0, October 2026}
\begin{document}
\maketitle
\begin{abstract}
We study the fastest transport exponent $\alpha_u^+$ of quasi-periodic Schr\"odinger operators $(Hu)(n)=u(n+1)+u(n-1)+v(\theta+n\alpha)u(n)$, with $v$ real analytic and positive Lyapunov exponent on the spectrum, at the worst phase rather than a typical one. Let $\beta(\alpha)=\limsup_nq_n^{-1}\log q_{n+1}$. For $0<\beta<\infty$ we prove, uniformly in the phase, that $\alpha_u^+$ is at most the maximum over the spectrum of $1-\ell/(2\omega\beta)$ when $\beta\ge\ell$, with an explicit formula for $\beta<\ell$, where $\ell$ is the Lyapunov exponent and $\omega$ the acceleration. So finite $\beta$ keeps $\alpha_u^+$ below $1$ at every phase. The factor $2\omega$ is attained, by $2\lambda\cos2\pi sx$. For the almost Mathieu operator the phase $1/4$ has exponent at least $1-\log\lambda/(2\beta_e)$, where $\beta_e$ is the limsup along even denominators; when $\beta_e=\beta$ is finite and at least $\log\lambda$ this is the exact worst exponent, and it holds on a dense $G_\delta$ set of phases, while almost every phase has $1-\log\lambda/\beta$. For potentials with $v(x+1/2)=-v(x)$, a half-period sign flip closes the central gap of every even-denominator approximant and gives the uniform lower bound $1-\ell(0)/(2\beta_e)$. In the Aubry--Andr\'e XY chain, for even sources and the receivers we specify, at almost every phase the light cone at a fixed site follows the typical law, and when $\beta_e=\beta$ is finite and at least $\log\lambda$ the light cone uniform over sites follows the worst-case law; for every supercritical analytic field the uniform light cone obeys the all-phase upper bound; for an explicit computable chain the two exponents differ and both have Turing degree $\mathbf 0''$. Finally, the Lyapunov exponent of a computable analytic potential is computable at every computable irrational frequency, uniformly in energy, and at one explicit computable frequency and the phase $1/4$ the transport exponent has degree exactly $\mathbf 0''$ for every computable even potential in an explicit open neighbourhood of $4\cos2\pi x$.
\end{abstract}

{\small\tableofcontents}

\section{Introduction}\label{sec:intro}

Consider a one-frequency quasi-periodic Schr\"odinger operator on $\ell^2(\Z)$,
\begin{equation}\label{eq:qp}
(H_{v,\alpha,\theta}u)(n)=u(n+1)+u(n-1)+v(\theta+n\alpha)\,u(n),
\end{equation}
with a real-analytic sampling function $v$ on $\T=\R/\Z$, an irrational frequency $\alpha$ and a phase $\theta\in\T$. The almost Mathieu operator, $v(x)=2\lambda\cos2\pi x$, is the Aubry--Andr\'e model of physics~\cite{AubryAndre1980}. Let $p_n/q_n$ be the continued-fraction convergents of $\alpha$ and
\begin{equation}\label{eq:beta}
\beta(\alpha)=\limsup_{n\to\infty}\frac{\log q_{n+1}}{q_n}\in[0,\infty].
\end{equation}
When the Lyapunov exponent $L$ is positive on the spectrum, transport is a competition between the localization strength $L$ and the Liouville strength $\beta(\alpha)$. The companion papers~\cite{DMamo,DMuniv} determine the fastest transport exponent $\au$ at almost every phase: it equals $\max\{0,1-\ell_*/\beta(\alpha)\}$, where $\ell_*$ is the minimum of $L$ on the spectrum, for the almost Mathieu operator and for every supercritical analytic potential of type~I.

This paper asks what happens at the worst phase. The question has a physical reading. Moving the launch site of a packet by $x$ sites moves the phase by $x\alpha$, so the worst phase governs transport uniformly over launch sites. Through the Jordan--Wigner transformation it bounds the uniform light cone of the quasi-periodic XY spin chain, and for the Aubry--Andr\'e chain, at almost every phase, the typical law governs the light cone seen from a fixed site. Three points summarize the answers.
\begin{itemize}
\item At the worst phase the localization strength is divided by at most twice the acceleration: when $\max_\Sigma\ell\le\beta(\alpha)<\infty$, transport is never faster than $1-\ell/(2\om\beta)$, maximized over the spectrum, for every analytic potential with $L>0$ on the spectrum and uniformly in the phase. The factor $2\om$ is attained.
\item The worst value is generic in the sense of Baire: for the almost Mathieu operator, when $\beta_e(\alpha)=\beta(\alpha)\in[\log\lambda,\infty)$, the exponent is $1-\log\lambda/(2\beta)$ on a residual set of phases and $1-\log\lambda/\beta$ at almost every phase. In general the typical value is $\max\{0,1-\log\lambda/\beta\}$.
\item In the Aubry--Andr\'e XY chain, for the observables specified below, at almost every phase the fixed-site light cone follows the typical law, and when $\beta_e=\beta$ is finite and at least $\log\lambda$ the uniform light cone follows the worst-case law; for every supercritical analytic field the uniform light cone obeys the all-phase upper law. For an explicit computable chain both exponents are different reals of Turing degree $\mathbf 0''$. The Lyapunov exponent itself is computable, so the uncomputability sits entirely in the frequency.
\end{itemize}

\subsection{Transport exponents}

Let $\psi_t=e^{-itH}\delta_0$, $a(j,T)=\frac2T\int_0^\infty e^{-2t/T}|\langle\delta_j,\psi_t\rangle|^2\,dt$, $P(N,T)=\sum_{|j|>N}a(j,T)$ and $M_p(T)=\sum_j|j|^pa(j,T)$. The upper transport exponents are $\beta^+(p)=\limsup_{T\to\infty}\log M_p(T)/(p\log T)$, and $\au=\sup_{p>0}\beta^+(p)$ is the fastest one, as in~\cite{DT2007,DMamo,DMuniv}. A tail $P(T^a,T)$ decaying faster than every power forces $\au\le a$, and a tail of size $T^{-C}$ at distance $T^a$ forces $\au\ge a$. We write $\au(\theta)$ when the phase matters, and
\[
\aunif=\inf\Bigl\{a\ge0:\ \sup_\theta P_\theta(T^a,T)=O(T^{-m})\text{ for every }m>0\Bigr\}
\]
for the uniform exponent. It bounds $\sup_\theta\au(\theta)$, and it is also the exponent of all packets of one operator, uniformly in the launch site (Section~\ref{sec:prelim}). Let
\begin{equation}\label{eq:Ga}
\Ga=\bigl\{\theta\in\T:\ |\sin(2\pi q_n\theta)|\ge q_n^{-2}\text{ for all sufficiently large }n\bigr\},
\end{equation}
a set of full measure~\cite[Lemma~2.1]{DMamo}. Let $\delta_n=q_n\alpha-p_n$, $b_n=q_n^{-1}\log|\delta_n|^{-1}$, and let $\beta_e(\alpha)$ and $\beta_o(\alpha)$ be the limsups of $b_n$ along even and along odd $q_n$, with value $0$ along a finite set; then $\beta(\alpha)=\max\{\beta_e,\beta_o\}$.

\subsection{The all-phase law}

Write $\ell(E)$ for the Lyapunov exponent at energy $E$ and $\Sigma$ for the spectrum. The acceleration $\om(E)$, from Avila's global theory~\cite{Avila2015}, is the right derivative at $y=0$ of $L(E,y)/(2\pi)$, where $L(E,y)$ is the Lyapunov exponent with phase complexified by $iy$; it is a positive integer at every spectral energy where $\ell>0$. The operator is supercritical if $\ell>0$ on $\Sigma$, and of type~I if moreover $\om=1$ on $\Sigma$~\cite{GeJitomirskaya2024}. For $\ell>0$, $w\ge1/2$ and $b\ge0$ let
\[
f(b;\ell,w)=1-\frac{\ell}{2wb}\quad(b\ge\ell),\qquad f(b;\ell,w)=\frac{(2w-1)b}{(2w-1)b+\ell}\quad(b<\ell).
\]

\begin{maintheorem}[All-phase upper law]\label{thm:A}
Let $v$ be real analytic with $L>0$ on the spectrum, and let $0<\beta(\alpha)<\infty$. Put $\FU(\beta)=\max_{E\in\Sigma}f(\beta;\ell(E),\om(E))$. For every $a>\FU(\beta)$ and every $m>0$,
\[
\sup_{\theta\in\T}P_\theta(T^a,T)=O(T^{-m})\qquad(T\to\infty).
\]
Hence, at every phase, $\au(\theta)\le\aunif\le\FU(\beta)<1$. If $\beta\ge\max_\Sigma\ell$, then
\[
\FU(\beta)=1-\frac1{2\beta}\min_{E\in\Sigma}\frac{\ell(E)}{\om(E)}\le1-\frac{\ell_*}{2\om^+\beta},\qquad\om^+=\max_\Sigma\om .
\]
\end{maintheorem}

For type~I potentials the bound is $1-\ell_*/(2\beta)$ when $\beta\ge\ell_*$, and $\beta/(\beta+\ell_*)$ when $\beta<\ell_*$ (Corollary~\ref{cor:typeI}).

\begin{remark}\label{rem:A}
At $\beta(\alpha)=0$ the same argument gives $0$; this endpoint is known~\cite{DT2007},~\cite[Remark~1.12]{JitomirskayaMavi2017},~\cite[Corollary~5.10]{DMuniv}. At $\beta(\alpha)=\infty$, Jitomirskaya and Zhang proved $\au=1$ at every phase~\cite[Theorem~3(2)]{JZ2022}. With Theorem~\ref{thm:A}: $\au=1$ at some, equivalently at every, phase if and only if $\beta(\alpha)=\infty$. The lower bound $\au(\theta)\ge\max\{0,1-\ell_*/\beta\}$ holds at every phase~\cite[Theorem~4.1]{DMuniv}, so Theorem~\ref{thm:A} brackets the exponent at every phase.
\end{remark}

\subsection{Sharpness}

\begin{maintheorem}[Sharpness]\label{thm:B}
\begin{enumerate}
\item Let $\lambda>1$ and $L=\log\lambda$. For the almost Mathieu operator, $\au(H_{\lambda,\alpha,1/4})\ge\max\{0,1-L/(2\beta_e(\alpha))\}$ for every irrational $\alpha$. If $\beta_e(\alpha)=\beta(\alpha)\in[L,\infty)$, then
\[
\aunif=\sup_\theta\au(\theta)=\au(1/4)=1-\frac{L}{2\beta(\alpha)} .
\]
If $\beta_e(\alpha)\ge L$, then $\au(1/4)=\max\{0,1-L/\beta_o(\alpha),1-L/(2\beta_e(\alpha))\}$.
\item For $v(x)=2\lambda\cos2\pi sx$ with $\lambda>1$ and an integer $s\ge2$, the acceleration is $s$ on $\Sigma$, and for every $\beta_*>\log\lambda$ there is an $\alpha$ with $\beta(\alpha)=\beta_*$ and $\aunif=\sup_\theta\au(\theta)=1-\ell_*/(2\om^+\beta(\alpha))$. So the factor $2\om^+$ in Theorem~\ref{thm:A} cannot be improved in the class of analytic potentials.
\item Let $v$ be real analytic with $L>0$ on $\Sigma$ and $v(x+1/2)=-v(x)$. Then $0\in\Sigma$, every even-denominator periodic approximant has a closed gap at $E=0$ at some phase, and $\aunif\ge\max\{0,1-\ell(0)/(2\beta_e(\alpha))\}$.
\end{enumerate}
\end{maintheorem}

\begin{remark}\label{rem:B}
At almost every phase the exponent of the almost Mathieu operator is $\max\{0,1-L/\beta\}$~\cite[Theorem~A]{DMamo}. So when $\beta_e=\beta\in[L,\infty)$ the worst phase sees exactly half the typical localization strength, and when $\beta_e=\beta\in(L/2,\infty)$ the phase $1/4$ is strictly faster than almost every phase, including in the window $e^{\beta}<\lambda<e^{2\beta}$ where Avila, You and Zhou proved Anderson localization at almost every phase~\cite{AYZ2017}. For the almost Mathieu operator the closed central gap of even-denominator approximants is classical~\cite{BellissardSimon1982}, and the value of the Chambers polynomial at $0$ that expresses it is due to Lamoureux and Mingo~\cite{LamoureuxMingo2007}.
\end{remark}

For odd-harmonic type~I potentials, Theorem~\ref{thm:B}(3) and Theorem~\ref{thm:A} give $1-\ell(0)/(2\beta)\le\aunif\le1-\ell_*/(2\beta)$ when $\beta_e=\beta$ is finite and at least $\ell_*$. When moreover $\beta\ge\ell(0)$, the uniform exponent of packets filtered ever more narrowly near the closed gap at $E=0$ converges to $1-\ell(0)/(2\beta)$ (Corollary~\ref{cor:window}); this is a statement about the uniform exponent in the limit of narrow filters, not about each phase or each fixed filter.

\subsection{Measure versus category}

\begin{maintheorem}[Category]\label{thm:C}
Let $\lambda>1$ and $\beta_e(\alpha)=\beta(\alpha)\in[\log\lambda,\infty)$. Then $\au(H_{\lambda,\alpha,\theta})=1-\log\lambda/(2\beta(\alpha))$ for every $\theta$ in a dense $G_\delta$ set.
\end{maintheorem}

Since $\au=1-\log\lambda/\beta$ at almost every phase~\cite[Theorem~A]{DMamo}, the exponent takes one value on a set of full measure and a strictly larger value on a residual set. The lower bound on a residual set holds whenever $\beta_e>\log\lambda/2$, and a residual lower bound $1-\ell(0)/(2\beta_e)$ holds for every supercritical odd-harmonic analytic potential (Corollary~\ref{cor:G2}). The Baire argument behind it is standard; the content lies in lower-bound mechanisms that depend on the phase only through an open set of local phases.

\subsection{Light cones}

Consider the XY chain $\sum_n(X_nX_{n+1}+Y_nY_{n+1})+\sum_nh_nZ_n$ with $h_n=v(\theta+n\alpha)$. For an even observable $A$ near a site $x$ and receivers $B$ beyond distance $T^a$ from $x$, either even observables on the two sides, such as parity strings, or arbitrary observables on a half-line, let $\Gamma_x(\theta)$ be the least $a$ for which $\|[\tau_t(A),B]\|$ decays faster than every power of $T$ uniformly for $|t|\le T$, and $\Gamma_{\mathrm{unif}}(\theta)$ the same uniformly over $x$ (Definition~\ref{def:LR}).

\begin{maintheorem}[Light cones]\label{thm:D}
\begin{enumerate}
\item For the Aubry--Andr\'e chain, $h_n=2\lambda\cos2\pi(\theta+n\alpha)$ with $\lambda>1$, and for almost every $\theta$, the light-cone exponent at every site is $\Gamma_x(\theta)=\max\{0,1-\log\lambda/\beta(\alpha)\}$.
\item For every real-analytic $v$ with $L>0$ on $\Sigma$ and $\beta(\alpha)<\infty$, $\Gamma_{\mathrm{unif}}(\theta)\le\FU(\beta(\alpha))$ at every phase. For the Aubry--Andr\'e chain with $\beta_e(\alpha)=\beta(\alpha)\in[\log\lambda,\infty)$, $\Gamma_{\mathrm{unif}}(\theta)=1-\log\lambda/(2\beta(\alpha))$ at every phase.
\item There are a computable $\alpha$ and a computable phase $\theta_*$ such that the chain with $h_n=4\cos2\pi(\theta_*+n\alpha)$ has $\Gamma_{\mathrm{unif}}=1-1/(2s_*)$ and $\Gamma_x=1-1/s_*$ at every site, where $s_*\in[23/12,2]$ has Turing degree $\mathbf 0''$; both exponents have degree $\mathbf 0''$. In the same chain at the phase $1/4$, every site has exponent $1-1/(2s_*)$.
\end{enumerate}
\end{maintheorem}

For the pair $Z_0$, $P_R$, an order-one commutator is equivalent to a one-body mass in $R$ bounded away from both $0$ and $1$: $\|[\tau_t(Z_0),P_R]\|=4\sqrt{p(1-p)}$, where $p$ is the one-body mass in $R$ at time $2t$ (Proposition~\ref{prop:comm}). The bounds concern even sources and the stated receivers; they are not conventional exponential Lieb--Robinson bounds for arbitrary observables.

\begin{remark}\label{rem:D}
Reductions of XY dynamics to one-body dynamics go back to Lieb, Schultz and Mattis~\cite{LSM1961}, and were used by Hamza, Sims and Stolz for disordered chains~\cite{HSS2012}. Damanik, Lemm, Lukic and Yessen identified anomalous Lieb--Robinson exponents of the Fibonacci XY chain with upper transport exponents~\cite{DLLY2016}. The infinite-temperature OTOC for one receiver site is contained in the formula of Riddell and S{\o}rensen~\cite{RiddellSorensen2020}.
\end{remark}

\subsection{Computability}

\begin{maintheorem}[Computability]\label{thm:E}
\begin{enumerate}
\item For every computable real-analytic $v$, the Lyapunov exponent is a computable function of the frequency and the energy on $(\R\setminus\Q)\times\R$. At every computable irrational frequency, $\ell$, $\ell_*$, the density of states and the spectrum are computable, with no hypothesis of supercriticality.
\item Let $\alpha_*$ be the explicit computable frequency of~\cite[Theorem~D(3)]{DMamo}. For every even analytic $v$ with computable coefficients and $\sup_{|\operatorname{Im}z|\le1/4}|v(z)-4\cos2\pi z|<1/4$, the exponent $\au(H_{v,\alpha_*,1/4})$ and every moment exponent equal $1-\ell_*(v)/\beta(\alpha_*)\in[0.3199,0.6924]$, which has Turing degree exactly $\mathbf 0''$; the value depends continuously on $v$.
\end{enumerate}
\end{maintheorem}

Part~1 is an algorithmic consequence of Avila's global theory; we make no priority claim for it, and computing spectra is not new (Remark~\ref{rem:compprior}). Paper A gave an exponent of degree $\mathbf 0''$ that is robust along a coupling interval~\cite[Theorem~D]{DMamo}; part~2 extends this to the computable members of an explicit neighbourhood of $4\cos2\pi x$ in the even analytic potentials, open relative to that class. Computable coefficients are needed: a noncomputable coupling in the same set gives the computable value $49/100$ (Remark~\ref{rem:noncomp}). Section~\ref{sec:computability} also shows that, for these transport laws at the prescribed good phases, the frequency is the only fragile datum, and that in the supercritical analytic class uncomputability is carried by $\limsup$-type quantities.

\subsection{Relation to earlier work}

For phases outside the full-measure sets of the companion papers, little was known at positive finite $\beta(\alpha)$. Damanik and Tcheremchantsev proved vanishing upper exponents at every phase when $\beta(\alpha)=0$, for trigonometric polynomial potentials~\cite[Theorems~4--5]{DT2007}, with methods that also control individual times~\cite{DT2008}. Jitomirskaya and Mavi proved arithmetic upper bounds that reach $\beta=0$ for analytic potentials~\cite[Theorem~1.7 and Remark~1.12]{JitomirskayaMavi2017}. Han and Jitomirskaya proved upper bounds under a strongly bounded discrepancy condition~\cite{HanJitomirskaya2019}, and Bradshaw, de Jong, Liu, Wang, Wang and Yang proved phase-uniform subpower bounds for weakly Liouville frequencies, which have $\beta=0$~\cite{BdJLWWY2025}. Jitomirskaya, Liu and Tcheremchantsev bounded fractal dimensions of spectral measures near the transition for $\alpha$-Diophantine phases~\cite{JLT2025}. Jitomirskaya and Zhang settled $\beta=\infty$~\cite{JZ2022}. In the literature we checked we did not find an all-phase upper law at positive finite $\beta$; this is a statement about our search, not a priority claim. The structural input is Avila's global theory~\cite{Avila2015}. The connection between the acceleration and zeros of finite-volume determinants is due to Han and Schlag~\cite{HanSchlag2022}, for Diophantine frequencies (phase restrictions enter only their application to localization), and the acceleration counts the vanishing dual Lyapunov exponents by Ge, Jitomirskaya, You and Zhou~\cite{GJYZ2023}; we prove the finite-scale count we need directly. The interpolation at every phase rests on the product lemma of Avila and Jitomirskaya~\cite[Lemma~9.7]{AJ2009}, used in the hierarchical analysis of Jitomirskaya and Liu~\cite{JL2018}. Reflection resonances of the phase were developed by Jitomirskaya and Liu~\cite{JLphase} and Liu~\cite{Liu2022}, and Jitomirskaya, Liu and Mi gave a sharp criterion for semi-uniform dynamical localization, with almost Mathieu examples that are dynamically localized but not semi-uniformly~\cite{JLM2024}. The distinction between sitewise and uniform behaviour is therefore not new; the exact power laws are.

\subsection{How the proofs work}

\emph{The upper bound} (Sections~\ref{sec:zeros}--\ref{sec:upperlaw}) combines two single-scale estimates through the covering argument of~\cite{DMamo,DMuniv}. At a resonant scale $q=q_n$, the rational discriminant $t_q(x,E)$ is, as a function of the local phase $x$, close to a trigonometric polynomial of bounded degree. A Jensen identity bounds the number of its zeros in a strip of width proportional to $q$ by about $2\om(E)$ (Lemma~\ref{lem:zerocount}). So the set where the discriminant is below an order-one threshold is covered by at most $2\om$ arcs, each of length about $e^{-\ell q/(2\om)}$ (Lemma~\ref{lem:sublevel}); the $q$-step resolvent identity of~\cite{DMuniv} then confines the packet within about $e^{(b-\min\{\ell,b\}/(2\om))q}$ sites at every phase (Proposition~\ref{prop:Comega}). For the speed, determinants of intervals of length $k$ are, up to an exponentially smaller error, trigonometric polynomials of degree about $\om k$ in the starting phase (Lemma~\ref{lem:RN}), and $q$ consecutive orbit points interpolate polynomials of degree $q/2$ at every phase (Lemma~\ref{lem:TI}). So every window of $2q$ sites contains a regular interval of length about $q/(2\om)$, and a flux argument bounds the speed by $e^{-\ell q/(2\om)}$ (Theorem~\ref{thm:V}). Optimizing the two radii over the time gives $f$.

\emph{Sharpness} (Section~\ref{sec:sharp}). A half-period sign flip closes the central gap of every even-denominator approximant of an odd-harmonic potential (Proposition~\ref{prop:G}). Next to a closed gap the band is wide, its discriminant has derivative about $e^{\ell q/2}$, and an elliptic run of the Weyl solution from~\cite{DMamo} carries a polynomial tail to distance $e^{(b-\ell/2)q}$.

\emph{Computability} (Section~\ref{sec:computability}). Where the Lyapunov exponent is positive, quantization makes the complexified cocycle uniformly hyperbolic at small positive heights. There a finite certificate bounds the Lyapunov exponent from below, and convexity carries the bound to height zero exactly. Where it is zero, nonnegativity and finite-scale upper bounds suffice.

\subsection{Organization}

Section~\ref{sec:prelim} fixes notation and lists the imported results. Sections~\ref{sec:zeros}--\ref{sec:upperlaw} prove Theorem~\ref{thm:A}. Section~\ref{sec:amoproof} gives a second proof for the almost Mathieu operator. Section~\ref{sec:sharp} proves Theorem~\ref{thm:B}, Section~\ref{sec:category} Theorem~\ref{thm:C}, Section~\ref{sec:xy} Theorem~\ref{thm:D}, and Section~\ref{sec:computability} Theorem~\ref{thm:E}. Section~\ref{sec:discussion} discusses open problems.

\section{Setting and imported results}\label{sec:prelim}

\subsection{Operators, Lyapunov exponents and acceleration}

Throughout, $v:\T\to\R$ is real analytic. It extends holomorphically to a strip $|\operatorname{Im}x|<\delta_v$, and we fix $0<\delta_1<\delta_v$. For an irrational $\alpha$ and a phase $\theta\in\T$ we write $H_\theta=H_{v,\alpha,\theta}$ for the operator~\eqref{eq:qp}, and $H_{\lambda,\alpha,\theta}$ for the almost Mathieu operator, $v(x)=2\lambda\cos2\pi x$; for it we write $L=\log\lambda$. The spectrum $\Sigma$ of $H_\theta$ does not depend on $\theta$. We put
\[
J=[-\|v\|_\infty-3,\ \|v\|_\infty+3],
\]
an interval that contains $\Sigma$ with room to spare.

For $E\in\C$ and $x\in\C$ with $|\operatorname{Im}x|<\delta_v$ let $T(E,x)=\bigl(\begin{smallmatrix}E-v(x)&-1\\1&0\end{smallmatrix}\bigr)$ and $A_r(E,x)=T(E,x+(r-1)\alpha)\cdots T(E,x)$. For a rational frequency $p/q$ we write $A^{(p/q)}_r$ for the same products with $\alpha$ replaced by $p/q$. We write $\ell(E)$ for the Lyapunov exponent at $E\in\C$, and, for real $E$ and $|y|<\delta_v$, $L(E,y)$ for the Lyapunov exponent of the cocycle $x\mapsto T(E,x+iy)$ over the rotation by $\alpha$; so $L(E,0)=\ell(E)$. The function $L(E,\cdot)$ is even and convex~\cite[Section~2.1]{Avila2015}. The \emph{acceleration} is
\[
\om(E)=\lim_{y\downarrow0}\frac{L(E,y)-\ell(E)}{2\pi y}.
\]
We use three facts from Avila's global theory~\cite{Avila2015}, in the form recorded in~\cite[Section~5.1]{DMuniv}.
\begin{itemize}
\item[(A1)] \emph{Quantization}~\cite[Theorem~5]{Avila2015}. The acceleration of an analytic $SL(2,\C)$ cocycle over an irrational rotation is an integer. Hence $L(E,\cdot)$ is convex and piecewise affine with slopes in $2\pi\Z$.
\item[(A2)] \emph{Regularity}~\cite[Definition~3 and Theorem~6]{Avila2015}. A cocycle is regular if its complexified Lyapunov exponent is affine for heights near zero. If its Lyapunov exponent is positive, it is regular if and only if it is uniformly hyperbolic.
\item[(A3)] \emph{Uniform hyperbolicity}~\cite[Section~2.1]{Avila2015}. A real energy lies outside $\Sigma$ if and only if its cocycle is uniformly hyperbolic. Uniform hyperbolicity is an open condition, and the Lyapunov exponent is continuous on the uniformly hyperbolic cocycles.
\end{itemize}
The operator is \emph{supercritical} if $\ell>0$ on $\Sigma$. Then $\ell_*=\min_\Sigma\ell>0$, and we put $\ell^+=\max_\Sigma\ell$. At a supercritical spectral energy the acceleration is an integer $\om(E)\ge1$: if it were $0$, evenness would make $L(E,\cdot)$ constant near $0$, the cocycle would be regular with positive Lyapunov exponent, hence uniformly hyperbolic by (A2), and $E\notin\Sigma$ by (A3)~\cite[proof of Lemma~5.2]{DMuniv}. The operator is \emph{supercritical and type~I} if moreover $\om=1$ on $\Sigma$~\cite[Definition~2.1]{GeJitomirskaya2024},~\cite[Definition~5.1]{DMuniv}. The almost Mathieu operator with $\lambda>1$ is supercritical and type~I, with $\ell=L$ and $\om=1$ on $\Sigma$~\cite[Section~5.1]{DMuniv}.

We also use the continuity of $E\mapsto L(E,y)$ for each fixed $|y|<\delta_v$, which holds for analytic $SL(2,\C)$ cocycles at every irrational frequency~\cite{JKS2009,AJS2014}; for $y=0$ it is~\cite[Theorem~1]{BourgainJitomirskaya2002}. Since $L(E,\cdot)$ is convex, pointwise convergence on an open interval is locally uniform, so $(E,y)\mapsto L(E,y)$ is jointly continuous on $\R\times(-\delta_v,\delta_v)$. Finally, by the Thouless formula~\cite{CraigSimon1983},
\begin{equation}\label{eq:thouless}
\ell(E+i\eta)\ge\ell(E)\qquad(E\in\R,\ \eta\ge0),
\end{equation}
as in~\cite[proof of Lemma~5.8]{DMuniv}. We use~\eqref{eq:thouless} only at energies $E+i\eta$ with the same real part.

\subsection{Continued fractions}

Let $p_n/q_n$ be the convergents of $\alpha$ and put
\[
\delta_n=q_n\alpha-p_n,\qquad D_n=|\delta_n|^{-1},\qquad b_n=\frac{\log D_n}{q_n},\qquad \beta(\alpha)=\limsup_{n\to\infty}b_n=\limsup_{n\to\infty}\frac{\log q_{n+1}}{q_n}.
\]
The two limsups agree because $q_{n+1}<D_n<2q_{n+1}$~\cite[eq.~(5)]{DMuniv}. At a fixed scale we drop the index and write $q,p,\delta,D,b$. We also use the limsups along even and odd denominators,
\[
\beta_e(\alpha)=\limsup_{n:\,q_n\text{ even}}b_n,\qquad\beta_o(\alpha)=\limsup_{n:\,q_n\text{ odd}}b_n,
\]
with the convention that a limsup over a finite set is $0$. Thus $\beta=\max\{\beta_e,\beta_o\}$.

\subsection{Transport exponents}

For a vector $\psi\in\ell^2(\Z)$ let $\psi_t=e^{-itH}\psi$ and
\begin{equation}\label{eq:abel}
\begin{gathered}
a_\psi(j,T)=\frac2T\int_0^\infty e^{-2t/T}|\langle\delta_j,\psi_t\rangle|^2\,dt,\\
P_\psi(N,T)=\sum_{|j|>N}a_\psi(j,T),\qquad M_{p,\psi}(T)=\sum_j|j|^pa_\psi(j,T).
\end{gathered}
\end{equation}
For $\psi\ne0$ the upper transport exponents are
\begin{equation}\label{eq:exponents}
\beta^+_\psi(p)=\limsup_{T\to\infty}\frac{\log M_{p,\psi}(T)}{p\log T},\qquad\alpha^+_{u,\psi}=\sup_{p>0}\beta^+_\psi(p).
\end{equation}
They do not change when $\psi$ is multiplied by a nonzero constant. When $\psi=\delta_0$ we drop the index, and we write $P_\theta(N,T)$ and $\au(\theta)=\au(H_{v,\alpha,\theta})$ when the phase matters. The fastest exponent $\au$ is that of Damanik and Tcheremchantsev~\cite{DT2007}. Always $0\le\au\le1$. We use two facts from~\cite{DMuniv}: if $0<a<1$ and $P_\psi(T^a,T)=O(T^{-m})$ for every $m$, for a unit vector $\psi$ with finite moments, then $\alpha^+_{u,\psi}\le a$~\cite[Lemma~2.3]{DMuniv}; and the Parseval formula~\cite[Lemma~2.1]{DMuniv}
\begin{equation}\label{eq:parseval}
P_\psi(N,T)=\frac1{\pi T}\int_\R\bigl\|\un_{\{|j|>N\}}(H-E-i/T)^{-1}\psi\bigr\|^2\,dE,\qquad\frac1{\pi T}\int_\R\bigl\|(H-E-i/T)^{-1}\psi\bigr\|^2\,dE=\|\psi\|^2.
\end{equation}
Conversely, $M_{p}(T)\ge N^pP(N,T)$, so a tail $P(N_j,T_j)\ge T_j^{-C}$ along $T_j\to\infty$ with $\log N_j/\log T_j\to a$ gives $\au\ge a$.

\begin{definition}[Uniform exponent]\label{def:unif}
For $v$ and $\alpha$ fixed, let
\[
\aunif=\inf\Bigl\{a\ge0:\ \sup_{\theta\in\T}P_\theta(T^a,T)=O(T^{-m})\text{ as }T\to\infty,\text{ for every }m>0\Bigr\}.
\]
\end{definition}

The set in the definition contains every $a>1$, because hopping has range one, so $\aunif\le1$. If $\sup_\theta P_\theta(T^a,T)=O(T^{-m})$ for every $m$, then $\au(\theta)\le a$ for every $\theta$. Hence
\begin{equation}\label{eq:unifsup}
\sup_\theta\au(\theta)\le\aunif.
\end{equation}
Translation by $x$ sites conjugates $H_\theta$ to $H_{\theta+x\alpha}$, and for fixed $N,T$ the map $\theta\mapsto P_\theta(N,T)$ is continuous. Since the orbit $\theta+x\alpha$ is dense, $\sup_{x\in\Z}P_{\theta+x\alpha}(N,T)=\sup_\varphi P_\varphi(N,T)$ for every $\theta$. So $\aunif$ is also the exponent of the packets launched from all sites of one fixed operator $H_\theta$, uniformly in the launch site, for any $\theta$. We do not know whether equality holds in~\eqref{eq:unifsup} in general.

For $0<\beta<\infty$ we write $F_\beta(\ell)=\max\{0,1-\ell/\beta\}$, with $F_0\equiv0$ and $F_\infty\equiv1$, as in~\cite{DMamo,DMuniv}.

\subsection{Rational discriminants}

Let $p/q$ be a convergent of $\alpha$ and $\Delta_\varphi(E)=\tr A^{(p/q)}_q(E,\varphi)$ the discriminant of the $p/q$-periodic operator at phase $\varphi$. It is invariant under $\varphi\mapsto\varphi+1/q$, so
\[
t_q(x,E)=\Delta_{x/q}(E)
\]
is well defined, $1$-periodic and analytic in $x$ on $|\operatorname{Im}x|<q\delta_v$, and a monic real polynomial of degree $q$ in $E$~\cite[Section~5.2]{DMuniv}. For real $x$ its roots are $q$ simple real numbers, the \emph{band centres}~\cite[Lemma~5.3]{DMuniv}. For $H_\theta$ the \emph{local phase} at site $j$ is $x_j=q(\theta+j\alpha)\bmod1=x_0+j\delta$. For the almost Mathieu operator, Chambers' formula~\cite[Lemma~3.1]{DMamo} gives
\begin{equation}\label{eq:chambers}
\Delta_\varphi(E)=P_q(E)-2\lambda^q\cos(2\pi q\varphi)
\end{equation}
with a real monic polynomial $P_q=P_{p/q}$ of degree $q$ that does not depend on $\varphi$.

\subsection{Imported results}

We use the companion papers~\cite{DMamo} (cited as A) and~\cite{DMuniv} (cited as B), version~1.0.0, with their numbering. From A: Lemma~2.1 (the set $\Ga$ of~\eqref{eq:Ga} has full measure), Lemma~2.3 (uniform growth on a neighbourhood $U_h$ of $\Sigma$), Lemma~2.5 (rational spectra), Lemma~3.1 (Chambers' formula), Proposition~3.3 and equation~(9) (the conserved polynomial and its normalization $J_n$, $\sigma_n$), Lemma~3.4, Lemmas~4.1--4.3 (barriers), Proposition~4.4, Proposition~7.1, Theorem~7.3, Lemma~8.2 (elliptic blocks), Lemmas~9.2, 9.3 and~9.5 (insertion, perturbed rotations, elliptic runs), Theorems~A, B and~D, and Lemma~10.1. From B: Lemmas~2.1, 2.3, 2.7 and~2.8, Theorem~4.1 (the lower bound), Lemmas~5.2--5.5, Corollary~5.6, Lemma~5.7 (the $q$-step identity), Theorem~5.9 (the per-denominator estimate), Corollary~5.10, the proof of Theorem~5.15 (its far-energy, forbidden-decay and cutting paragraphs, which do not use its hypothesis~(W)), Lemma~5.20, Definition~5.21, Lemma~5.22, Theorem~5.26, Proposition~5.27, Lemma~5.28, Theorem~5.29, Remark~5.30, Lemmas~7.3, 7.5 and~7.6 (factorization, truncation and the mean lower bound), and Theorems~A, B and~D. Each use is stated where it occurs.

One further input is the trigonometric-polynomial product lemma of Avila and Jitomirskaya~\cite[Lemma~9.7]{AJ2009}; see also~\cite[Lemma~B.1]{JL2018}. For every irrational $\alpha$, every $x\in\R$ and every $n$, if $0\le l_0\le q_n-1$ minimizes $|\sin\pi(x+l\alpha)|$ over $0\le l\le q_n-1$, then
\begin{equation}\label{eq:AJ}
\Bigl|\sum_{0\le l\le q_n-1,\ l\ne l_0}\log|\sin\pi(x+l\alpha)|+(q_n-1)\log2\Bigr|\le C\log q_n
\end{equation}
with an absolute constant $C$. No condition on $x$ or on the arithmetic of $\alpha$ is needed. We also use the elementary three-gap facts: for $0<|d|<q_n$, $\|d\alpha\|\ge\|q_{n-1}\alpha\|\ge1/(2q_n)$, where $\|\cdot\|$ is the distance to $\Z$; so the points $\theta+s\alpha$, $s$ ranging over $q_n$ consecutive integers, are at mutual distance at least $1/(2q_n)$.

\section{Counting zeros near the real axis}\label{sec:zeros}

In this section $H_{v,\alpha}$ is supercritical. The key fact is that, for all large convergent scales and uniformly for $E$ in a fixed energy neighbourhood $J$ of the spectrum, a truncation of the rational discriminant $t_q(\cdot,E)$ has at most $2\omt(E)$ zeros in a strip whose width is proportional to $q$ (Lemma~\ref{lem:zerocount}); at each fixed spectral energy this sharpens to $2\om(E)$, with parameters and a threshold depending on the energy (Remark~\ref{rem:twice}). We state it with finite-height quantities, which is what the dynamical argument needs.

\subsection{The acceleration at a finite height}

For real $E$ and $0<y<\delta_v$ put
\[
\om_y(E)=\frac{L(E,y)-\ell(E)}{2\pi y}.
\]
For $\rho>0$ write $\Sigma_\rho=\{E\in\R:\dist(E,\Sigma)\le\rho\}$.

\begin{lemma}\label{lem:omy}
\begin{enumerate}
\item For fixed $y$, $E\mapsto\om_y(E)$ is continuous on $\R$.
\item For fixed $E$, $y\mapsto\om_y(E)$ is nondecreasing, and $\om_y(E)\to\om(E)$ as $y\downarrow0$.
\item On $\Sigma$, $\om$ is upper semicontinuous with values in $\{1,2,\dots\}$. In particular $\om^+=\max_\Sigma\om$ exists.
\end{enumerate}
\end{lemma}

\begin{proof}
Part~1 is the continuity of $L(\cdot,y)$ and of $\ell$. Part~2: the secant slope of a convex function from $0$ is nondecreasing, and its limit is the right derivative. Part~3: $\om=\inf_{y>0}\om_y$ is an infimum of continuous functions, hence upper semicontinuous; its values on $\Sigma$ are integers by (A1) and at least $1$ by Section~\ref{sec:prelim}. An upper semicontinuous function on a compact set attains its maximum.
\end{proof}

\begin{lemma}[Uniform complexified growth]\label{lem:growthy}
Let $0<y_0\le\delta_1$ and $a>0$. There are $\rho>0$ and $C$ such that
\[
\|A_r(z,x+iy)\|\le Ce^{(L(E,y)+a)r}
\]
for every real $E\in J$, every $z\in\C$ with $|z-E|\le\rho$, every real $x$, every $|y|\le y_0$ and every $r\ge0$. For all large $q$ the rational products $A^{(p/q)}_r(z,\varphi+iy)$, $0\le r\le q$, obey the same bound with $C$ replaced by $C'=C\exp(C\|v'\|_{y_0})$, where $\|v'\|_{y_0}=\max_{|\operatorname{Im}x|\le y_0}|v'(x)|$.
\end{lemma}

\begin{proof}
This is the proof of~\cite[Lemma~5.5(2)]{DMuniv}, with the type~I identity for $L(E,y)$ replaced by the joint continuity of $(E,y)\mapsto L(E,y)$. Fix $(E_0,y_0')\in J\times[-y_0,y_0]$ and pick $r_0$ with $r_0^{-1}\int_\T\log\|A_{r_0}(E_0,x+iy_0')\|\,dx<L(E_0,y_0')+a/4$. The integrand is jointly continuous in $(z,y,x)$, and Birkhoff averages over the uniquely ergodic rotation by $r_0\alpha$ converge uniformly. This gives the bound $Ce^{(L(E_0,y_0')+3a/4)r}$ on a complex neighbourhood of $(E_0,y_0')$, and joint continuity of $L$ replaces $L(E_0,y_0')$ by $L(E,y)+a/4$ nearby. Compactness of $J\times[-y_0,y_0]$ gives $\rho$ and $C$. The rational products are handled as in~\cite[Lemma~5.5(3)]{DMuniv}: the potentials of the rational and irrational products differ by at most $\|v'\|_{y_0}|\delta|$ per site, and expanding the product in these errors costs the factor $\exp(C\|v'\|_{y_0}q|\delta|)\le\exp(C\|v'\|_{y_0})$.
\end{proof}

\begin{lemma}[A limit of maxima]\label{lem:monolimit}
Let $f:[0,\infty)\times(0,\infty)\times[1/2,\infty)\to\R$ be continuous and nondecreasing in its third argument, and fix $b\ge0$. Then
\[
\lim_{y\downarrow0,\ \rho\downarrow0}\ \max_{E\in\Sigma_\rho}f\bigl(b,\ell(E),\max\{1,\om_y(E)\}\bigr)=\max_{E\in\Sigma}f\bigl(b,\ell(E),\om(E)\bigr),
\]
where $\rho$ ranges over values small enough that $\ell>0$ on $\Sigma_\rho$.
\end{lemma}

\begin{proof}
The right side is a maximum: $E\mapsto f(b,\ell(E),\om(E))$ is upper semicontinuous on $\Sigma$, because $\om$ is integer valued and upper semicontinuous and $f$ is continuous and monotone in its last argument. Call it $F$. Since $\om_y\ge\om\ge1$ on $\Sigma$, the left side is at least $F$. Suppose it is not at most $F$ in the limit. Then there are $\eps>0$, $y_k,\rho_k\to0$ and $E_k\in\Sigma_{\rho_k}$ with $f(b,\ell(E_k),\max\{1,\om_{y_k}(E_k)\})\ge F+\eps$. Pass to a subsequence with $E_k\to E_0\in\Sigma$. For fixed $y>0$ and large $k$, $\om_{y_k}\le\om_y$ by Lemma~\ref{lem:omy}(2), so by monotonicity and continuity
\[
F+\eps\le\lim_kf\bigl(b,\ell(E_k),\max\{1,\om_y(E_k)\}\bigr)=f\bigl(b,\ell(E_0),\max\{1,\om_y(E_0)\}\bigr).
\]
Letting $y\downarrow0$ gives $F+\eps\le f(b,\ell(E_0),\om(E_0))\le F$, a contradiction.
\end{proof}

\subsection{A Jensen identity and the zero count}

For a real trigonometric polynomial $\mathcal T(x)=\sum_{|k|\le K'}c_ke^{2\pi ikx}$ of exact degree $K'\ge1$, with zeros $\zeta_1,\dots,\zeta_{2K'}$ in $\C/\Z$ counted with multiplicity, and for real $y$, put $m_y(\mathcal T)=\int_0^1\log|\mathcal T(x+iy)|\,dx$.

\begin{lemma}[Jensen identity]\label{lem:jensen}
$m_y(\mathcal T)=\log|c_{K'}|+\pi\sum_k|y-\operatorname{Im}\zeta_k|$. Hence, for $Y>0$,
\[
m_Y(\mathcal T)+m_{-Y}(\mathcal T)-2m_0(\mathcal T)=2\pi\sum_{k:\,|\operatorname{Im}\zeta_k|<Y}\bigl(Y-|\operatorname{Im}\zeta_k|\bigr).
\]
\end{lemma}

\begin{proof}
By~\cite[Lemma~7.3(1)]{DMuniv}, $|\mathcal T(x)|=|c_{K'}|4^{K'}\prod_k|\sin\pi(x-\zeta_k)|$ for complex $x$, and $\int_0^1\log|\sin\pi(x+iy-\zeta)|\,dx=\pi|y-\operatorname{Im}\zeta|-\log2$~\cite[proof of Lemma~7.3(2)]{DMuniv}. Sum over $k$. For the second identity note that $|Y-b|+|Y+b|-2|b|$ equals $2(Y-|b|)$ if $|b|<Y$ and $0$ otherwise.
\end{proof}

Let $\Gamma=\max\{\log\|T(E,x)\|:E\in J,\ |\operatorname{Im}x|\le\delta_1\}$ and $K_Z=\lceil(\Gamma+2)/(\pi\delta_1)\rceil$. For $E\in J$ write $t_q(x,E)=\sum_kc_{q,k}(E)e^{2\pi ikx}$ and let
\[
T_{q,E}(x)=\sum_{|k|\le K_Z}c_{q,k}(E)e^{2\pi ikx}
\]
be its truncation. For real $E$ it is a real trigonometric polynomial.

\begin{lemma}[Zero count]\label{lem:zerocount}
Fix $y_0\in(0,\delta_1/2]$, $s\in(0,1)$, $a>0$ and $\eta>0$, and put
\[
\omt(E)=\max\Bigl\{1,\ \frac{\om_{y_0}(E)}{s}+\frac{2a+\eta}{2\pi sy_0}\Bigr\}.
\]
For all large $q$ and every $E\in J$:
\begin{enumerate}
\item $|t_q(x,E)-T_{q,E}(x)|\le e^{-q}$ for $|\operatorname{Im}x|\le qy_0$;
\item the number $N(E)$ of zeros $\zeta$ of $T_{q,E}$ in $\C/\Z$, counted with multiplicity, with $|\operatorname{Im}\zeta|\le(1-s)qy_0$, satisfies $N(E)\le2\omt(E)$.
\end{enumerate}
\end{lemma}

\begin{proof}
1. For $|\operatorname{Im}x|\le q\delta_1$ every factor of the rational product defining $t_q(x,E)$ has a phase with imaginary part at most $\delta_1$, so $|t_q(x,E)|\le2e^{\Gamma q}$. Shifting the contour of the Fourier integral gives $|c_{q,k}(E)|\le2e^{\Gamma q-2\pi|k|q\delta_1}$~\cite[proof of Lemma~7.5]{DMuniv}. For $|\operatorname{Im}x|\le qy_0\le q\delta_1/2$,
\[
\sum_{|k|>K_Z}|c_{q,k}(E)|e^{2\pi|k|qy_0}\le\sum_{|k|>K_Z}2e^{\Gamma q-\pi|k|q\delta_1}\le8e^{(\Gamma-\pi(K_Z+1)\delta_1)q}\le8e^{-2q}\le e^{-q}
\]
for large $q$, since $\pi K_Z\delta_1\ge\Gamma+2$.

2. If $T_{q,E}$ is constant there is nothing to count. Otherwise apply Lemma~\ref{lem:jensen} with $Y=qy_0$. On the line $\operatorname{Im}x=\pm Y$ we have $t_q(x,E)=\tr A^{(p/q)}_q(E,x/q)$ with $\operatorname{Im}(x/q)=\pm y_0$. So Lemma~\ref{lem:growthy}, part~1 and the evenness of $L(E,\cdot)$ give $m_{\pm Y}(T_{q,E})\le\log\bigl(2C'e^{(L(E,y_0)+a)q}+e^{-q}\bigr)\le q\bigl(L(E,y_0)+2a\bigr)$ for large $q$. The mean lower bound~\cite[Lemma~7.6]{DMuniv} gives $m_0(T_{q,E})\ge q(\ell(E)-\eta)$; its proof applies verbatim to the truncation degree $K_Z$, since the truncation error on $\R$ is at most $e^{-q}$ by part~1 and~\cite[Lemma~7.3(5)]{DMuniv} loses only $K_Z\log4$. Each zero counted in $N(E)$ contributes at least $sY$ to the sum in Lemma~\ref{lem:jensen}. Therefore
\[
2\pi sqy_0N(E)\le m_Y+m_{-Y}-2m_0\le2q\bigl(L(E,y_0)-\ell(E)+2a+\eta\bigr)=2q\bigl(2\pi y_0\om_{y_0}(E)+2a+\eta\bigr),
\]
which is $N(E)\le2\omt(E)$.
\end{proof}

All thresholds depend only on $v$, $\alpha$ and the fixed parameters, not on $E\in J$.

\begin{remark}[Twice the acceleration]\label{rem:twice}
Fix $E\in\Sigma$. Choosing $y_0,a,\eta$ small and $s$ close to $1$, depending on $E$, makes $2\omt(E)<2\om(E)+1$ by Lemma~\ref{lem:omy}(2), so $N(E)\le2\om(E)$ for all large $q$, with a threshold depending on $E$. In words: at a fixed spectral energy, in a strip of width proportional to $q$, the truncated discriminant has at most twice the acceleration many zeros. We do not claim a count that is uniform over energies with different accelerations; the arguments below keep the finite-height quantity $\omt$, which is uniform. Han and Schlag connected the acceleration with zero counts of finite-volume Dirichlet determinants~\cite[Theorem~1.2]{HanSchlag2022}, under Diophantine assumptions, and Ge, Jitomirskaya, You and Zhou identified the acceleration with the number of vanishing dual Lyapunov exponents~\cite{GJYZ2023}. Lemma~\ref{lem:zerocount} is a finite-scale statement for rational discriminants at every irrational frequency; it is proved directly from the Jensen identity and the mean lower bound.
\end{remark}

\subsection{Sublevel sets}

Write $\mathsf g(u)=|\sin\pi u|e^{-\pi|\operatorname{Im}u|}$ for $u\in\C/\Z$ and $c_0=2e^{-\pi}$, as in~\cite[Section~7.2]{DMuniv}.

\begin{lemma}[Sublevel sets]\label{lem:sublevel}
In the setting of Lemma~\ref{lem:zerocount}, let $q$ be large, $E\in J$ and $\Lambda'>0$, and put
\[
\eps_X=4^{K_Z}(\Lambda'+e^{-q})\,e^{-m_0(T_{q,E})}\,0.49^{-2K_Z}.
\]
If $\eps_X<c_0^{2K_Z}$, then $X=\{x\in\T:|t_q(x,E)|<\Lambda'\}$ is empty when $N(E)=0$, and otherwise lies in the union of $N(E)$ arcs of length $2r$, centred at the real parts of the zeros counted in $N(E)$, with $r=\eps_X^{1/N(E)}/c_0$. In particular $X$ is covered by at most $2\omt(E)$ arcs of length at most $2\eps_X^{1/(2\omt(E))}/c_0$.
\end{lemma}

\begin{proof}
Write $\mathcal T=T_{q,E}$, of exact degree $K'\le K_Z$. For real $x$, \cite[Lemma~7.3(3)]{DMuniv} holds with equality: $|\mathcal T(x)|=4^{K'}e^{m_0(\mathcal T)}\prod_k\mathsf g(x-\zeta_k)$. A zero not counted in $N(E)$ has $|\operatorname{Im}\zeta_k|>(1-s)qy_0\ge1$ for large $q$, and then $\mathsf g(x-\zeta_k)\ge0.49$~\cite[proof of Lemma~7.3(2)]{DMuniv}. If $x\in X$, then $|\mathcal T(x)|<\Lambda'+e^{-q}$ by Lemma~\ref{lem:zerocount}(1), so the product of $\mathsf g(x-\zeta_k)$ over the counted zeros is less than $\eps_X$. If $N(E)=0$ the product is $1>\eps_X$, which is impossible. Otherwise some counted zero has $\mathsf g(x-\zeta_k)<\eps_X^{1/N(E)}<c_0$. Since $\mathsf g(u)\ge c_0\min\{\dist(u,\Z),1\}$~\cite[Lemma~7.3(2)]{DMuniv}, $\dist(x-\operatorname{Re}\zeta_k,\Z)\le\dist(x-\zeta_k,\Z)<\eps_X^{1/N(E)}/c_0$. The last sentence follows from $N(E)\le2\omt(E)$ and $\eps_X<1$.
\end{proof}

At an order-one threshold, Lemma~\ref{lem:sublevel} covers the sublevel set by at most $2\omt(E)$ arcs of length at most about $(e^{-m_0(T_{q,E})})^{1/(2\omt(E))}$, where $m_0(T_{q,E})\approx q\ell(E)$; it gives upper bounds on the arcs, not their actual lengths. At a fixed spectral energy the limiting worst-cluster scale is $e^{-\ell(E)q/(2\om(E))}$. This is the price of an exponentially tight cluster of zeros, which need not coincide, and it is the source of the factor $2\om$ below.

\section{Velocity and confinement at every phase}\label{sec:velocity}

In this section $H_{v,\alpha}$ is supercritical. We prove two single-scale estimates that hold at every phase: a bound $e^{-\ell(E)q/(2\om(E))}$ on the speed (Theorem~\ref{thm:V}) and a bound on the distance the packet can reach before the drift stops it (Proposition~\ref{prop:Comega}). Both are stated with finite-height slack; the limiting rates are approached as the slack parameters tend to zero.

\subsection{Interpolation at every phase}

\begin{lemma}[Full-period interpolation]\label{lem:TI}
There is an absolute constant $C$ such that the following holds for every irrational $\alpha$, every $n$ with $q=q_n\ge8$, every $\theta\in\R$ and every set $S$ of $2K+1$ consecutive integers with $2K+1\in\{q-1,q\}$. Let $\varphi_s=\theta+s\alpha$ and
\[
l_s(\varphi)=\prod_{s'\in S\setminus\{s\}}\frac{\sin\pi(\varphi-\varphi_{s'})}{\sin\pi(\varphi_s-\varphi_{s'})}\qquad(s\in S).
\]
Then $\sup_{\varphi\in\R}|l_s(\varphi)|\le q^C$, and every trigonometric polynomial $f$ of degree at most $K$ satisfies $f=\sum_{s\in S}f(\varphi_s)\,l_s$.
\end{lemma}

\begin{proof}
Let $M\supset S$ be a set of $q$ consecutive integers, so $|M\setminus S|\le1$. By~\eqref{eq:AJ}, applied with $x=\theta+m_0\alpha-\varphi$ where $m_0=\min M$, the product $\prod_{m\in M,\,m\ne m_*}|\sin\pi(\varphi-\varphi_m)|$ is at most $2^{-(q-1)}q^{C}$, where $m_*$ minimizes $|\sin\pi(\varphi-\varphi_m)|$ over $M$. The numerator of $|l_s(\varphi)|$ is the full product over $M$ divided by the factors with $m$ in $E_s=(M\setminus S)\cup\{s\}$, a set of at most two indices. If $m_*\in E_s$, its factor cancels. The other factors of $E_s$ are at least $\sin(\pi/(4q))\ge1/(2q)$, since distinct nodes are at distance at least $1/(2q)$ and so at most one node is within $1/(4q)$ of $\varphi$. Hence the numerator is at most $2^{-(q-1)}q^C(2q)^2$. The denominator is at least $\prod_{m\in M\setminus\{s\}}|\sin\pi(\varphi_s-\varphi_m)|$, since the possible extra factor is at most one, and this is at least $2^{-(q-1)}q^{-C}$ by~\eqref{eq:AJ}, whose minimizing index is now $m=s$. Dividing gives the bound.

A product of $2K$ factors $\sin\pi(\varphi-a_i)$ is a combination of $e^{2\pi ij\varphi}$ with $|j|\le K$, so each $l_s$ is a trigonometric polynomial of degree $K$ with $l_s(\varphi_{s'})=\delta_{ss'}$. The nodes are distinct modulo one, and a nonzero trigonometric polynomial of degree at most $K$ has at most $2K$ zeros on $\T$. This gives the interpolation formula.
\end{proof}

The lemma uses no condition on the phase, and this is the point: at every phase, $2K+1\approx q$ nodes of one orbit segment control trigonometric polynomials of degree $K\approx q/2$, with only a polynomial loss.

\subsection{Regular intervals}

For an interval of length $j$ whose first site has phase $\varphi$, write $F_j(z;\varphi)=\det(z-H_{[0,j-1]}(\varphi))$, the $(1,1)$ entry of $A_j(z,\varphi)$. Fix $y_0\in(0,\delta_1]$ and $a\in(0,\min\{\ell_*,1\}/100)$. For a convergent denominator $q$ put $K=\lfloor(q-1)/2\rfloor$ and, for real $E$,
\[
\omh(E)=\max\Bigl\{1,\ \om_{y_0}(E)+\frac{3a}{2\pi y_0}\Bigr\},\qquad k(E)=\Bigl\lfloor\frac K{\omh(E)}\Bigr\rfloor .
\]

\begin{lemma}[Regular intervals]\label{lem:RN}
There is $\rho>0$ such that for all large $q$, every phase $\theta$, every real $E$ with $\dist(E,\Sigma)\le\rho$, every $z=E+i\eta$ with $0\le\eta\le\rho$, and every set $W$ of $2q$ consecutive sites, there are $k'\in\{k-2,k-1,k\}$, $k=k(E)$, and an interval $I\subset W$ of length $k'$ with
\[
|\det(z-H_I)|\ge e^{(\ell(E)-4a)k'}.
\]
\end{lemma}

\begin{proof}
Take $\rho$ no larger than the $\rho$ of Lemma~\ref{lem:growthy} with these $y_0,a$, and so small that $\ell_*/2\le\ell\le\ell^++1$ on $\Sigma_\rho$. By~\eqref{eq:thouless}, $\ell(z)\ge\ell(E)$. Since $\ell(z)=\inf_rr^{-1}\int_\T\log\|A_r(z,\varphi)\|\,d\varphi$, some real $\varphi_*$ has $\|A_k(z,\varphi_*)\|\ge e^{k\ell(E)}$. The entries of $A_k(z,\varphi_*)$ are, up to sign, $F_k(z;\varphi_*)$, $F_{k-1}(z;\varphi_*)$, $F_{k-1}(z;\varphi_*+\alpha)$ and $F_{k-2}(z;\varphi_*+\alpha)$. So for some $k'\in\{k,k-1,k-2\}$ and some real $\varphi'$, $|F_{k'}(z;\varphi')|\ge e^{k'\ell(E)}/2$.

Fix this $k'$ and let $F=F_{k'}(z;\cdot)$, a $1$-periodic function holomorphic on $|\operatorname{Im}\varphi|<\delta_v$. By Lemma~\ref{lem:growthy}, $|F(\varphi+iy)|\le\|A_{k'}(z,\varphi+iy)\|\le Ce^{(L(E,y)+a)k'}$ for $|y|\le y_0$. Shifting the Fourier integral to $\operatorname{Im}\varphi=\mp y_0$ gives $|\hat F(m)|\le Ce^{(L(E,y_0)+a)k'-2\pi|m|y_0}$. So the truncation $F_K$ of $F$ to degree $K$ satisfies on $\R$
\[
|F-F_K|\le C_{y_0}e^{(L(E,y_0)+a)k'-2\pi Ky_0}=C_{y_0}e^{(\ell(E)+2\pi y_0\om_{y_0}(E)+a)k'-2\pi Ky_0}\le C_{y_0}e^{(\ell(E)-2a)k'},
\]
because $2\pi Ky_0\ge2\pi y_0\omh(E)k'\ge(2\pi y_0\om_{y_0}(E)+3a)k'$.

Let $W=[w,w+2q-1]$ and $S=[w,w+2K]$. The intervals $[s,s+k'-1]$, $s\in S$, lie in $W$, since $2K+k'\le3K<2q$, and their first sites have the phases $\varphi_s=\theta+s\alpha$. Suppose $|F(\varphi_s)|\le e^{(\ell(E)-4a)k'}$ for every $s\in S$. Then $|F_K(\varphi_s)|\le2C_{y_0}e^{(\ell(E)-2a)k'}$, and Lemma~\ref{lem:TI} gives $\sup_\R|F_K|\le(2K+1)q^C2C_{y_0}e^{(\ell(E)-2a)k'}$. Hence $\sup_\R|F|\le q^{C+2}e^{(\ell(E)-2a)k'}<e^{k'\ell(E)}/2$ for large $q$, because $k'\ge k-2\ge q/(3\omh_{\max})-3$ with $\omh_{\max}=\max_{\Sigma_\rho}\omh<\infty$. This contradicts the choice of $\varphi'$. So some $s\in S$ has $|F(\varphi_s)|>e^{(\ell(E)-4a)k'}$, and $I=[s,s+k'-1]$ works.
\end{proof}

\begin{lemma}[Node inequality]\label{lem:node}
Let $I=[x_1,x_2]$ be as in Lemma~\ref{lem:RN} and $x=x_1+\lfloor(k'-1)/2\rfloor$. If $(H-z)g=0$ on $I$, then for $y\in\{x,x+1\}$,
\[
|g(y)|^2\le2\gamma_0^2\bigl(|g(x_1-1)|^2+|g(x_2+1)|^2\bigr),\qquad\gamma_0=Ce^{2\ell^++2}e^{-(\ell(E)/2-5a)(k(E)-2)}.
\]
\end{lemma}

\begin{proof}
The Poisson formula gives $g(y)=-G_I(x_1,y)g(x_1-1)-G_I(y,x_2)g(x_2+1)$ with $G_I=(H_I-z)^{-1}$, and Cramer's rule gives $|G_I(y,x_2)|=|\det(z-H_{[x_1,y-1]})|/|\det(z-H_I)|$ and the mirror formula for $G_I(x_1,y)$. The numerators are determinants of length at most $k'/2+1$, bounded by $Ce^{(\ell(E)+a)(k'/2+1)}$ by Lemma~\ref{lem:growthy} at height zero. With the lower bound of Lemma~\ref{lem:RN}, both Green's function entries are at most $Ce^{-(\ell(E)/2-5a)k'+\ell(E)+a}\le\gamma_0$, using $k'\ge k-2$, $\ell(E)/2-5a>0$ and $\ell(E)\le\ell^++1$. Finally $(u+v)^2\le2u^2+2v^2$.
\end{proof}

The bound uses the common length $k(E)-2$ rather than the length $k'$ selected in each window, so that $\gamma_0$ does not depend on the window.

\subsection{Velocity}

\begin{theorem}[Velocity at every phase]\label{thm:V}
Let $y_0,a$ be as above and $\rho$ as in Lemma~\ref{lem:RN}. There are $n_0$, $T_0$, $C$ and $\mu>0$ such that for all $n\ge n_0$, every phase $\theta$, every $T\ge T_0$, every real $E$ with $\dist(E,\Sigma)\le\rho$ and every $N\ge1$, the resolvent column $\psi_E(j)=\langle\delta_j,(H_\theta-E-i/T)^{-1}\delta_0\rangle$ satisfies
\[
W^+_E(N):=\sum_{j>N}|\psi_E(j)|^2\le\exp\Bigl(-\frac{N/(4q_n)-2}{1+\bar\kappa_n(E)T}\Bigr)W^+_E(0),
\]
and the mirror statement for $j<-N$, where
\[
\bar\kappa_n(E)=4\gamma_0^2\le C\exp\Bigl(-\frac{(\ell(E)-10a)\,q_n}{2\omh(E)}\Bigr).
\]
Consequently
\[
P_\theta(N,T)\le2\sup_{\dist(E,\Sigma)\le\rho}\exp\Bigl(-\frac{N/(4q_n)-2}{1+\bar\kappa_n(E)T}\Bigr)+Ce^{-\mu N}.
\]
\end{theorem}

\begin{proof}
Choose $T_0\ge1/\rho$, so that $z=E+i/T$ satisfies the hypothesis of Lemma~\ref{lem:RN}. Write $q=q_n$.

\emph{Right half-line.} The vector $g(k)=\psi_E(k)$, $k\ge0$, is square summable and solves $(H-z)g=0$ at every $k\ge1$. Let $W_i=[1+2(i-1)q,\,2iq]$, $i\ge1$. Lemma~\ref{lem:RN} gives a regular interval $I_i\subset W_i$; let $x_i$ be its node, as in Lemma~\ref{lem:node}. Put $s_i=\sum_{k>x_i}|g(k)|^2$ and $M_i=s_i-s_{i+1}$. The flux identity
\[
\frac1T\sum_{k>m}|g(k)|^2=\operatorname{Im}\bigl(g(m)\overline{g(m+1)}\bigr)\qquad(m\ge0),
\]
which follows by summing the imaginary part of the recursion multiplied by $\overline{g(k)}$ over $k>m$, holds for every real potential~\cite[proof of Lemma~9.5]{DMamo}. For $i\ge2$ the two sites adjacent to $I_i$ lie in $(x_{i-1},x_i)$ and $(x_i,x_{i+1})$, so Lemma~\ref{lem:node} gives $|g(x_i)||g(x_i+1)|\le2\gamma_0^2(M_{i-1}+M_i)$, and therefore
\[
s_i\le T|g(x_i)||g(x_i+1)|\le\frac{\bar\kappa T}2\,(s_{i-1}-s_{i+1}),\qquad\bar\kappa=4\gamma_0^2.
\]
Since $s_{i+1}\le s_i$, this gives $s_{i+1}\le\frac{\bar\kappa T}{\bar\kappa T+2}s_{i-1}\le\exp\bigl(-\frac1{1+\bar\kappa T}\bigr)s_{i-1}$. Iterating from $i=2$, $s_{2l+1}\le\exp(-l/(1+\bar\kappa T))s_1\le\exp(-l/(1+\bar\kappa T))W^+_E(0)$. Since $x_{2l+1}\le2(2l+1)q$, the choice $l=\lfloor N/(4q)\rfloor-1$ gives $x_{2l+1}\le N$ and $l\ge N/(4q)-2$, which is the claim when $l\ge0$; when $l<0$ the claim is trivial.

\emph{Left half-line.} The reflection $(Ru)(n)=u(-n)$ conjugates $H_{v,\alpha,\theta}$ to $H_{\tilde v,\alpha,-\theta}$ with $\tilde v(x)=v(-x)$. Reversing transfer products and using the evenness of $L(E,\cdot)$, $\tilde v$ has the same complexified Lyapunov exponent $L(E,y)$ for every $y$, hence the same $\ell$, $\om_{y_0}$ and $\omh$. Apply the right half-line argument to it.

\emph{The rate.} $\bar\kappa=4\gamma_0^2\le Ce^{-(\ell(E)-10a)(k(E)-2)}$, and $k(E)-2\ge K/\omh(E)-3\ge(q-2)/(2\omh(E))-3$.

\emph{The outside probability.} By~\eqref{eq:parseval},
\[
P_\theta(N,T)=\frac1{\pi T}\int_\R\|\un_{\{|j|>N\}}\psi_E\|^2\,dE,\qquad\frac1{\pi T}\int_\R\|\psi_E\|^2\,dE=1.
\] For $\dist(E,\Sigma)\le\rho$ use the two half-line bounds and $W^+_E(0)+W^-_E(0)\le\|\psi_E\|^2$. For $\dist(E,\Sigma)>\rho$ the Combes--Thomas bound~\cite[Theorem~5.15(3)]{DMuniv} gives $\|\un_{\{|j|>N\}}\psi_E\|^2\le64g^{-2}g'^{-1}e^{-g'N/4}$ with $g=\dist(E,\Sigma)$ and $g'=\min\{g,1\}\ge\min\{\rho,1\}$, which integrates to at most $Ce^{-\mu N}$. All constants are independent of $\theta$.
\end{proof}

The theorem bounds Abel averages at the time scale $T$. A bound at each individual time needs a logarithmic allowance in the radius (Lemma~\ref{lem:abel} below); the exponents are not affected.

\subsection{Confinement}

\begin{proposition}[Confinement at every phase]\label{prop:Comega}
Let $y_0\in(0,\delta_1/2]$, $s,a,\eta$ be as in Lemma~\ref{lem:zerocount}, and let $\omt_{\max}=\max_J\omt$. Fix $c>0$ and $0<d<B$, and let $\zeta>0$ and $\xi>0$ satisfy
\[
2a+\eta+\zeta\le\xi<\frac{\min\{c,\ell_*,1\}}{1000\,\omt_{\max}} .
\]
For all large $n$ with $b_n\ge c$, every phase $\theta$, every $T\in[e^{dq_n},e^{Bq_n}]$ and every real $E$ with $\dist(E,\Sigma)\le e^{-\xi q_n}$,
\[
\bigl\|\un_{\{|j|>R(E)\}}\psi_E\bigr\|\le2T^2\exp\bigl(-3\xi q_ne^{\xi q_n}\bigr),\qquad R(E)=\exp\Bigl[\Bigl(b_n-\frac{\min\{\ell(E),b_n\}}{2\omt(E)}+7\xi\Bigr)q_n\Bigr],
\]
where $\psi_E=(H_\theta-E-i/T)^{-1}\delta_0$.
\end{proposition}

\begin{proof}
We follow the confinement steps of~\cite[proof of Theorem~5.15]{DMuniv}, which do not use its hypothesis~(W), and replace its sublevel geometry by Lemma~\ref{lem:sublevel}. Write $q=q_n$, $b=b_n$, $\delta=\delta_n$, $D=e^{bq}$, $x_j=q\theta+j\delta$, $z=E+i/T$ and $d_j=t_q(x_j,z)$. Let $N_\Omega=\lfloor D/16\rfloor$ and $\Omega=[-N_\Omega,N_\Omega]\cap\Z$; the map $j\mapsto x_j\bmod1$ is injective on $\Omega$, and its image lies in an arc of length at most $1/8$. Apply the $q$-step identity~\cite[Lemma~5.7]{DMuniv} with its parameter $a$ equal to $\zeta/2$; it applies because $1/T\le e^{-dq}\le1/q$. Its error coefficient satisfies $\eps_{\mathrm{res}}\le e^{(\ell(E)+\zeta-b)q}$ for large $q$~\cite[(E2) in the proof of Theorem~5.15]{DMuniv}. Put $\Lambda=e^{3\xi q}(2+2\eps_{\mathrm{res}})$, so that $\vartheta=(2+\eps_{\mathrm{res}})/(\Lambda-\eps_{\mathrm{res}})\le e^{-3\xi q}$.

\emph{Comparison.} For real $x$, $t_q(x,\cdot)$ is monic with real roots $E_k$~\cite[Lemma~5.3]{DMuniv}, so $|t_q(x,E+i/T)|=\prod_k|E+i/T-E_k|\ge|t_q(x,E)|$. Hence $\mathcal A=\{j\in\Omega:|d_j|<\Lambda\}$ is contained in $\mathcal I=\{j\in\Omega:|t_q(x_j,E)|<\Lambda\}$.

\emph{Sublevel geometry.} We have $\Lambda\le4e^{3\xi q}\max\{1,e^{(\ell(E)+\zeta-b)q}\}$ and $m_0(T_{q,E})\ge q(\ell(E)-\eta)$ by~\cite[Lemma~7.6]{DMuniv}, as in the proof of Lemma~\ref{lem:zerocount}. So the number $\eps_X$ of Lemma~\ref{lem:sublevel}, with $\Lambda'=\Lambda$, satisfies
\[
\eps_X\le C^{K_Z}e^{(3\xi+\zeta+\eta)q}e^{-\min\{\ell(E),b\}q},
\]
which is less than $c_0^{2K_Z}$ for large $q$, because $\ell(E)\ge2\ell_*/3$ for large $q$ and so $\min\{\ell(E),b\}\ge\min\{2\ell_*/3,c\}>4\xi\ge3\xi+\zeta+\eta$. By Lemma~\ref{lem:sublevel}, $\{x:|t_q(x,E)|<\Lambda\}$ lies in at most $2\omt(E)$ arcs of length $2r$ with
\[
r\le c_0^{-1}C^{K_Z}e^{(3\xi+\zeta+\eta)q}e^{-\min\{\ell(E),b\}q/(2\omt(E))};
\]
here we used $N(E)\le2\omt(E)$ on the factor $e^{-\min\{\ell(E),b\}q}<1$ and $N(E)\ge1$ on the remaining factor. Each such arc meets the arc $\{x_j:j\in\Omega\}$ in one arc, which pulls back to an integer interval of at most $2rD+1$ sites. So $\mathcal I\subset\mathcal I'$, a union of at most $2\omt(E)$ integer intervals with
\[
\#\mathcal I'\le2\omt(E)(2rD+1)\le e^{(b-\min\{\ell(E),b\}/(2\omt(E))+4\xi)q}
\]
for large $q$, since $3\xi+\zeta+\eta<4\xi$ and $b-\min\{\ell(E),b\}/(2\omt(E))\ge0$.

\emph{Forbidden decay.} Let $q'=q+1$ and $\mathcal B=\mathcal A\cup[-q',q']\cup(\Z\setminus[-N_\Omega+q,N_\Omega-q])$. By the forbidden-decay paragraph of~\cite[proof of Theorem~5.15]{DMuniv}, which uses only the $q$-step identity and $|d_j|\ge\Lambda$ off $\mathcal A$, $|\psi_E(j)|\le\vartheta^kT$ whenever $\dist(j,\mathcal B)\ge kq$.

\emph{Cutting.} The set $\mathcal J=\mathcal I'\cup[-q',q']$ has at most $2\omt(E)+1$ components and $\#\mathcal J\le\#\mathcal I'+2q'+1$. Let $M_0=q\lceil e^{\xi q}\rceil$. The $M_0$-neighbourhood of $\mathcal J$ has at most $\#\mathcal J+(2\omt(E)+1)(2M_0+1)$ sites. So the first site $c_+>q'$ at distance at least $M_0$ from $\mathcal J$ satisfies
\[
c_+\le q'+1+\#\mathcal J+(2\omt(E)+1)(2M_0+1)\le e^{(b-\min\{\ell(E),b\}/(2\omt(E))+6\xi)q}.
\]
This is far below $N_\Omega-q-M_0$, because $\min\{\ell(E),b\}/(2\omt(E))\ge\min\{\ell_*,c\}/(3\omt_{\max})>100\xi$ and $N_\Omega\ge e^{bq}/16-1$. Hence $\dist(c_+,\mathcal B)\ge M_0$ and $|\psi_E(c_+)|\le e^{-3\xi qe^{\xi q}}T$. The cutting estimate of~\cite[proof of Theorem~5.15]{DMuniv} gives $\|\un_{(c_+,\infty)}\psi_E\|\le T|\psi_E(c_+)|$. The left side is the same. Since $c_\pm\le R(E)$, the claim follows.
\end{proof}

Only the total size of the allowed set enters the position of the first cut, not the location of its components.

\section{The all-phase upper law}\label{sec:upperlaw}

\subsection{The rate function}

For $\ell>0$, $w\ge1/2$ and $b\ge0$ put
\[
\varrho(b;\ell,w)=b-\frac{\min\{\ell,b\}}{2w},\qquad\nu(\ell,w)=\frac\ell{2w},\qquad f(b;\ell,w)=\frac{\varrho(b;\ell,w)}{\varrho(b;\ell,w)+\nu(\ell,w)} .
\]
Here $\varrho$ is the exponent of the confinement radius of Proposition~\ref{prop:Comega} and $\nu$ the exponent of the velocity in Theorem~\ref{thm:V}, both per unit of $q$. Explicitly,
\begin{equation}\label{eq:f}
f(b;\ell,w)=\begin{cases}1-\dfrac{\ell}{2wb},&b\ge\ell,\\[2mm]\dfrac{(2w-1)b}{(2w-1)b+\ell},&0\le b<\ell.\end{cases}
\end{equation}

\begin{lemma}\label{lem:f}
The function $f$ is continuous, nondecreasing in $b$ and in $w$, nonincreasing in $\ell$, and $0\le f<1$, with $f(0;\ell,w)=0$. The function $\varrho$ is nondecreasing in $w$ and nonincreasing in $\ell$. For $\varrho_0\ge0$ and $\nu_0>0$,
\begin{equation}\label{eq:tauopt}
\sup_{\tau>0}\frac{\min\{\varrho_0,(\tau-\nu_0)_+\}}{\tau}=\frac{\varrho_0}{\varrho_0+\nu_0},
\end{equation}
attained at $\tau=\varrho_0+\nu_0$.
\end{lemma}

\begin{proof}
Since $w\ge1/2$, $\varrho\ge0$, and $\varrho+\nu=b$ when $b\ge\ell$. The formulas~\eqref{eq:f} agree at $b=\ell$, and the monotonicity is read off from them. For~\eqref{eq:tauopt}: the function vanishes for $\tau\le\nu_0$, equals $(\tau-\nu_0)/\tau$, which increases, for $\nu_0\le\tau\le\varrho_0+\nu_0$, and equals $\varrho_0/\tau$, which decreases, afterwards.
\end{proof}

\subsection{The theorem}

\begin{theorem}[All-phase upper law]\label{thm:U}
Let $v$ be real analytic, $\alpha$ irrational with $\beta=\beta(\alpha)<\infty$, and $H_{v,\alpha}$ supercritical. Put
\[
\FU(\beta)=\max_{E\in\Sigma}f\bigl(\beta;\ell(E),\om(E)\bigr).
\]
For every $s'>\FU(\beta)$ and $m>0$,
\[
\sup_{\theta\in\T}P_\theta(T^{s'},T)=O(T^{-m})\qquad(T\to\infty).
\]
Hence $\sup_\theta\au(\theta)\le\aunif\le\FU(\beta)<1$. If $\beta\ge\ell^+$, then
\[
\FU(\beta)=1-\frac1{2\beta}\min_{E\in\Sigma}\frac{\ell(E)}{\om(E)}\le1-\frac{\ell_*}{2\om^+\beta}.
\]
\end{theorem}

The maximum defining $\FU$ exists because $E\mapsto f(\beta;\ell(E),\om(E))$ is upper semicontinuous on $\Sigma$ (proof of Lemma~\ref{lem:monolimit}). The last formula holds because $f(\beta;\ell(E),\om(E))=1-\ell(E)/(2\om(E)\beta)$ when $\beta\ge\ell(E)$.

\begin{proof}
It suffices to treat $s'<1$, since $P_\theta(N,T)$ is nonincreasing in $N$. Fix $s'\in(\FU(\beta),1)$, $m>0$ and $\eps'>0$ with $\FU(\beta)<s'-5\eps'$. The parameters are chosen in the following order; none depends on $\theta$.

\emph{Height and energy neighbourhood.} By Lemma~\ref{lem:monolimit}, applied with $b=\beta$, there are $y_0\in(0,\delta_1/2]$ and $\rho_0>0$ with $2\ell_*/3\le\ell\le\ell^++1$ on $\Sigma_{\rho_0}$ and
\[
\max_{E\in\Sigma_{\rho_0}}f\bigl(\beta;\ell(E),\max\{1,\om_{y_0}(E)\}\bigr)<s'-4\eps'.
\]
Let $W=2\max_J\max\{1,\om_{y_0}\}+1$. We shall take $s\in[1/2,1)$ and $a,\eta\le\pi y_0/3$, so that $\omt\le W$ and $\omh\le W$ on $J$.

\emph{The integer $K$ and the slack $\xi$.} Fix an integer $K\ge1$ with $K\ell_*/(4W)\ge4$, and then $\xi>0$ with
\[
8K\xi<\eps',\qquad\xi<\frac{s'}{10K},\qquad\xi<\frac{\min\{s'/(2K),\ell_*,1\}}{1000W}.
\]

\emph{Slack in the zero count.} Choose $s\in[1/2,1)$, $a\in(0,\min\{\ell_*,1,\pi y_0\}/100)$, $\eta\in(0,\pi y_0/3]$, $\zeta>0$ and $B_0>\beta$ with $2a+\eta+\zeta\le\xi$ and
\[
\max_{E\in\Sigma_{\rho_0}}f\bigl(B_0;\ell(E)-10a,\Omega(E)\bigr)<s'-3\eps',\qquad\Omega(E)=\max\{\omt(E),\omh(E)\}.
\]
This is possible: as $s\to1$ and $a,\eta\to0$, $\Omega$ converges to $\max\{1,\om_{y_0}\}$ uniformly on $J$, and $f$ is uniformly continuous on $[\beta,\beta+1]\times[\ell_*/2,\ell^++2]\times[1,W]$. With these choices $\nu(\ell(E)-10a,\Omega(E))\ge\ell_*/(4W)$ on $\Sigma_{\rho_0}$. Let $\rho\le\rho_0$ be as in Lemma~\ref{lem:RN} for $y_0$ and $a$, and fix the thresholds of Theorem~\ref{thm:V} and of Proposition~\ref{prop:Comega}, the latter with $c=s'/(2K)$, $d=1/K$ and $B=(B_0+1)/s'$.

\emph{Covering.} For large $T$ let $j$ be the least index with $q_j\ge K\log T$. As $T\to\infty$, $j\to\infty$.

\emph{Case A: $q_j\le T^{s'}/(\log T)^2$.} By Theorem~\ref{thm:V} at scale $j$, for $\dist(E,\Sigma)\le\rho$,
\[
\bar\kappa_j(E)T\le CT\exp\Bigl(-\frac{\ell_*q_j}{4W}\Bigr)\le CT^{1-K\ell_*/(4W)}\le1
\]
for large $T$. So, with $N=T^{s'}\ge8q_j(m\log T+2)$, which holds for large $T$,
\[
P_\theta(T^{s'},T)\le2\exp\Bigl(-\frac{N/(4q_j)-2}{2}\Bigr)+Ce^{-\mu N}\le2e^{-1}T^{-m}+Ce^{-\mu T^{s'}}=O(T^{-m}).
\]

\emph{Case B: $q_j>T^{s'}/(\log T)^2$.} Put $n=j-1$, $q=q_n<K\log T$ and $b=b_n$. Since $q_{n+1}<D_n=e^{bq}$, we have $bq>s'\log T-2\log\log T$. Since $b\le B_0$ for large $n$, $q\ge(s'\log T-2\log\log T)/B_0$, so $\log\log T=o(q)$, and $b>s'/(2K)=c$ for large $T$. With $\tau=\log T/q$ we get $1/K<\tau\le(B_0+o(1))/s'$, so $T\in[e^{dq},e^{Bq}]$. Split the energy axis.

\emph{Far energies}, $\dist(E,\Sigma)\ge e^{-\xi q}$. The Combes--Thomas bound~\cite[Theorem~5.15(3)]{DMuniv} with $N=T^{s'}\ge e^{s'q/K}\ge e^{2\xi q}$ shows that these energies contribute at most $Ce^{-e^{\xi q}/8}$ to~\eqref{eq:parseval}, exactly as in~\cite[proof of Theorem~5.15]{DMuniv}.

\emph{Near energies}, $\dist(E,\Sigma)\le e^{-\xi q}$. For such $E$ consider two radii:
\[
N_1(E)=8q(m\log T+2)\bigl(1+\bar\kappa_n(E)T\bigr),\qquad R(E)\text{ of Proposition~\ref{prop:Comega}}.
\]
Theorem~\ref{thm:V} gives the relative bound $\|\un_{\{|j|>N_1(E)\}}\psi_E\|^2\le T^{-2m}\|\psi_E\|^2$, and Proposition~\ref{prop:Comega} the absolute bound $\|\un_{\{|j|>R(E)\}}\psi_E\|^2\le4T^4e^{-6\xi qe^{\xi q}}$. Let $\mathcal E_1$ be the set of near energies with $N_1(E)\le R(E)$ and $\mathcal E_2$ the rest; both are Borel sets. We first show that $\min\{N_1(E),R(E)\}\le N$ for every near energy, where $N=T^{s'}$ is the radius used for the far energies, and then bound the whole tail at this one radius.

\emph{The radius.} Fix a near energy $E$ and write $\nu=\nu(\ell(E)-10a,\Omega(E))$ and $\varrho'=\varrho(b;\ell(E),\Omega(E))$. Since $\nu\le(\ell(E)-10a)/(2\omh(E))$, Theorem~\ref{thm:V} gives $\bar\kappa_n(E)T\le Ce^{(\tau-\nu)q}$, so $\log N_1(E)\le(\tau-\nu)_+q+O(\log\log T)$. Since $\varrho$ is nondecreasing in $w$ and $\Omega\ge\omt$, $\log R(E)\le(\varrho'+7\xi)q$. Hence, by~\eqref{eq:tauopt} and $\tau>1/K$,
\[
\frac{\log\min\{N_1(E),R(E)\}}{\log T}\le\frac{\min\{\varrho',(\tau-\nu)_+\}+7\xi}{\tau}+o(1)\le\frac{\varrho'}{\varrho'+\nu}+7K\xi+o(1).
\]
Now $\varrho(b;\ell,w)\le\varrho(b;\ell-10a,w)$, so $\varrho'/(\varrho'+\nu)\le f(b;\ell(E)-10a,\Omega(E))\le f(B_0;\ell(E)-10a,\Omega(E))<s'-3\eps'$. The $o(1)$ is uniform in $E$. Therefore $\min\{N_1(E),R(E)\}\le T^{s'-3\eps'+7K\xi+o(1)}\le T^{s'}=N$ for every near energy $E$, once $T$ is large.

\emph{The tail at radius $N$.} On $\mathcal E_1$ this gives $N_1(E)\le N$, and on $\mathcal E_2$ it gives $R(E)\le N$. Since the tails are nonincreasing in the radius, the relative bound holds at radius $N$ on $\mathcal E_1$ and the absolute bound holds at radius $N$ on $\mathcal E_2$; the far energies were estimated at radius $N$ above. So~\eqref{eq:parseval} gives
\begin{align*}
P_\theta(T^{s'},T)&\le\frac1{\pi T}\Bigl[\int_{\mathcal E_1}T^{-2m}\|\psi_E\|^2\,dE+\int_{\mathcal E_2}4T^4e^{-6\xi qe^{\xi q}}\,dE\Bigr]+Ce^{-e^{\xi q}/8}\\
&\le T^{-2m}+CT^3e^{-6\xi qe^{\xi q}}+Ce^{-e^{\xi q}/8}.
\end{align*}
Since $q\ge c'\log T$ with $c'=s'/(2B_0)$, the last two terms are $O(T^{-m})$, so $P_\theta(T^{s'},T)=O(T^{-m})$.

All constants and thresholds are independent of $\theta$. This proves the tail bound. The bound $\aunif\le\FU(\beta)$ follows from Definition~\ref{def:unif}, and $\sup_\theta\au\le\aunif$ is~\eqref{eq:unifsup}. Finally $\FU(\beta)<1$ because $f<1$ and the maximum is attained.
\end{proof}

\subsection{Consequences}

\begin{corollary}[Finite Liouville strength caps the speed]\label{cor:U1}
Let $v$ be real analytic with $H_{v,\alpha}$ supercritical. If $\beta(\alpha)<\infty$, then $\sup_\theta\au(H_{v,\alpha,\theta})\le\FU(\beta(\alpha))<1$.
\end{corollary}

\begin{remark}[The endpoint]\label{rem:endpoint}
The converse is due to Jitomirskaya and Zhang. If $\beta(\alpha)=\infty$, the Lyapunov exponent has no zeros, since it is positive on $\Sigma$ and, by (A3), off $\Sigma$, and their theorem gives quasiballistic transport, $\beta^+(p)=1$ for every $p>0$, at every phase~\cite[Theorem~3(2)]{JZ2022}; the numbering is the same in the published version and in arXiv:1510.07086v2. Together with Corollary~\ref{cor:U1}: for supercritical analytic potentials, at every phase, $\au=1$ if and only if $\beta(\alpha)=\infty$.
\end{remark}

\begin{remark}[An all-phase bracket]\label{rem:bracket}
For $0<\beta(\alpha)<\infty$ the lower bound $\au(\theta)\ge F_\beta(\ell_*)$ holds at every phase~\cite[Theorem~4.1]{DMuniv}; it needs only a Lipschitz potential whose Lyapunov exponent is continuous near a spectral minimizer, which holds for analytic $v$~\cite{BourgainJitomirskaya2002}. So at every phase
\[
F_\beta(\ell_*)\le\au(\theta)\le\FU(\beta),\qquad\text{and for }\beta\ge\ell^+:\quad1-\frac{\ell_*}\beta\le\au(\theta)\le1-\frac1{2\beta}\min_\Sigma\frac\ell\om .
\]
\end{remark}

For type~I potentials $\om=1$ on $\Sigma$, and since $f$ is nonincreasing in $\ell$, $\FU(\beta)=f(\beta;\ell_*,1)$.

\begin{corollary}[Type~I]\label{cor:typeI}
If $H_{v,\alpha}$ is supercritical and type~I and $\beta(\alpha)<\infty$, then $\aunif\le f(\beta;\ell_*,1)$. That is, $\aunif\le1-\ell_*/(2\beta)$ if $\beta\ge\ell_*$, and $\aunif\le\beta/(\beta+\ell_*)$ if $\beta<\ell_*$. In particular, for the almost Mathieu operator with $\lambda>1$ and $L=\log\lambda$,
\[
\aunif\le1-\frac{L}{2\beta}\quad(\beta\ge L),\qquad\aunif\le\frac{\beta}{\beta+L}\quad(0\le\beta<L).
\]
\end{corollary}

For comparison, at almost every phase the exponent is $F_\beta(\ell_*)$~\cite[Theorem~A]{DMuniv}. So, for type~I potentials with $\ell_*\le\beta<\infty$, the worst phase costs at most a factor $2$ in the localization strength; Section~\ref{sec:sharp} shows that this factor is attained. For $\beta<\ell_*$ the typical exponent is $0$ and the worst-phase bound is $\beta/(\beta+\ell_*)$, which is not a halving of $\ell_*$.

\subsection{An example with acceleration two}

\begin{proposition}\label{prop:twoharm}
Let $v(x)=2\lambda\cos2\pi x+2\mu\cos4\pi x$ with $\lambda\in\R$ and $|\mu|>1$, and $\alpha$ irrational. Then $\ell(E)\ge\log|\mu|$ for every real $E$, so $H_{v,\alpha}$ is supercritical. On $\Sigma$, $\om\in\{1,2\}$, and $\om(E)=2$ if and only if $\ell(E)=\log|\mu|$.
\end{proposition}

\begin{proof}
Let $w=e^{2\pi ix}$ and $w_k=we^{2\pi ik\alpha}$. Then $w_k^2T(E,x+k\alpha)$ is a matrix polynomial in $w_k$ whose value at $w_k=0$ is $\operatorname{diag}(-\mu,0)$, because the coefficient of $e^{-4\pi ix}$ in $E-v$ is $-\mu$. So $w^{2n}A_n(E,x)$ is, up to a unimodular constant, a matrix polynomial $M(w)$ with $\|M(0)\|=|\mu|^n$. For $y\ge0$, $|w|=e^{-2\pi y}\le1$ on the line $\operatorname{Im}x=y$, and subharmonicity of $\log\|M\|$ gives $\int_0^1\log\|A_n(E,x+iy)\|\,dx\ge n\log|\mu|+4\pi ny$. Hence $L(E,y)\ge\log|\mu|+4\pi|y|$ for all real $E$ and $y$, using evenness, and in particular $\ell(E)\ge\log|\mu|$. On the other hand $\|T(E,x+iy)\|\le|\mu|e^{4\pi y}+O(e^{2\pi y})$, so $L(E,y)\le\log|\mu|+4\pi y+o(1)$ as $y\to\infty$; $v$ is entire. By (A1), $L(E,\cdot)$ is convex and piecewise affine with slopes in $2\pi\Z$. The two bounds force every slope on $[0,\infty)$ to be at most $4\pi$, and $L(E,y)-4\pi y\to\log|\mu|$. At $E\in\Sigma$, $\om(E)\ge1$, so $\om(E)\in\{1,2\}$. If $\om(E)=2$, the slope is $4\pi$ on all of $(0,\infty)$, so $L(E,y)=\ell(E)+4\pi y$ and the limit gives $\ell(E)=\log|\mu|$. Conversely, if $\ell(E)=\log|\mu|$, then $L(E,y)\ge\ell(E)+4\pi y$ and slopes at most $4\pi$ force equality, so $\om(E)=2$.
\end{proof}

Suppose $\om=2$ somewhere on $\Sigma$. Then $\ell_*=\log|\mu|$, and $\ell>\log|\mu|$ where $\om=1$, so $\min_\Sigma\ell/\om=\log|\mu|/2$. Theorem~\ref{thm:U} gives, for $\beta(\alpha)\ge\ell^+$,
\[
\sup_\theta\au(\theta)\le1-\frac{\log|\mu|}{4\beta(\alpha)}.
\]
If $\lambda\ne0$ the potential is primitive. Conditionally on the weak separation hypothesis WXS$(v,\alpha)$ of~\cite[Definition~7.1]{DMuniv}, the exponent at almost every phase is $1-\log|\mu|/\beta(\alpha)$~\cite[Theorem~D]{DMuniv}. Whether the worst phase reaches $1-\log|\mu|/(4\beta)$, or only $1-\log|\mu|/(2\beta)$, is open (Section~\ref{sec:discussion}).

\section{A second proof for the almost Mathieu operator}\label{sec:amoproof}

For the almost Mathieu operator the upper bound of Corollary~\ref{cor:typeI} has a second proof, with different ingredients for both halves. The velocity comes from reflection pairs of interpolation nodes instead of full-period interpolation, and the confinement from the conserved polynomial of~\cite{DMamo} instead of the zero count. We include it because it shows where the factor $1/2$ comes from in this model: at a turning phase every node has a mirror partner with almost the same cosine, so only about $q/2$ distinct nodes are available in a period.

In this section $\lambda>1$, $L=\log\lambda$, $H=H_{\lambda,\alpha,\theta}$, and $\|x\|$ is the distance from $x$ to $\Z$.

\subsection{Reflection half-sets}

\begin{lemma}[Uniform nodes at every phase]\label{lem:L}
There is an absolute constant $C_L$ such that the following holds for every irrational $\alpha$, every $n$ with $q=q_n\ge8$, every $\psi\in\R$ and every integer $m_0$. Put $\varphi_m=\psi+m\alpha$ and $M=\{m_0,\dots,m_0+q-1\}$. There is a set $\mathcal J\subset M$ of $\lfloor(q-1)/2\rfloor$ consecutive integers such that for every $j\in\mathcal J$ and every $c\in[-1,1]$,
\[
\prod_{i\in\mathcal J\setminus\{j\}}\frac{|c-\cos2\pi\varphi_i|}{|\cos2\pi\varphi_j-\cos2\pi\varphi_i|}\le q^{C_L}.
\]
The same holds, with a larger absolute constant, for every $\mathcal J'\subset\mathcal J$ with $|\mathcal J\setminus\mathcal J'|\le3$.
\end{lemma}

\begin{proof}
Write $\mathsf s(x)=|\sin\pi x|$.

\emph{Reflection centre.} Let $c_*$ minimize $\|2\psi+c\alpha\|$ over the integers $c\in R_c=[2m_0,2m_0+2q-2]$, and let $\eps_*\in(-1/2,1/2]$ represent $2\psi+c_*\alpha$ modulo one. Since $R_c$ contains $q$ consecutive integers, the three-gap facts give $|\eps_*|\le2/q$. Put $\mu=c_*/2\in[m_0,m_0+q-1]$. If $\mu\le m_0+q/2-1$, let $\mathcal J$ be the set of integers in $(\mu,\mu+q/2)$; otherwise the set of integers in $(\mu-q/2,\mu)$. In both cases $\mathcal J\subset M$ is an interval of integers, and after removing one endpoint if necessary, $|\mathcal J|=\lfloor(q-1)/2\rfloor$.

\emph{The involution.} For $m\in M$ let $\iota(m)$ be the element of $M$ with $\iota(m)\equiv c_*-m\pmod q$. Both $m+\iota(m)$ and $c_*$ lie in $R_c$, so $m+\iota(m)=c_*+t_mq$ with $t_m\in\{-1,0,1\}$, and
\begin{equation}\label{eq:pair}
\varphi_m+\varphi_{\iota(m)}\equiv\eps_*+t_m\delta_n\pmod1,\qquad|\eps_*+t_m\delta_n|\le3/q.
\end{equation}
For distinct $j,m\in\mathcal J$ we have $0<|j+m-c_*|<q$, because both lie on the same side of $\mu$ at distance less than $q/2$. So $\iota(\mathcal J)\cap\mathcal J=\emptyset$. Let $E_j=M\setminus\bigl((\mathcal J\setminus\{j\})\cup(\iota(\mathcal J)\setminus\{\iota(j)\})\bigr)$; then $|E_j|\le6$.

\emph{Separation.} Distinct $m,m'\in M$ have $\|\varphi_m-\varphi_{m'}\|\ge1/(2q)$. For distinct $j,m\in\mathcal J$ put $d=j+m-c_*$, so $0<|d|<q$ and $\|d\alpha\|\ge1/(2q)$. Since $j+m\in R_c$, minimality gives $\|\eps_*+d\alpha\|\ge\|\eps_*\|$. Hence
\[
\|\varphi_j+\varphi_m\|=\|\eps_*+d\alpha\|\ge\max\{\|\eps_*\|,\|d\alpha\|-\|\eps_*\|\}\ge\frac{\|d\alpha\|}2\ge\frac1{4q}.
\]
This lower bound holds even when $\eps_*=0$; it is what fails if both members of a mirror pair are kept.

\emph{Splitting.} Use $|\cos2\pi a-\cos2\pi b|=2\mathsf s(a+b)\mathsf s(a-b)$. For real $a$,
\[
S(a):=\sum_{i\in\mathcal J\setminus\{j\}}\log|\cos2\pi a-\cos2\pi\varphi_i|=(|\mathcal J|-1)\log2+S_-(a)+S_+(a),
\]
with $S_-(a)=\sum_{i\in\mathcal J\setminus\{j\}}\log\mathsf s(a-\varphi_i)$ and, by~\eqref{eq:pair}, $S_+(a)=\sum_{i'\in\iota(\mathcal J)\setminus\{\iota(j)\}}\log\mathsf s(a+\eta_{i'}-\varphi_{i'})$ with $|\eta_{i'}|\le3/q$.

\emph{Upper bound at every $a$.} Let $m_*$ minimize $\mathsf s(a-\varphi_m)$ over $M$. For $m\ne m_*$, $\|a-\varphi_m\|\ge1/(4q)$, so $\mathsf s(a-\varphi_m)\ge1/(2q)$, and the ordered distances give $\sum_{m\ne m_*}1/\mathsf s(a-\varphi_m)\le Cq\log q$. Drop the $m_*$-term wherever it occurs; it is at most zero. For the shifted terms with $i'\ne m_*$, $\log\mathsf s(a+\eta-\varphi_{i'})\le\log\mathsf s(a-\varphi_{i'})+\pi|\eta|/\mathsf s(a-\varphi_{i'})$, and these corrections add up to at most $C\log q$. Adding back the at most six omitted terms of $E_j$, each at least $-\log(2q)$, and using~\eqref{eq:AJ},
\[
S_-(a)+S_+(a)\le\sum_{m\in M,\,m\ne m_*}\log\mathsf s(a-\varphi_m)+C\log q\le-(q-1)\log2+C'\log q.
\]

\emph{Lower bound at $a=\varphi_j$.} Now $m_*=j$. The unshifted part is at least $\sum_{m\in M,\,m\ne j}\log\mathsf s(\varphi_j-\varphi_m)\ge-(q-1)\log2-C\log q$ by~\eqref{eq:AJ}, since removing terms that are at most zero only increases a sum. Each shift changes a term by $\log\mathsf s(\varphi_j+\varphi_{\iota^{-1}(i')})-\log\mathsf s(\varphi_j-\varphi_{i'})$. Both arguments have sine at least $1/(2q)$ by the separation step. For the indices with $\|\varphi_j-\varphi_{i'}\|\le12/q$, of which there are at most $49$ by the separation $1/(2q)$, the change is at least $-\log(2q)$. For the others $\mathsf s(\varphi_j-\varphi_{i'})\ge24/q$, so $\pi|\eta_{i'}|/\mathsf s(\varphi_j-\varphi_{i'})\le\pi/8<1/2$, and the change is at least $-2\pi|\eta_{i'}|/\mathsf s(\varphi_j-\varphi_{i'})$, and these add up to at least $-C\log q$. So $S_-(\varphi_j)+S_+(\varphi_j)\ge-(q-1)\log2-C''\log q$.

\emph{Conclusion.} Every $c\in[-1,1]$ is $\cos2\pi a$ for a real $a$, and the logarithm of the product in the lemma is $S(a)-S(\varphi_j)\le(C'+C'')\log q$. Removing at most three further nodes changes only the number of omitted factors by a bounded amount.
\end{proof}

\subsection{Velocity \texorpdfstring{$\lambda^{-q/2}$}{lambda to the -q/2}}

Fix $0<\zeta<L/8$ and $h=L+\zeta$. By Lemma~\ref{lem:growthy} at height zero and continuity of $\ell$ there are $r_0>0$, a real neighbourhood $U'_h$ of $\Sigma$ and $C_h$ such that $\|A_r(z,x)\|\le C_he^{hr}$ for all $z=E+i\eta$ with $E\in U'_h$ and $0\le\eta\le r_0$, all real $x$ and all $r\ge0$. For $q=q_n$ put $k_1=\lfloor(\lfloor(q-1)/2\rfloor-2)/2\rfloor$ and $k=2k_1+1$, so that $k+1\le\lfloor(q-1)/2\rfloor$ and $k\ge q/2-4$.

\begin{lemma}[Regular intervals for the almost Mathieu operator]\label{lem:RNamo}
For all large $n$, every phase $\theta$, every $z=E+i\eta$ with $E\in U'_h$ and $0\le\eta\le r_0$, and every set $W=[w,w+q+2k_1-1]$ of consecutive sites, there is $x\in[w+k_1,w+k_1+q-1]$ such that $I=[x-k_1,x+k_1]\subset W$ satisfies $|\det(z-H_I)|\ge e^{(L-\zeta)k}$. For every solution of $(H-z)g=0$ on $I$ and $y\in\{x,x+1\}$,
\[
|g(y)|^2\le2\gamma_0^2\bigl(|g(x-k_1-1)|^2+|g(x+k_1+1)|^2\bigr),\qquad\gamma_0=C_he^he^{-(L/2-2\zeta)k}.
\]
\end{lemma}

\begin{proof}
The determinant of $z-H$ on an interval of odd length $k$ centred at a site with phase $\varphi$ is a polynomial $Q_k(z;\cos2\pi\varphi)$ of degree at most $k$ in $\cos2\pi\varphi$: it is a trigonometric polynomial of degree $k$ in $\varphi$, and reversing the interval about its centre changes $\varphi$ into $-\varphi$ without changing the determinant. Apply Lemma~\ref{lem:L} with $\psi=\theta$ and $M=[w+k_1,w+k_1+q-1]$, so that $\varphi_m$ is the centre phase of $[m-k_1,m+k_1]\subset W$, and take $\mathcal J'\subset\mathcal J$ with $k+1$ elements. If $|Q_k(z;\cos2\pi\varphi_m)|\le e^{(L-\zeta)k}$ for all $m\in\mathcal J'$, Lagrange interpolation through these $k+1$ nodes gives $\sup_{[-1,1]}|Q_k(z;\cdot)|\le(k+1)q^{C}e^{(L-\zeta)k}$, so $\int_\T\log|\det(z-H_{[0,k-1]}(\theta'))|\,d\theta'\le(L-\zeta)k+C\log q<kL$ for large $n$. This contradicts Herman's bound $\int_\T\log|\det(z-H_{[0,k-1]}(\theta'))|\,d\theta'\ge kL$, valid for every $z\in\C$~\cite[proof of Lemma~6.2]{DMamo}: the determinant is a trigonometric polynomial in $\theta'$ whose extreme coefficient has modulus $\lambda^k$. So some $m=x\in\mathcal J'$ gives a regular interval. The node inequality follows from the Poisson formula and Cramer's rule as in Lemma~\ref{lem:node}: the numerators have length at most $k_1+1$ and are at most $C_he^{h(k_1+1)}$, and $hk_1-(L-\zeta)(2k_1+1)\le-(L/2-2\zeta)k$.
\end{proof}

\begin{theorem}[Velocity for the almost Mathieu operator]\label{thm:Vamo}
For every $0<\zeta<L/8$ there are $n_0$, $T_0$, $C$ and $\mu>0$ such that for all $n\ge n_0$, every phase $\theta$, every $T\ge T_0$ and every $N\ge1$,
\[
P_\theta(N,T)\le2\exp\Bigl(-\frac{N/(4q_n)-2}{1+\kappa_nT}\Bigr)+Ce^{-\mu N},\qquad\kappa_n=4C_h^2e^{2h}e^{-(L-4\zeta)k}\le C'\lambda^{-q_n/2}e^{2\zeta q_n}.
\]
\end{theorem}

\begin{proof}
The proof of Theorem~\ref{thm:V}, with Lemma~\ref{lem:RNamo} in place of Lemmas~\ref{lem:RN} and~\ref{lem:node}, windows $W_i=[1+(i-1)(q+2k_1),\,i(q+2k_1)]$ of length at most $2q$, $T_0\ge1/r_0$, and the Combes--Thomas bound for energies outside $U'_h$. The left half-line is the right half-line of $H_{\lambda,\alpha,-\theta}$, since the cosine is even.
\end{proof}

\subsection{Confinement from the conserved polynomial}

We use the conserved polynomial $Q=P_q(H)=V+K$ of~\cite[Proposition~3.3]{DMamo}, with $V(j)=2\lambda^q\cos(2\pi(q\theta+j\delta))$, $K$ of range $q$, and its normalization $\|K\|\le J_n$, $\sigma_n=J_n/\lambda^{q_n}$~\cite[eq.~(9)]{DMamo}.

\begin{proposition}[Confinement for the almost Mathieu operator]\label{prop:Camo}
Fix $c>0$ and $0<\eps<\min\{L,c\}/4$, and use~\cite[eq.~(9)]{DMamo} with $h=L+\eps/8$. There are an absolute constant $C_0$ and a constant $C$ such that for all large $n$ with $b_n\ge c$, every phase $\theta$ and all $T>0$,
\[
P_\theta(R_n,T)\le C(1+T^2)\exp\Bigl(-\frac{e^{\eps q_n/16}}{q_n}\Bigr),\qquad R_n=\bigl\lceil C_0\sigma_n^{1/2}|\delta_n|^{-1}\bigr\rceil+6G_n,\quad G_n=\lceil q_ne^{\eps q_n/2}\rceil,
\]
and $R_n\le e^{(b_n-\min\{L,b_n\}/2+\eps)q_n}$.
\end{proposition}

\begin{proof}
Write $q=q_n$, $\delta=\delta_n$, $G=G_n$, $\kappa=1/\sigma_n$ and $u=q\theta$. Apply the barrier lemmas of~\cite[Section~4]{DMamo} to $Q'=Q/J_n=V'+K'$, where $\|K'\|\le1$, $K'$ has range $q$ and $V'(j)=2\kappa\cos(2\pi(u+j\delta))$, with $d_0=V'(0)$, $W=8$ and $V_0=16$. Dividing $Q$ by $J_n$ changes neither its spectral projections nor its commutation with $H$.

\emph{Sublevel geometry.} With $x=2\pi(u+j\delta)$ and $x_0=2\pi u$, the set $\{j:|V'(j)-d_0|<32\}=\{j:|\cos x-\cos x_0|<16/\kappa\}$ meets one period, $|j|\le1/(2|\delta|)$, in at most two arcs, each of $x$-length at most $C\kappa^{-1/2}$, uniformly in $x_0$. On the two complementary gaps $V'-d_0$ has a fixed sign and modulus at least $32$. Let $\mathcal J_0$ be the arc containing $0$, merged with the other arc if one of the two gaps between them has fewer than $4G$ sites. Since $\kappa\ge e^{(\min\{L,c\}-\eps/4)q}$ is huge and the period has $|\delta|^{-1}$ sites, the two gaps cannot both be short, and after the merging rule each gap adjacent to $\mathcal J_0$ has at least $4G$ sites. (One of the gaps can be short near a turning value, so both gaps need not be long; this is why the merging rule is used.) Thus $\mathcal J_0$ is an integer interval containing $0$ of length at most $C_0\kappa^{-1/2}|\delta|^{-1}+4G$, contained in $[-(R_n-G),R_n-G]$ after enlarging $C_0$, and followed on each side by at least $4G$ sites on which $V'-d_0$ has a fixed sign and modulus at least $32$.

\emph{Outer barrier.} Write $\mathcal J_0=[e_-,e_+]$, $N_+=e_++2G$ and $N_-=|e_-|+2G$, both at most $R_n$. The intervals $I_+=[e_++1,e_++3G]$ and $I_-=[e_--3G,e_--1]$ lie in the adjacent gaps, and the pairs $\{N_+,N_++1\}$ and $\{-N_--1,-N_-\}$ are at distance at least $G-1$ from their complements. \cite[Lemma~4.1]{DMamo}, with $J=1$, $V_0-W=8$ and $\ell=\lfloor(G-2)/q\rfloor-2$, gives $\|\un_B\Pi\|\le2\cdot8^{-(\ell+2)}\le\exp(-e^{\eps q/2}/q)$ for the set $B$ of these four sites and $\Pi=\un_{[d_0-W,d_0+W]}(Q')$.

\emph{Inner core.} For $|j|\le r$, $|V'(j)-d_0|\le4\pi\kappa|j\delta|+4\pi^2\kappa j^2\delta^2\le2=W/4$ provided $r\le\min\{1/(8\pi\kappa|\delta|),(8\pi^2\kappa)^{-1/2}|\delta|^{-1}\}$. From $J_n\ge C_hq^{5/2}e^{(L+\eps/8)q}|\delta|$ we get, at every scale,
\[
\kappa|\delta|\le C^{-1}q^{-5/2}e^{-\eps q/8},\qquad\kappa^{-1/2}|\delta|^{-1}\ge C^{1/2}q^{5/4}e^{(b_n+\eps/8)q/2}.
\]
Since $b_n\ge c>4\eps$, we may take $r\ge e^{\eps q/16}$ for large $q$. \cite[Lemma~4.2]{DMamo} with $\ell_0=\lfloor r/q\rfloor-2$ gives $\|(1-\Pi)\delta_0\|\le3\cdot6^{-(\ell_0+2)}$.

\emph{Cutoff.} The proof of~\cite[Lemma~4.3]{DMamo} works verbatim for the asymmetric window $\chi=\un_{[-N_-,N_+]}$, since $[H,\chi]$ still has norm one and lives on the two boundary bonds. It bounds the Abel-averaged mass outside $[-N_-,N_+]$, which dominates $P_\theta(R_n,T)$, by $8\|(1-\Pi)\delta_0\|^2+T^2\|\un_B\Pi\|^2$. This gives the bound. All constants are independent of $\theta$. Finally $\sigma_n\le e^{-(\min\{L,b_n\}-\eps/4)q}$ by~\cite[Lemma~3.4]{DMamo}, which gives the bound on $R_n$.
\end{proof}

\begin{corollary}\label{cor:amo2}
For the almost Mathieu operator with $\beta(\alpha)<\infty$, the conclusion $\sup_\theta P_\theta(T^a,T)=O(T^{-m})$ for every $a>f(\beta;L,1)$ and $m>0$ follows from Theorem~\ref{thm:Vamo} and Proposition~\ref{prop:Camo}, without Sections~\ref{sec:zeros} and~\ref{sec:velocity}.
\end{corollary}

\begin{proof}
Run the covering argument of Theorem~\ref{thm:U}. In Case~A use Theorem~\ref{thm:Vamo}. In Case~B both Theorem~\ref{thm:Vamo} and Proposition~\ref{prop:Camo} bound the same outside probability, uniformly in energy, so $P_\theta(\min\{N_1,R_n\},T)\le T^{-m}$ with $N_1=8q(m\log T+2)(1+C'Te^{-(L/2-2\zeta)q})$, and
\[
\frac{\log\min\{N_1,R_n\}}q\le\min\bigl\{\varrho(b;L,1)+\eps,\ (\tau-L/2)_++2\zeta\bigr\}+o(1).
\]
By~\eqref{eq:tauopt} the exponent is at most $f(b;L,1)+K(\eps+2\zeta)+o(1)$, and $\eps,\zeta$ can be taken small.
\end{proof}

\section{Sharpness: closed gaps and the exceptional phases}\label{sec:sharp}

The factor $2\om$ in Theorem~\ref{thm:U} comes from the worst possible cluster of zeros. This section shows that it is attained. The mechanism is a gap of a periodic approximant that closes, or nearly closes, at a particular phase: next to a closed gap the band is wide, of width about $e^{-\ell q/2}$ rather than $e^{-\ell q}$, and a packet launched where the local phase sits at that phase moves with the corresponding speed.

\subsection{Closed central gaps}

Call $v$ \emph{odd-harmonic} if $v(x+1/2)=-v(x)$. The cosine is odd-harmonic. Write $T(y)=\bigl(\begin{smallmatrix}y&-1\\1&0\end{smallmatrix}\bigr)$ and $K_0=\operatorname{diag}(1,-1)$, so that $K_0T(y)K_0=-T(-y)$.

\begin{proposition}[Closed central gaps]\label{prop:G}
Let $v$ be continuous and odd-harmonic, let $p/q$ be in lowest terms with $q$ even, and put $s=(-1)^{q/2}$ and $v_k=v(\varphi+kp/q)$. Let $M_1(\varphi)=T(-v_{q/2-1})\cdots T(-v_0)$ be the transfer matrix at $E=0$ over half a period, and $g(\varphi)=(M_1)_{11}-(M_1)_{22}$. Then:
\begin{enumerate}
\item $\Delta_\varphi(0)=s\bigl(g(\varphi)^2+2\bigr)$ for every $\varphi$;
\item $g(\varphi+p/q)=-g(\varphi)$, so $g$ has a zero $\varphi_0$, and the monodromy of the $p/q$-periodic operator at $E=0$ and phase $\varphi_0$ is $sI$: the gap at $E=0$ is closed;
\item $\Delta_\varphi(-E)=\Delta_\varphi(E)$ for all $\varphi$ and $E$.
\end{enumerate}
\end{proposition}

\begin{proof}
Since $q$ is even, $p$ is odd, so $v_{k+q/2}=v(\varphi+kp/q+1/2)=-v_k$. Hence the second half period contributes $\prod_kT(v_k)=\prod_k(-K_0T(-v_k)K_0)=(-1)^{q/2}K_0M_1K_0$, and the monodromy is $M=sK_0M_1K_0M_1$. If $M_1=\bigl(\begin{smallmatrix}a&b\\c&d\end{smallmatrix}\bigr)$, then $K_0M_1K_0M_1=\bigl(\begin{smallmatrix}a^2-bc&b(a-d)\\c(d-a)&d^2-bc\end{smallmatrix}\bigr)$, whose trace is $(a-d)^2+2$ because $ad-bc=1$. This is part~1. At a zero of $g$ the matrix $K_0M_1K_0M_1$ is the identity, which gives the closed gap in part~2.

Shifting $\varphi$ by $p/q$ moves the half period by one site: $M_1(\varphi+p/q)=T(-v_{q/2})M_1(\varphi)T(-v_0)^{-1}=T(y)M_1(\varphi)T(-y)^{-1}$ with $y=v_0$. A direct computation shows that the diagonal of $T(y)M_1T(-y)^{-1}$ is $(d-yb,\ a-yb)$, so its diagonal difference is $d-a=-g(\varphi)$. The intermediate value theorem gives a zero.

For part~3, odd-harmonicity gives
\[
A^{(p/q)}_q(-E,\varphi+1/2)=\prod_kT(-(E-v_k))=(-1)^qK_0A^{(p/q)}_q(E,\varphi)K_0,
\]
so $\Delta_{\varphi+1/2}(-E)=\Delta_\varphi(E)$. Since $1/2=(q/2)\cdot(1/q)$ is a period of $\varphi\mapsto\Delta_\varphi$, part~3 follows.
\end{proof}

\begin{corollary}\label{cor:closedgap}
For the almost Mathieu operator and even $q$, the Chambers polynomial satisfies $P_q(0)=s(2\lambda^q+2)$ and $P_q(-E)=P_q(E)$.
\end{corollary}

\begin{proof}
By~\eqref{eq:chambers}, $s\Delta_\varphi(0)=sP_q(0)-2s\lambda^q\cos(2\pi q\varphi)$, whose minimum over $\varphi$ is $sP_q(0)-2\lambda^q$. By Proposition~\ref{prop:G}(1)--(2), this minimum is $2$. Evenness follows from part~3 and~\eqref{eq:chambers}.
\end{proof}

\begin{remark}[Priority]\label{rem:LM}
The closed middle gap of even-denominator almost Mathieu approximants is classical; see Bellissard and Simon~\cite[Theorem~3]{BellissardSimon1982}. The identity of Corollary~\ref{cor:closedgap} is due to Lamoureux and Mingo~\cite[Corollary~2.6 and Theorem~2.7]{LamoureuxMingo2007}; their value $2(-1)^{q/2}(1+(\lambda'/2)^q)$ is ours with $\lambda'=2\lambda$. Proposition~\ref{prop:G} gives an elementary proof that applies to every continuous odd-harmonic potential.
\end{remark}

\subsection{A wide band at the turning phase}

For even $q=q_n$ write $s=(-1)^{q/2}$ and $\Delta_t(E)=P_q(E)-2s\lambda^q$, the discriminant at the turning phases $\varphi$ with $\cos(2\pi q\varphi)=s$.

\begin{proposition}\label{prop:wideband}
Let $\lambda>1$, $h>L$ and $0<\eta\le1/4$, and let $C_h$ be the constant of~\cite[Proposition~3.3(2)]{DMamo} for $r=2$, so that $\sup_{U_h}|P_q''|\le C_hq^2e^{hq}$. For all large $n$ with $q=q_n$ even there is $E^{(1)}\ge0$ with
\[
|\Delta_t(E^{(1)})|=2-\eta\qquad\text{and}\qquad|P_q'(E^{(1)})|\le q\sqrt{2C_h\eta}\,e^{hq/2}.
\]
\end{proposition}

\begin{proof}
At a phase with $\cos(2\pi q\varphi)=0$, \eqref{eq:chambers} gives $t_q=P_q$, so every root of $P_q$ is real and simple~\cite[Lemma~5.3]{DMuniv}. Away from its roots, $(\log|P_q|)''=-\sum_j(E-r_j)^{-2}<0$, so every critical point of $P_q$ is a strict local maximum of $|P_q|$. By Corollary~\ref{cor:closedgap}, $P_q'(0)=0$ and $sP_q$ has a strict local maximum at $0$ with value $2\lambda^q+2$. Let $f(E)=s(P_q(0)-P_q(E))$ and let $E^{(1)}$ be the first $E>0$ with $|f(E)|=\eta$. On $[0,E^{(1)}]$, $sP_q>0$, and $f\ge0$: otherwise $sP_q$ would exceed its value at $0$ somewhere in $(0,E^{(1)})$, forcing a positive local minimum of $|P_q|$ in between. So $f(E^{(1)})=\eta$ and $s\Delta_t=2-f\in[2-\eta,2]$ on $[0,E^{(1)}]$. This interval therefore lies in the spectrum of a turning approximant, hence in $U_h$ for large $n$~\cite[Lemma~2.5]{DMamo}, where $|f''|\le M:=C_hq^2e^{hq}$. Since $f(0)=f'(0)=0$ and $f'(E^{(1)})\ge0$, we have $f'(E^{(1)})\le ME^{(1)}$. With $t=f'(E^{(1)})/M\in[0,E^{(1)}]$, Taylor's formula gives $0\le f(E^{(1)}-t)\le\eta-f'(E^{(1)})^2/(2M)$. Hence $|P_q'(E^{(1)})|=f'(E^{(1)})\le\sqrt{2M\eta}$.
\end{proof}

\subsection{The lower bound at the phase \texorpdfstring{$1/4$}{1/4}}

\begin{theorem}\label{thm:quarter}
Let $\lambda>1$ and $\alpha$ irrational. For every $h>L$ and $\eta\in(0,1/4]$ there are $c_\eta>0$ and $n_0$ such that for every $n\ge n_0$ with $q_n$ even and $b_n\ge h/2+(h-L)$, and every $T$ with $e^{(h/2+(h-L))q_n}\le T\le|\delta_n|^{-1}$,
\begin{equation}\label{eq:quarter}
P_\varphi\bigl(\lfloor c_\eta Te^{-hq_n/2}\rfloor,\,T\bigr)\ge T^{-2}e^{-2hq_n}
\end{equation}
holds for $\varphi=1/4$, and, with modified constants, for every phase $\varphi$ with $\|\varphi-1/4\|\le2|\delta_n|$ and for every phase $\varphi$ with $|\cos(2\pi q_n\varphi)-(-1)^{q_n/2}|\le\eta\lambda^{-q_n}/16$. Consequently
\[
\au(H_{\lambda,\alpha,1/4})\ge\max\Bigl\{0,\ 1-\frac{L}{2\beta_e(\alpha)}\Bigr\},
\]
read as $1$ if $\beta_e(\alpha)=\infty$.
\end{theorem}

\begin{proof}
Write $q=q_n$, $p=p_n$ (odd), $\delta=\delta_n$, $s=(-1)^{q/2}$ and $\eps=1/T$, so $\eps\ge|\delta|$. We first treat $\varphi=1/4$.

\emph{Step 0: the turning block.} The first $q$-block of $H_{1/4}$ to the right, sites $1,\dots,q$, has phases $x_0+k\alpha$ with $x_0=1/4+\alpha$. Compare it with the rational block $R(E)=A^{(p/q)}_q(E,x_0')$, $x_0'=1/4+p/q$. Since $\cos(2\pi qx_0')=\cos(\pi q/2)=s$, \eqref{eq:chambers} gives $\tr R(E)=\Delta_t(E)$. The phase discrepancy at site $k$ is $(k+1)\delta/q$, so the potentials differ by $|w_k|\le4\pi\lambda|\delta|$.

\emph{Step 1: an energy window.} Let $E^{(1)}$ be as in Proposition~\ref{prop:wideband} and $I_*=[E^{(1)}-r,E^{(1)}+r]$ with $r=\eta e^{-hq/2}/(C_1q)$ and $C_1$ large. For $E\in I_*$, which lies in $U_h$, $|P_q'(E)|\le C_2qe^{hq/2}$ and $|\Delta_t(E)-\Delta_t(E^{(1)})|\le\eta/2$, so $|\tr R(E)|\le2-\eta/2$, $|\sin\varrho_R|\ge\sqrt\eta/2$ for the rotation angle $\varrho_R$, and by~\cite[Lemma~8.2(1)]{DMamo} the block sum satisfies $S_R=|P_q'(E)|/|\sin\varrho_R|\le C_3qe^{hq/2}$.

\emph{Step 2: insertion.} Let $B_R\in SL(2,\R)$ map $i$ to the fixed point of $R(E)$. The insertion lemma~\cite[Lemma~9.2]{DMamo} with the shifts $w_k$ gives $A_k(E,x_0)B_R=A^{(p/q)}_k(E,x_0')B_R\Pi_k$ with
\[
\|\Pi_k-I\|\le y_0:=\exp(4\pi\lambda|\delta|S_R)-1\le C q e^{-(h-L)q},
\]
using $|\delta|\le e^{-(h/2+(h-L))q}$. So $|\tr A_q(E,x_0)|\le2-\eta/2+2y_0\le2-\eta/4$. By~\cite[Lemma~9.3]{DMamo} the fixed point of $B_R^{-1}A_q(E,x_0)B_R$ is within hyperbolic distance $12y_0/(\sqrt\eta/2)$ of $i$; let $N\in SL(2,\R)$ map $i$ to it, so $\|N\|^2\le2$, and put $B_A=B_RN$, which maps $i$ to the fixed point of $A_q(E,x_0)$. The block sum of $A_q(E,x_0)$ in this frame is $S_0\le2(1+y_0)^2S_R\le3C_3qe^{hq/2}$.

\emph{Step 3: an elliptic run.} Apply~\cite[Lemma~9.5]{DMamo} to $g(k)=G(k;E+i\eps)=\langle\delta_k,(H_{1/4}-E-i\eps)^{-1}\delta_0\rangle$, $k\ge0$, which solves the half-line equation for $k\ge1$, with ellipticity margin $\eta'=\eta/4$, frame $B_A$, block sum $S_0$, and $\xi=(4\pi\lambda+1)\eps$. Its run length satisfies $m_*\ge c_\eta Te^{-hq/2}/q\ge8$ for large $n$, because $Te^{-hq/2}\ge e^{(h-L)q}$. With $m_0=\lceil m_*/2\rceil$ the lemma gives $\sum_{k>m_0q}|g(k)|^2\ge\eta'^{5/2}2^{-22}\|\Phi_0\|^2/\xi$ and $|g(0)|\le\|B_A\|\,\|\Phi_0\|$. By~\cite[Lemma~8.2(4)]{DMamo} and the growth bound, $\|B_A\|^2\le1+2\|A_q(E,x_0)\|/|\sin\varrho_A|\le C_\eta e^{hq}$, and $|G(0;E+i\eps)|\ge\operatorname{Im}G(0;E+i\eps)\ge\eps/D_\lambda^2$ with $D_\lambda=4(1+\lambda)+1$. Hence
\[
\sum_{k>m_0q}|G(k;E+i\eps)|^2\ge c'_\eta\,\eps\,e^{-hq}\qquad(E\in I_*).
\]

\emph{A common cutoff.} The run length depends on $E$, so we fix one radius before integrating: $N_*=\lfloor c_\eta Te^{-hq/2}\rfloor$, with $c_\eta$ reduced so that $N_*\le m_0(E)q$ for every $E\in I_*$. By monotonicity of tails, $\sum_{k>N_*}|G(k;E+i\eps)|^2\ge c'_\eta\eps e^{-hq}$ on $I_*$.

\emph{Step 4: Parseval.} By~\eqref{eq:parseval},
\[
P_{1/4}(N_*,T)\ge\frac1{\pi T}\int_{I_*}\sum_{k>N_*}|G(k;E+i/T)|^2\,dE\ge\frac{2r}{\pi T}\,c'_\eta T^{-1}e^{-hq}\ge\frac{c}{q}T^{-2}e^{-3hq/2}\ge T^{-2}e^{-2hq}
\]
for large $n$. This is~\eqref{eq:quarter}.

\emph{Nearby phases.} If $\|\varphi-1/4\|\le2|\delta|$, compare the first block of $H_\varphi$ with the same rational block $R$. The discrepancies become $(\varphi-1/4)+(k+1)\delta/q$, of size at most $3|\delta|$, so $|w_k|\le12\pi\lambda|\delta|$, and Steps~1--4 hold with modified constants. If $|\cos(2\pi q\varphi)-s|\le\eta\lambda^{-q}/16$, compare the first block, with phases $\varphi+\alpha+k\alpha$, with $A^{(p/q)}_q(E,\varphi+p/q)$; its trace $P_q(E)-2\lambda^q\cos(2\pi q\varphi)$ is within $\eta/8$ of $\Delta_t(E)$, the derivative $P_q'$ is unchanged, and the discrepancies are again $(k+1)\delta/q$. So Steps~1--4 hold with $\eta/2$ replaced by $3\eta/8$ and the same exponents.

\emph{The exponent.} At $T=T_n=|\delta_n|^{-1}=e^{b_nq}$, \eqref{eq:quarter} gives a tail at least $T_n^{-2-2h/b_n}$ beyond $N_*\ge cT_n^{1-h/(2b_n)}-1$. Since $M_p(T)\ge N^pP(N,T)$, $\beta^+(p)\ge1-h/(2b_n)-(2+2h/b_n)/p$ along even $n$ with $b_n\to\beta_e$, provided $\beta_e>h/2+(h-L)$. Let $p\to\infty$ and then $h\downarrow L$. If $\beta_e\le L/2$ there is nothing to prove; if $\beta_e=\infty$ the same argument with $b_n\to\infty$ gives $1$.
\end{proof}

The theorem gives a polynomially small tail at the critical distance, not a fixed fraction of the packet.

\subsection{Exact laws for the almost Mathieu operator}

\begin{corollary}[Exact uniform law]\label{cor:exactunif}
Let $\lambda>1$ and let $\alpha$ satisfy $\beta_e(\alpha)=\beta(\alpha)\in[L,\infty)$, for instance when all denominators with large $b_n$ are even. Then
\[
\aunif=\sup_\theta\au(H_{\lambda,\alpha,\theta})=\au(H_{\lambda,\alpha,1/4})=1-\frac{L}{2\beta(\alpha)} .
\]
\end{corollary}

\begin{proof}
Theorem~\ref{thm:quarter} gives $\au(1/4)\ge1-L/(2\beta)$, Corollary~\ref{cor:typeI} gives $\aunif\le1-L/(2\beta)$, and~\eqref{eq:unifsup} closes the chain.
\end{proof}

\begin{remark}\label{rem:aevalue}
At almost every phase the exponent is $\max\{0,1-L/\beta(\alpha)\}$~\cite[Theorem~A]{DMamo}. So under the hypothesis of Corollary~\ref{cor:exactunif}, with $\beta<\infty$, the worst phase has exactly half the typical localization strength.
\end{remark}

\begin{corollary}[The exponent at the phase $1/4$]\label{cor:quarterexact}
Let $\lambda>1$ and $\beta_e(\alpha)\ge L$. Then
\[
\au(H_{\lambda,\alpha,1/4})=\max\Bigl\{0,\ 1-\frac L{\beta_o(\alpha)},\ 1-\frac L{2\beta_e(\alpha)}\Bigr\},
\]
read as $1$ if $\beta(\alpha)=\infty$.
\end{corollary}

\begin{proof}
\emph{Lower bound.} The lower bound $\au\ge\max\{0,1-L/\beta\}$ holds at every phase~\cite[Theorem~A]{DMamo}, and Theorem~\ref{thm:quarter} gives $1-L/(2\beta_e)$. Since $\beta=\max\{\beta_e,\beta_o\}$, the larger of the two bounds is the displayed maximum. If $\beta=\infty$ the value is $1$, by~\cite[Theorem~A(2)]{DMamo} or Theorem~\ref{thm:quarter}.

\emph{Upper bound.} Let $\beta<\infty$ and $a>\max\{0,1-L/\beta_o,1-L/(2\beta_e)\}$, and run the covering argument of Theorem~\ref{thm:U} at the phase $1/4$, with $B_0^e>\beta_e$ and $B_0^o>\beta_o$ chosen so that $f(B_0^e;L,1)<a$ and $\max\{0,1-L/B_0^o\}<a$. Case~A is unchanged. In Case~B with $q_n$ even, the per-energy argument of Theorem~\ref{thm:U}, with $\om=1$ and $\ell=L$ on $\Sigma$, gives exponent at most $f(B_0^e;L,1)+o(1)<a$; here $f(b;L,1)\le\max\{1/2,1-L/(2b)\}$ and $1/2\le1-L/(2\beta_e)$ because $\beta_e\ge L$. In Case~B with $q_n$ odd, apply~\cite[Proposition~7.1]{DMamo} at the single scale $n$. Its proof uses the phase only through the sine condition $|\sin2\pi q_n\theta|\ge q_n^{-2}$ at that scale, via~\cite[Proposition~4.4 and Theorem~5.4]{DMamo}, which are stated for every phase satisfying it, and through $b_n\le B_0$, which holds for large odd $n$ with $B_0=B_0^o$. At $\theta=1/4$ and odd $q_n$, $|\sin(\pi q_n/2)|=1$. So $P_{1/4}(T^a,T)\le CT^{-m}$ in this case too.
\end{proof}

\begin{corollary}[Exceptional phases]\label{cor:exceptional}
If $\lambda>1$ and $\beta_e(\alpha)=\beta(\alpha)\in(L/2,\infty)$, then
\[
\au(H_{\lambda,\alpha,1/4})\ge1-\frac L{2\beta(\alpha)}>\max\Bigl\{0,1-\frac L{\beta(\alpha)}\Bigr\}=\au(H_{\lambda,\alpha,\theta})\quad\text{for almost every }\theta .
\]
In particular, if $e^{\beta(\alpha)}<\lambda<e^{2\beta(\alpha)}$, the phase $1/4$ has a positive transport exponent, while for almost every phase the operator has Anderson localization by the theorem of Avila, You and Zhou~\cite{AYZ2017}.
\end{corollary}

\begin{proof}
Theorem~\ref{thm:quarter} and~\cite[Theorem~A]{DMamo}.
\end{proof}

\begin{remark}[Reflection resonance and related work]\label{rem:resonance}
For even $q_n$ and odd $p_n$,
\[
2\cdot\frac14+\frac{q_n}2\,\alpha=\frac12+\frac{p_n}2+\frac{\delta_n}2\equiv\frac{\delta_n}2\pmod1 .
\] So $1/4$ is a reflection-resonant phase at every even scale, with reflection centres at distance $q_n/4$, and its phase-resonance exponent in the sense of Jitomirskaya and Liu~\cite{JLphase} is at least $2\beta_e$. The phase $1/4$ is not of the completely resonant type $2\theta\in\alpha\Z+\Z$ treated by Liu~\cite[Theorem~1.1]{Liu2022}, and that localization theorem does not apply to it. That localization is not uniform over phases or sites for the almost Mathieu operator is known in other forms: Jitomirskaya, Liu and Mi gave a sharp palindromic criterion for semi-uniform dynamical localization, with almost Mathieu examples that are dynamically localized but not semi-uniformly~\cite{JLM2024}. Corollaries~\ref{cor:exactunif}--\ref{cor:exceptional} give the corresponding exact power laws for transport.
\end{remark}

\begin{remark}[The plateau radius]\label{rem:plateau}
At even $q_n$ with $b_n>L$ and $T=|\delta_n|^{-1}$, Theorem~\ref{thm:quarter} reaches distance $e^{(b_n-h/2)q_n}$, and Proposition~\ref{prop:Camo} confines the packet, at every phase, within $e^{(b_n-L/2+\eps)q_n}$. So the fastest-tail scale at such a scale is matched up to factors $e^{o(q_n)}$. The mass at that scale is only shown to be polynomially small; whether a fixed fraction of the packet reaches it is open.
\end{remark}

\subsection{Odd-harmonic and type I potentials}

The proof of Theorem~\ref{thm:quarter} uses the almost Mathieu operator only through a wide band at one local phase. We isolate that input. For a real-analytic supercritical $v$, $\ell_0>0$, $\xi>0$, $\eta\in(0,1]$ and a set $\mathcal N'$ of indices, consider:

\begin{itemize}
\item[$(\mathrm W_v)$] for every $n\in\mathcal N'$ there are $x'\in\T$ and a real $E'$ with $\dist(E',\Sigma)\le e^{-\xi q_n}$, $\ell(E')\le\ell_0+\xi$, $|t_{q_n}(x',E')|\le2-\eta$ and $|\partial_Et_{q_n}(x',E')|\le e^{(\ell(E')/2+\xi)q_n}$.
\end{itemize}

\begin{theorem}[A lower bound from a wide band]\label{thm:Wv}
Let $v$ be real analytic with $H_{v,\alpha}$ supercritical, and assume $(\mathrm W_v)$ along an infinite set $\mathcal N'$ with $b_n\to\beta'$ as $n\to\infty$ in $\mathcal N'$, where $\beta'>\ell_0/2$. There are constants $c,C>0$, depending only on $v$, $\alpha$, $\ell_0$ and $\eta$, such that the following holds if $\xi$ is small enough. For every phase $\theta$ and every large $n\in\mathcal N'$ there is a site $x_n$ with $|x_n|\le|\delta_n|^{-1}$ and
\[
P_{\theta+x_n\alpha}\bigl(cT_n^{1-(\ell_0/2+C\xi)/b_n},\,T_n\bigr)\ge T_n^{-C},\qquad T_n=|\delta_n|^{-1}.
\]
Hence $\aunif\ge1-\ell_0/(2\beta')-C\xi$.
\end{theorem}

\begin{proof}
This is the proof of Theorem~\ref{thm:quarter} with the following changes. Restrict to $n\in\mathcal N'$ with $b_n>\ell_0/2+C\xi$; otherwise the claimed exponent is not positive. The rational block is the $p_n/q_n$-periodic block at the phase $x'/q_n$, with trace $t_{q_n}(x',E')$ of modulus at most $2-\eta$, and~\cite[Lemma~8.2(1)]{DMamo} gives $S_R\le Ce^{(\ell(E')/2+\xi)q_n}$. The energy window around $E'$ has width $e^{-(\ell(E')/2+C\xi)q_n}/(Cq_n^2)$; on it the bound $|\partial_E^2t_q|\le8q^2e^{q(\ell(E')+\xi)}$~\cite[Corollary~5.6]{DMuniv} keeps the trace elliptic with a fixed margin and the derivative of size $e^{(\ell(E')/2+C\xi)q}$. The local phases $q_n(\theta+j\alpha)=q_n\theta+j\delta_n$ sweep $\T$ as $j$ ranges over an interval of length $|\delta_n|^{-1}$, so there is a site $x_n$, $|x_n|\le|\delta_n|^{-1}$, at which the first block to the right has local phase within $|\delta_n|$ of $x'$. Closeness of local phases identifies the block only up to a cyclic shift of the rational potential; cyclic shifts have the same trace and the same energy derivative, and uniform growth controls their frames. The insertion error is at most $\exp(C\|v'\||\delta_n|S_R)-1\to0$, since $b_n>\ell(E')/2+2\xi$. The elliptic run lemma~\cite[Lemma~9.5]{DMamo} holds for Lipschitz potentials, $\|A_{q}\|\le Ce^{(\ell(E')+\xi)q}$ by Lemma~\ref{lem:growthy}, and the rest of the proof is unchanged. The last statement follows as in Section~\ref{sec:prelim}: if $\sup_\varphi P_\varphi(T^a,T)=O(T^{-m})$ for all $m$, then $a\ge1-(\ell_0/2+C\xi)/\beta'$.
\end{proof}

\begin{theorem}[Odd-harmonic potentials]\label{thm:oddharm}
Let $v$ be real analytic and odd-harmonic, with $H_{v,\alpha}$ supercritical. Then $0\in\Sigma$, and
\[
\aunif\ge\max\Bigl\{0,\ 1-\frac{\ell(0)}{2\beta_e(\alpha)}\Bigr\},
\]
read as $1$ if $\beta_e(\alpha)=\infty$. More precisely, for every $\eps>0$ there is $C$ such that for every phase $\theta$ and all large even $q_n$ with $b_n\ge\ell(0)/2+\eps$ there is a site $x_n$, $|x_n|\le|\delta_n|^{-1}$, with $P_{\theta+x_n\alpha}(T_n^{1-(\ell(0)/2+\eps)/b_n},T_n)\ge T_n^{-C}$.
\end{theorem}

\begin{proof}
\emph{$0\in\Sigma$.} Reduced fractions $p/q$ with even $q$ approach $\alpha$. By Proposition~\ref{prop:G}, $0$ lies in the spectrum of the $p/q$-periodic operator at some phase, and every point of the phase-union rational spectrum is within $C|\alpha-p/q|^{1/2}$ of $\Sigma$~\cite[Lemma~2.7(1)]{DMuniv}.

\emph{$(\mathrm W_v)$ at even scales.} Let $q=q_n$ be even and $\varphi_0$ the phase of Proposition~\ref{prop:G}(2), and write $\Delta=\Delta_{\varphi_0}$ and $s=(-1)^{q/2}$. Then $\Delta(0)=2s$, $\Delta$ is even, and the bands of the periodic operator adjacent to $0$ lie on both sides of $0$, with $s\Delta\le2$ on them. Let $f=2-s\Delta\ge0$ on these bands, so $f(0)=f'(0)=0$. Let $E^{(1)}>0$ be the first point with $f(E^{(1)})=\eta$. Every band of a $q$-periodic operator with hopping one has width at most $2\pi/q$~\cite[proof of Lemma~5.4]{DMuniv}, so $E^{(1)}\to0$, $\ell(E^{(1)})\to\ell(0)$, and by~\cite[Lemma~2.7(1)]{DMuniv} the interval $[0,E^{(1)}]$ is within $C(|\delta|/q)^{1/2}\le e^{-\xi q}$ of $\Sigma$ when $b_n$ is bounded below and $\xi$ is small. On it~\cite[Corollary~5.6]{DMuniv} gives $|f''|\le M=8q^2e^{q(\ell(0)+\xi)}$ for large $q$. With $t=f'(E^{(1)})/M\le E^{(1)}$, Taylor's formula on $[E^{(1)}-t,E^{(1)}]$ gives $|\Delta'(E^{(1)})|^2\le2M\eta$. So $(\mathrm W_v)$ holds at every large even $q_n$ with $x'=q\varphi_0$, $E'=E^{(1)}$, $\ell_0=\ell(0)+o(1)$ and slack $\xi$, after relabelling $\xi$. Theorem~\ref{thm:Wv}, applied along even scales with $b_n\to\beta_e$, and $\xi\to0$, give the claim. If $\beta_e\le\ell(0)/2$ there is nothing to prove.
\end{proof}

The sites $x_n$ may depend on the scale, so Theorem~\ref{thm:oddharm} bounds the uniform exponent, not the exponent of a prescribed phase; Section~\ref{sec:category} proves the same lower bound on a residual set of phases (Corollary~\ref{cor:G2}). For supercritical odd-harmonic type~I potentials the lower and upper bounds coincide when $\ell(0)=\ell_*$ and $\beta_e=\beta$ is finite and at least $\ell_*$; then the uniform exponent and the exponent on a residual set of phases both equal $1-\ell_*/(2\beta)$ (Corollaries~\ref{cor:oddbracket} and~\ref{cor:G2}).

\begin{corollary}[A bracket for odd-harmonic type I potentials]\label{cor:oddbracket}
Let $v$ be real analytic and odd-harmonic, with $H_{v,\alpha}$ supercritical and type~I, and $\beta_e(\alpha)=\beta(\alpha)\in[\ell_*,\infty)$. Then
\[
\max\Bigl\{0,\ 1-\frac{\ell(0)}{2\beta}\Bigr\}\le\aunif\le1-\frac{\ell_*}{2\beta},
\]
and the two sides coincide when $\ell(0)=\ell_*$.
\end{corollary}

\begin{proof}
Theorem~\ref{thm:oddharm} and Corollary~\ref{cor:typeI}.
\end{proof}

For the almost Mathieu operator $\ell(0)=\ell_*=L$, and Corollary~\ref{cor:oddbracket} recovers the uniform half of Corollary~\ref{cor:exactunif}. The class is larger than the almost Mathieu family: for instance $4\cos2\pi x+2\mu\cos6\pi x$ is odd-harmonic, and it is supercritical and type~I at every irrational frequency when $2|\mu|\cosh(3\pi/2)\le1/4$, by Theorem~\ref{thm:nbhd}.

The upper bound is local in energy for smooth filters. For a vector $\psi$ and $g\in C_c^\infty(\R)$ we use the unnormalized packet $g(H_\theta)\psi$ and the quantities~\eqref{eq:abel}.

\begin{corollary}[Smooth filters]\label{cor:filter}
Let $v$ be real analytic with $H_{v,\alpha}$ supercritical and type~I, $\beta(\alpha)<\infty$, $\ell_0>0$, and let $g\in C_c^\infty(\R)$ be real with $\ell(E)\ge\ell_0$ for every $E\in\Sigma\cap\supp g$. Then for every $s'>f(\beta;\ell_0,1)$ and $m>0$,
\[
\sup_\theta P_{g(H_\theta)\delta_0}(T^{s'},T)=O(T^{-m}).
\]
Consequently $\alpha^+_{u,g(H_\theta)\delta_0}\le f(\beta;\ell_0,1)$ for every $\theta$ with $g(H_\theta)\delta_0\ne0$.
\end{corollary}

\begin{proof}
Fix $\eps\in(0,\ell_0/2)$ and an open neighbourhood $U$ of $\Sigma\cap\supp g$ with $\ell\ge\ell_0-\eps$ on $U$. Choose $\chi\in C_c^\infty(U)$ with $\chi=1$ near $\Sigma\cap\supp g$. Then $(1-\chi)g$ vanishes on $\Sigma=\sigma(H_\theta)$, so $g(H_\theta)=(\chi g)(H_\theta)$; we may therefore assume $\supp g\subset U$. Fix $d_0>0$ such that the $2d_0$-neighbourhood of $\supp g$ lies in $U$.

Write $R_E=(H_\theta-E-i/T)^{-1}$ and $\psi=g(H_\theta)\delta_0$, so that $R_E\psi=g(H_\theta)R_E\delta_0$, and use~\eqref{eq:parseval} for $\psi$.

\emph{Energies with $\dist(E,\supp g)\ge d_0$.} Here $R_E\psi=h_E(H_\theta)\delta_0$ with $h_E(x)=g(x)/(x-E-i/T)$. These functions are smooth, supported in $\supp g$, and all their derivatives are bounded uniformly in $\theta$, $T\ge1$ and such $E$, by $C/(1+|E|)$. By~\cite[Lemma~2.8]{DMuniv}, $|\langle\delta_j,h_E(H_\theta)\delta_0\rangle|\le C_r(1+|E|)^{-1}|j|^{-r}$ for every $r$, so these energies contribute $O(T^{-m})$ at radius $T^{s'}$.

\emph{The remaining energies.} If $\dist(E,\supp g)<d_0$, then $E\in U$. By~\cite[Lemma~2.8]{DMuniv},
\[
\bigl\|\un_{\{|j|>N\}}\,g(H_\theta)\,\un_{\{|j|\le N/2\}}\bigr\|\le C_rN^{-r}
\]
uniformly in $\theta$, so
\[
\|\un_{\{|j|>N\}}R_E\psi\|^2\le2C_r^2N^{-2r}\|R_E\delta_0\|^2+2\|g\|_\infty^2\|\un_{\{|j|>N/2\}}R_E\delta_0\|^2.
\]
The first term integrates to at most $2C_r^2N^{-2r}$ by~\eqref{eq:parseval}. For the second, run the proof of Theorem~\ref{thm:U} at radius $N/2$ with $N=T^{s'}$, restricted to energies in $U$: far energies are handled by Combes--Thomas, and near energies by the per-energy minimum of the velocity and confinement radii. For a type~I potential, $\om_{y_0}=1$ on $\Sigma$ for small $y_0$~\cite[Lemma~5.2]{DMuniv}, so by Lemma~\ref{lem:omy}(1) the quantity $\Omega$ of that proof is close to $1$ on $\Sigma_\rho$ once $\rho$ and the slack parameters are small. The per-energy exponent is then at most $f(B_0;\ell_0-\eps-10a,\Omega)$, which is below $s'$ once $\eps$, $B_0-\beta$ and the slack parameters are small. This gives the bound, uniformly in $\theta$. The last statement follows from~\cite[Lemma~2.3]{DMuniv}, since the normalized packet has finite moments by~\cite[Lemma~2.8]{DMuniv}.
\end{proof}

For a rough filter, merely supported where $\ell\ge\ell_0$, the statement can fail: the filtered vector need not have rapidly decaying spatial tails. Smoothness is used exactly in the off-diagonal bounds above.

\begin{corollary}[An exact uniform law at the closed gap]\label{cor:window}
Let $v$ be real analytic and odd-harmonic, with $H_{v,\alpha}$ supercritical and type~I, and $\beta_e(\alpha)=\beta(\alpha)\in[\ell(0),\infty)$. For $w>0$ let $g_w\in C_c^\infty((-w,w))$ be real with $g_w=1$ near $0$, and let
\[
\gamma_w=\inf\Bigl\{a\ge0:\ \sup_\theta P_{g_w(H_\theta)\delta_0}(T^a,T)=O(T^{-m})\text{ for every }m>0\Bigr\}.
\]
Then
\[
1-\frac{\ell(0)}{2\beta}\le\gamma_w\le f\Bigl(\beta;\min_{\Sigma\cap[-w,w]}\ell,\,1\Bigr),\qquad\lim_{w\downarrow0}\gamma_w=1-\frac{\ell(0)}{2\beta}.
\]
\end{corollary}

\begin{proof}
The upper bound is Corollary~\ref{cor:filter}. As $w\downarrow0$, $\min_{\Sigma\cap[-w,w]}\ell\to\ell(0)$, since $0\in\Sigma$ and $\ell$ is continuous, and $f(\beta;\ell(0),1)=1-\ell(0)/(2\beta)$ because $\beta\ge\ell(0)$. For the lower bound, the energy windows $I_*$ in the proof of Theorem~\ref{thm:oddharm} shrink to $0$, so for large $n$ they lie in a closed interval $J_w$ with $g_w=1$ on a neighbourhood of $J_w$. For $E\in J_w$, $0<\eta\le1$, the functions $x\mapsto\chi(x)(1-g_w(x))/(x-E-i\eta)$, with $\chi$ smooth and equal to $1$ on a neighbourhood of the spectrum, have all derivatives bounded uniformly in $E$ and $\eta$, because the numerator vanishes near $J_w$. So~\cite[Lemma~2.8]{DMuniv} gives $\|\un_{\{|j|>N\}}R_E(1-g_w)(H_\theta)\delta_0\|\le C_rN^{-r}$, uniformly in $\theta$, $E\in J_w$ and $T\ge1$. Since $\|a-b\|^2\ge\frac12\|a\|^2-\|b\|^2$, integrating over $I_*$ as in Step~4 of Theorem~\ref{thm:quarter} gives $P_{g_w(H_{\theta+x_n\alpha})\delta_0}(N_n,T_n)\ge\frac12T_n^{-C}-C_rN_n^{-2r}$, where $N_n=T_n^{1-(\ell(0)/2+\eps)/b_n}\ge T_n^{1/4}$ for large $n$ and $\eps<\ell(0)/4$, since $b_n\to\beta\ge\ell(0)$. Choosing $r$ large gives $\sup_\theta P_{g_w(H_\theta)\delta_0}(N_n,T_n)\ge\frac14T_n^{-C}$, and letting $\eps\to0$ gives the lower bound.
\end{proof}

At a fixed width $w>0$ we do not claim equality. The corollary says that the uniform law is exact, energy by energy, at the energy where a gap closes at every even scale.

\subsection{Non-primitive potentials attain the factor \texorpdfstring{$2\om$}{2 omega}}

\begin{proposition}\label{prop:NP}
Let $\lambda>1$, $s\ge2$ an integer and $v(x)=2\lambda\cos(2\pi sx)$. Then $\ell=\log\lambda$ and $\om=s$ on $\Sigma$. For every $\beta_*\in(\log\lambda,\infty)$ there is an irrational $\alpha$ with $\beta(\alpha)=\beta_*$ and
\[
\aunif=\sup_\theta\au(H_{v,\alpha,\theta})=1-\frac{\log\lambda}{2s\beta(\alpha)}=1-\frac{\ell_*}{2\om^+\beta(\alpha)} .
\]
\end{proposition}

So the bound of Theorem~\ref{thm:U} is attained in the class of all analytic potentials: no bound of the form $1-\ell_*/(M\beta)$ with $M<2\om^+$ holds for all of them.

\begin{proof}
By~\cite[Remark~5.30]{DMuniv}, $L_v(E,y)=L_w(E,sy)$, where $w(x)=2\lambda\cos2\pi x$ is taken at frequency $s\alpha$; so $\ell=\log\lambda$ and $\om=s$ on $\Sigma$. The upper bound is Theorem~\ref{thm:U}, since $\beta_*\ge\ell^+=\log\lambda$.

\emph{The frequency.} Let $m=2s$. We build $\alpha=[0;a_1,a_2,\dots]$ so that every large partial quotient follows a denominator divisible by $m$. Take $a_1=m$, so $q_1=m$. Whenever $q_n\equiv0\pmod m$, let $a_{n+1}=\lceil e^{\beta_*q_n}\rceil$; then $q_{n+1}=a_{n+1}q_n+q_{n-1}\equiv q_{n-1}\pmod m$, which is invertible modulo $m$ because $\gcd(q_{n-1},q_n)=1$. Then choose a bounded $a_{n+2}\in\{1,\dots,m\}$ with $a_{n+2}q_{n+1}+q_n\equiv0\pmod m$, which is possible because $q_{n+1}$ is invertible modulo $m$; so $q_{n+2}\equiv0\pmod m$, and repeat. The partial quotients that adjust the congruence are bounded by $m$, so $b_n\to0$ along the indices with $m\nmid q_n$, while $b_n\to\beta_*$ along those with $m\mid q_n$. Hence $\beta(\alpha)=\beta_*$.

\emph{The lower bound.} By~\cite[proof of Theorem~5.29]{DMuniv}, $H_{v,\alpha,\theta}=H_{\lambda,s\alpha,s\theta}$. At each $n$ with $m\mid q_n$, $q'=q_n/s$ is a convergent denominator of $s\alpha$ with $q'(s\alpha)-p_n=\delta_n$~\cite[Proposition~5.27]{DMuniv}, so its ratio is $sb_n$, and $q'$ is even. Hence $\beta_e(s\alpha)\ge s\beta(\alpha)$; with $\beta(s\alpha)\le s\beta(\alpha)$~\cite[Lemma~5.28]{DMuniv}, $\beta_e(s\alpha)=\beta(s\alpha)=s\beta(\alpha)$. Theorem~\ref{thm:quarter} at frequency $s\alpha$ and phase $1/4$ gives $\au(H_{\lambda,s\alpha,1/4})\ge1-\log\lambda/(2s\beta(\alpha))$, and $H_{\lambda,s\alpha,1/4}=H_{v,\alpha,1/(4s)}$. The equalities follow from~\eqref{eq:unifsup}.
\end{proof}

The mechanism is a symmetry: translation of $v$ by $1/s$ forces clusters of $2s$ zeros of the discriminant, at the relevant exponential scale, at every scale whose denominator is divisible by $s$; the zeros in a cluster need not coincide.

\section{Measure versus category}\label{sec:category}

The lower bounds of Section~\ref{sec:sharp} depend on the phase only through the local phase $q_n\theta\bmod1$ at the scale used, and they are stable under small changes of it. A standard Baire argument then turns them into lower bounds on a residual set of phases.

\begin{theorem}[A Baire principle]\label{thm:G}
Let $v$ be continuous and $\alpha$ irrational. Suppose there are scales $n_j\to\infty$, times $T_j\to\infty$, radii $N_j$, a constant $C>0$, a number $a_*\in[0,1]$ with $\log N_j/\log T_j\to a_*$, and nonempty open sets $U_j\subset\T$ such that
\[
P_\theta(N_j,T_j)\ge T_j^{-C}\qquad\text{for every }\theta\text{ with }q_{n_j}\theta\bmod1\in U_j .
\]
Then $\{\theta:\au(H_{v,\alpha,\theta})\ge a_*\}$ contains a dense $G_\delta$ set.
\end{theorem}

\begin{proof}
The map $\theta\mapsto H_\theta$ is norm continuous, so for fixed $N,T$ the map $\theta\mapsto P_\theta(N,T)=1-\sum_{|j|\le N}a_\theta(j,T)$ is continuous, by dominated convergence in the Abel integral. The set $V_j=\{\theta:P_\theta(N_j,T_j)>T_j^{-C}/2\}$ is open and contains $\{\theta:q_{n_j}\theta\in U_j\}$, a nonempty open $1/q_{n_j}$-periodic set, which meets every interval of length greater than $1/q_{n_j}$. So $W_m=\bigcup_{j\ge m}V_j$ is open and dense, and $G=\bigcap_mW_m$ is a dense $G_\delta$ set. For $\theta\in G$, $\theta\in V_j$ for infinitely many $j$, and $M_p(T_j)\ge N_j^pP_\theta(N_j,T_j)\ge N_j^pT_j^{-C}/2$. Hence $\beta^+(p)\ge a_*-C/p$, and letting $p\to\infty$ gives $\au(\theta)\ge a_*$.
\end{proof}

The argument is the standard open-dense mechanism; we claim no novelty for it. Its content lies in the inputs.

\begin{corollary}[The almost Mathieu operator]\label{cor:G1}
Let $\lambda>1$ and $\alpha$ irrational with $\beta_e(\alpha)>L/2$. Then $\au(H_{\lambda,\alpha,\theta})\ge1-L/(2\beta_e(\alpha))$ for every $\theta$ in a dense $G_\delta$ set, read as $1$ if $\beta_e=\infty$. If moreover $\beta_e(\alpha)=\beta(\alpha)\in[L,\infty)$, then
\[
\au(H_{\lambda,\alpha,\theta})=1-\frac L{2\beta(\alpha)}\quad\text{on a dense }G_\delta\text{ set},
\]
which is strictly larger than the value $1-L/\beta(\alpha)$ at almost every phase.
\end{corollary}

\begin{proof}
If $\beta_e=\infty$ the value is $1$ at every phase by Remark~\ref{rem:endpoint}. Let $L/2<\beta_e<\infty$ and fix $h>L$ with $h/2+(h-L)<\beta_e$ and $\eta=1/4$. Along even scales $n_j$ with $b_{n_j}\to\beta_e$, Theorem~\ref{thm:quarter} gives~\eqref{eq:quarter} at $T_j=|\delta_{n_j}|^{-1}$ for every $\theta$ whose local phase lies in the open nonempty set $U_j=\{x\in\T:|\cos2\pi x-(-1)^{q_{n_j}/2}|<\eta\lambda^{-q_{n_j}}/16\}$. There $T_j^{-2}e^{-2hq_{n_j}}\ge T_j^{-C}$ with $C=2+2h/(h/2+h-L)$, and $\log N_j/\log T_j\to1-h/(2\beta_e)$. Theorem~\ref{thm:G} gives a dense $G_\delta$ set on which $\au\ge1-h/(2\beta_e)$. Intersect these sets along a sequence $h\downarrow L$; a countable intersection of dense $G_\delta$ sets is a dense $G_\delta$ set. If $\beta_e=\beta\ge L$, Corollary~\ref{cor:typeI} gives $\au\le1-L/(2\beta)$ at every phase. The almost-everywhere value is~\cite[Theorem~A]{DMamo}; for finite $\beta\ge L$ the two values differ.
\end{proof}

So for these frequencies the transport exponent of the almost Mathieu operator takes one value at almost every phase and a strictly larger value at a residual set of phases. At $\beta=\infty$ both values are $1$.

\begin{corollary}[Odd-harmonic potentials]\label{cor:G2}
Let $v$ be real analytic and odd-harmonic, with $H_{v,\alpha}$ supercritical, and $\beta_e(\alpha)>\ell(0)/2$. Then $\au(H_{v,\alpha,\theta})\ge1-\ell(0)/(2\beta_e(\alpha))$ on a dense $G_\delta$ set of phases. If moreover $H_{v,\alpha}$ is type~I and $\ell_*\le\beta(\alpha)<\infty$, then $\au\le1-\ell_*/(2\beta(\alpha))$ at every phase, and $\au=\max\{0,1-\ell_*/\beta(\alpha)\}$ at almost every phase~\cite[Theorem~A]{DMuniv}.
\end{corollary}

\begin{proof}
In the proof of Theorem~\ref{thm:oddharm}, the lower bound at scale $n$ holds for the packet launched at any site whose first block has local phase within $|\delta_n|$ of $x'_n=q_n\varphi_0$, with the same constants. Applied to the launch site $0$, it holds for every $\theta$ with $q_n\theta$ in a fixed nonempty open neighbourhood $U_n$ of $x'_n-q_n\alpha$, because the insertion error is controlled uniformly there. Theorem~\ref{thm:G} and a countable intersection give the lower bound. The upper bound is Corollary~\ref{cor:typeI}.
\end{proof}

The lower and upper bounds of Corollary~\ref{cor:G2} coincide only when $\ell(0)=\ell_*$ and $\beta_e=\beta$. We do not claim that the worst exponent is attained on a residual set for every odd-harmonic type~I potential.

\begin{corollary}[Non-primitive potentials]\label{cor:G3}
Let $v(x)=w(sx)$ with an integer $s\ge2$. If the conclusion of Corollary~\ref{cor:G1} or~\ref{cor:G2} holds for $w$ at frequency $s\alpha$, with $\beta_e(s\alpha)$ and $\beta(s\alpha)$ in place of $\beta_e(\alpha)$ and $\beta(\alpha)$, then it holds for $v$ at frequency $\alpha$ with the same numbers.
\end{corollary}

\begin{proof}
$H_{v,\alpha,\theta}=H_{w,s\alpha,s\theta}$~\cite[Theorem~5.29]{DMuniv}, and $\theta\mapsto s\theta$ is a continuous, open surjection of $\T$, so preimages of dense $G_\delta$ sets are dense $G_\delta$ sets and preimages of null sets are null.
\end{proof}

\section{Light cones in the XY chain}\label{sec:xy}

\subsection{The chain}

Consider the spin chain on $\Z$ with Hamiltonian
\[
H_{XY}=\sum_n\bigl(X_nX_{n+1}+Y_nY_{n+1}\bigr)+\sum_nh_nZ_n,\qquad h_n=v(\theta+n\alpha),
\]
where $X_n,Y_n,Z_n$ are the Pauli matrices at site $n$. The interactions are nearest-neighbour and bounded, so the dynamics $\tau_t$ exists on the quasi-local algebra, and $\tau_t(A)=\lim_{L\to\infty}\tau^L_t(A)$ in norm for local $A$, where $\tau^L$ is generated by the restriction of $H_{XY}$ to $\Lambda_L=[-L,L]$~\cite{LiebRobinson1972}. All identities below are proved in finite volume and pass to the limit.

On $\Lambda_L$ the Jordan--Wigner transformation~\cite{LSM1961} $c_j=(\prod_{-L\le k<j}Z_k)\sigma^-_j$, $\sigma^-=(X-iY)/2$, gives $c_j^*c_j=(1+Z_j)/2$ and $H_{XY}|_{\Lambda_L}=\sum_{a,b}c_a^*(M_L)_{ab}c_b+\mathrm{const}$ with $M_L=-2(S+S^*)+2\operatorname{diag}(h)$ on $\ell^2(\Lambda_L)$. The diagonal unitary $W=\operatorname{diag}((-1)^j)$ gives $WM_LW=2H_L$, where $H_L$ is the Dirichlet restriction of $H=H_{v,\alpha,\theta}$. Hence
\[
\tau^L_t(c_x)=\sum_yU^L_t(x,y)c_y,\qquad|U^L_t(x,y)|=\bigl|\langle\delta_y,e^{-2itH_L}\delta_x\rangle\bigr| .
\]
For the Majorana operators $\gamma_{2x-1}=c_x+c_x^*$ and $\gamma_{2x}=i(c_x-c_x^*)$, $\tau_t(\gamma_{2x-1})$ and $\tau_t(\gamma_{2x})$ are real linear combinations of Majoranas whose total squared weight on the sites of any set $R$ is $\sum_{y\in R}|U_t(x,y)|^2$. Spin time $t$ therefore corresponds to one-body time $2t$, and we write
\[
p^{(x)}_s(R)=\sum_{y\in R}\bigl|\langle\delta_y,e^{-isH}\delta_x\rangle\bigr|^2
\]
for the one-body mass in $R$ at one-body time $s$. These formulas are read in the CAR algebra in finite volume; the even polynomials in Majoranas and the local spin observables used below have norm limits as $L\to\infty$.

\subsection{Exact commutator identities}

\begin{proposition}\label{prop:comm}
Let $R$ be a finite set of sites not containing $0$, $P_R=\prod_{k\in R}Z_k$ and $p=p^{(0)}_{2t}(R)$. Then
\[
\bigl\|[\tau_t(Z_0),P_R]\bigr\|=4\sqrt{p(1-p)} .
\]
In particular $\|[\tau_t(Z_0),Z_N]\|=4|u|\sqrt{1-|u|^2}$ with $u=\langle\delta_N,e^{-2itH}\delta_0\rangle$.
\end{proposition}

\begin{proof}
$\tau_t(n_0)=f^*f$ with $n_0=c_0^*c_0$ and $f=\sum_yU_t(0,y)c_y$, a normalized mode. Since $Z_k=2n_k-1=-(-1)^{n_k}$, $P_R=\pm(-1)^{N_R}$ with $N_R=\sum_{k\in R}n_k$, and conjugation by $P_R$ changes the sign of $c_y$ for $y\in R$ and fixes the others. Write $f=f_{\mathrm{in}}+f_R$ with $f_R$ the part supported in $R$. Then $P_RfP_R=f_{\mathrm{in}}-f_R$ and
\[
[f^*f,P_R]=2\bigl(f_{\mathrm{in}}^*f_R+f_R^*f_{\mathrm{in}}\bigr)P_R .
\]
Write $f_R=ag_1$ and $f_{\mathrm{in}}=bg_2$ with normalized modes $g_1,g_2$ of disjoint supports, $|a|^2=p$ and $|b|^2=1-p$. The operator $\bar ba\,g_2^*g_1+\bar ab\,g_1^*g_2$ vanishes on the zero- and two-particle sectors of $g_1,g_2$ and acts on the one-particle sector as $\bigl(\begin{smallmatrix}0&\bar ab\\a\bar b&0\end{smallmatrix}\bigr)$, so its norm is $|a||b|=\sqrt{p(1-p)}$. Since $P_R$ is unitary and $Z_0=2n_0-1$, the claim follows.
\end{proof}

\begin{proposition}[Infinite-temperature OTOC]\label{prop:otoc}
In finite volume $\Lambda\supset R\cup\{0\}$, with $C_R(t)=\frac12\,2^{-|\Lambda|}\tr\bigl(|[\tau_t(Z_0),P_R]|^2\bigr)$ and $p$ as in Proposition~\ref{prop:comm},
\[
C_R(t)=4p(1-p).
\]
\end{proposition}

\begin{proof}
From the proof of Proposition~\ref{prop:comm}, $|[\tau_t(Z_0),P_R]|^2=16\,|a|^2|b|^2\bigl(n_{g_1}(1-n_{g_2})+n_{g_2}(1-n_{g_1})\bigr)$, conjugated by $P_R$. In the tracial state, $g_1,g_2$ are orthonormal modes and $\langle n_{g_1}(1-n_{g_2})\rangle=1/4$.
\end{proof}

So an order-one commutator, or OTOC, between $Z_0$ and the parity of a far region is equivalent to a fixed fraction, bounded away from $0$ and $1$, of the one-body packet lying in that region.

\begin{remark}\label{rem:otoc}
The tracial OTOC for a single receiver site is contained in the general Gaussian-state formula of Riddell and S{\o}rensen~\cite[Appendix~C, eq.~(C33)]{RiddellSorensen2020} (numbering of the PDF of arXiv:1908.03292v3), whose tracial specialization is $1-4p(1-p)$. The reduction of XY dynamics to one-body dynamics is classical~\cite{LSM1961}; Hamza, Sims and Stolz used it to derive dynamical localization of disordered XY chains from one-body localization~\cite{HSS2012}.
\end{remark}

\subsection{From Abel averages to every time}

Let $K$ be a Jacobi operator on $\ell^2(\Z)$ with off-diagonal entries $1$ and real diagonal, $\psi_t=e^{-itK}\delta_x$, and $\mathfrak f(t)=\|\un_R\psi_t\|$ with $R=\{y:|y-x|>N\}$.

\begin{lemma}\label{lem:abel}
Let $N\ge1$. Then $\mathfrak f$ is $1$-Lipschitz, and for $T\ge1$ and $|t|\le T$,
\[
\mathfrak f(t)^2\le\bigl(4e^3T\,P(N,T)\bigr)^{2/3},\qquad P(N,T)=\frac2T\int_0^\infty e^{-2s/T}\mathfrak f(s)^2\,ds .
\]
\end{lemma}

\begin{proof}
$\frac d{dt}\mathfrak f^2=\langle\psi_t,i[K,\un_R]\psi_t\rangle$, and $[K,\un_R]$ lives on the two bonds crossing the boundary of $R$. Since $N\ge1$, the two inner boundary sites $x\pm N$ are distinct, and Cauchy--Schwarz gives $|(\mathfrak f^2)'|\le2\mathfrak f$, so $\mathfrak f$ is $1$-Lipschitz. If $\mathfrak f(t_0)=\sigma\in(0,1]$ with $0\le t_0\le T$, then $\mathfrak f\ge\sigma/2$ on $[t_0,t_0+\sigma/2]$, and $P(N,T)\ge\frac2Te^{-2(t_0+1/2)/T}\frac{\sigma^2}4\frac\sigma2\ge e^{-3}\sigma^3/(4T)$. Negative times follow from $e^{itK}\delta_x=\overline{e^{-itK}\delta_x}$, as $K$ is real.
\end{proof}

For $N=0$ the Lipschitz constant is $\sqrt2$ instead of $1$, which changes only the numerical constant. Splitting the Abel integral at $kT$ gives, for every $k>0$,
\begin{equation}\label{eq:abelsplit}
P(N,T)\le\sup_{0\le t\le kT}\mathfrak f(t)^2+e^{-2k}.
\end{equation}

\subsection{Locality for even observables}

An observable is \emph{even} if it commutes with the spin flip $\prod_kZ_k$ over its support.

\begin{lemma}[Locality]\label{lem:locality}
Let $A$ be even, supported in $[-l,l]$, with $\|A\|\le1$, and let $\|B\|\le1$.
\begin{enumerate}
\item Let $R$ be disjoint from $[-l,l]$ such that $R_-=R\cap(-\infty,-l)$ and $R_+=R\cap(l,\infty)$ are intervals, and let $B$ lie in the algebra generated by the even observables supported in $R_-$ and those supported in $R_+$; for instance a parity string $P_R$, or a product $B_-B_+$ of even factors. Then
\[
\|[\tau_t(A),B]\|\le2^{4l+4}(4l+2)\max_{|z|\le l}\sqrt{p^{(z)}_{2t}(R)} .
\]
\item Let $B$ be any observable supported in $[r,\infty)$ with $r>l$. Then the same bound holds with $R$ replaced by the half-line $[r,\infty)$. The mirror statement holds on the left.
\end{enumerate}
\end{lemma}

\begin{proof}
In case~2 use Jordan--Wigner strings from the right end of the volume; the one-body matrix changes only by diagonal signs, so the moduli above are unchanged. Then $B$, of any parity, lies in the CAR algebra generated by the modes of $[r,\infty)$, because the strings of its sites run to the right end inside $[r,\infty)$. In case~1 either representation works: an even spin observable supported in an interval is an even element of the CAR algebra of that interval, since its strings cancel in pairs inside the interval; so $B$ is an even element of the CAR algebra of $R$. In both cases $A$ is an even polynomial in the $2(2l+1)$ Majoranas of $[-l,l]$: $A=\sum_Sc_S\gamma_S$ over even subsets $S$, with at most $2^{4l+2}$ terms and $|c_S|\le\|A\|$. For $\mu\in S$ split $G_\mu=\tau_t(\gamma_\mu)=G^{\mathrm{in}}_\mu+G^{\mathrm{out}}_\mu$ into its parts on Majoranas outside and inside $R$ (respectively the half-line). These are real linear combinations, so $\|G^{\mathrm{in}}_\mu\|\le1$ and $\|G^{\mathrm{out}}_\mu\|=\sqrt{p^{(z)}_{2t}(R)}$ for the site $z$ of $\mu$. A product of an even number of linear elements supported outside $R$ commutes with every even element of the CAR algebra of $R$, and, in case~2, with every element. So $[\tau_t(\gamma_S),B]=[\prod_\mu G_\mu-\prod_\mu G^{\mathrm{in}}_\mu,B]$, and telescoping gives $\|\prod G_\mu-\prod G^{\mathrm{in}}_\mu\|\le\sum_\mu\|G^{\mathrm{out}}_\mu\|$. Sum over $S$.
\end{proof}

Case~2 needs the whole half-line: for an odd observable $B$, the Jordan--Wigner string of $B$ runs through every site between $B$ and the end of the chain, and a bound through the mass on the support of $B$ alone is not available. The bounds concern even source observables and the receiver classes above; they are not the usual exponential Lieb--Robinson bounds for arbitrary observables.

\subsection{Light-cone exponents}

For a site $x$ and $N\ge0$ let $\mathcal B(x,N)$ be the set of observables $B$ with $\|B\|\le1$ of the two kinds in Lemma~\ref{lem:locality}, relative to the site $x$: elements of the algebra generated by even observables of $R_-$ and of $R_+$, for sets $R\subset\{y:|y-x|>N\}$ with $R_\pm$ intervals on the two sides of $x$; and arbitrary observables supported in $[x+N+1,\infty)$ or in $(-\infty,x-N-1]$. Let $\mathcal A_l(x)$ be the set of even observables supported in $[x-l,x+l]$ with $\|A\|\le1$.

\begin{definition}\label{def:LR}
The \emph{light-cone exponent at the site $x$} of the chain with phase $\theta$ is
\[
\Gamma_x(\theta)=\inf\Bigl\{a\ge0:\ \forall l,m:\ \sup_{|t|\le T}\ \sup_{A\in\mathcal A_l(x),\,B\in\mathcal B(x,T^a)}\|[\tau_t(A),B]\|=O(T^{-m})\Bigr\},
\]
and the \emph{uniform light-cone exponent} $\Gamma_{\mathrm{unif}}(\theta)$ is defined in the same way with an additional supremum over $x\in\Z$ inside the $O(T^{-m})$.
\end{definition}

Translation by $x$ sites maps the chain with phase $\theta$ to the chain with phase $\theta+x\alpha$, so $\Gamma_x(\theta)=\Gamma_0(\theta+x\alpha)$. By the ballistic bound $|\langle\delta_y,e^{-isH}\delta_x\rangle|\le(e|s|)^{|y-x|}/|y-x|!$, all these exponents are at most $1$.

\begin{theorem}[The fixed-site light cone]\label{thm:LRfixed}
Let $\lambda>1$, $\alpha$ irrational, and $h_n=2\lambda\cos2\pi(\theta+n\alpha)$. Let $\Ga^\Z=\bigcap_{x\in\Z}(\Ga-x\alpha)$, a set of full measure, and $\gamma=\max\{0,1-L/\beta(\alpha)\}$.
\begin{enumerate}
\item Let $\beta(\alpha)<\infty$ and $\theta\in\Ga^\Z$. For every $a>\gamma$, $l\ge0$ and $m>0$ there is $C$ with $\|[\tau_t(A),B]\|\le CT^{-m}$ for all $T\ge1$, $|t|\le T$, $A\in\mathcal A_l(0)$ and $B\in\mathcal B(0,T^a)$.
\item Let $\beta(\alpha)>L$ and $\theta$ arbitrary. For every $a<\gamma$ there are times $t_j\to\infty$ and finite sets $R_j\subset\{y:|y|>t_j^a\}$, each the union of two intervals, with
\[
\bigl\|[\tau_{t_j}(Z_0),P_{R_j}]\bigr\|\ge\sqrt{2c_B}=2^{-15.5}(1+4\pi\lambda)^{-1/2},\qquad c_B=\frac{2^{-32}}{1+4\pi\lambda}.
\]
\end{enumerate}
Consequently $\Gamma_x(\theta)=\max\{0,1-L/\beta(\alpha)\}$ for every site $x$ and every $\theta\in\Ga^\Z$, read as $1$ when $\beta(\alpha)=\infty$.
\end{theorem}

\begin{proof}
1. Fix $a'\in(\gamma,a)$. For $|z|\le l$ the phase $\theta+z\alpha$ lies in $\Ga$, and~\cite[Theorem~A]{DMamo} gives $P_{\theta+z\alpha}(S^{a'},S)=O(S^{-m'})$ for every $m'$. For large $T$, $T^a-l\ge(2T)^{a'}$. Lemma~\ref{lem:abel} at time scale $2T$ gives $\sup_{|s|\le2T}p^{(z)}_s(\{y:|y-z|>(2T)^{a'}\})=O(T^{-m''})$ for every $m''$. Lemma~\ref{lem:locality}, with finitely many $z$, finishes the proof; the constants depend on $\theta$ through the indices from which the sine conditions hold.

2. Choose $L<h<c<\beta$ with $1-h/c>a$, and scales $n_j$ with $b_{n_j}\ge c$ and $b_{n_j}\to\beta$. At $T_j=|\delta_{n_j}|^{-1}$ and $N_j=T_je^{-hq_{n_j}}=T_j^{1-h/b_{n_j}}$, \cite[Theorem~B]{DMamo} gives $P(N_j,T_j)\ge c_B$ at every phase. By~\eqref{eq:abelsplit} with $e^{-2k}=c_B/2$, some one-body time $s\le kT_j$ has $p^{(0)}_s(\{|y|>N_j\})\ge c_B/2$. Since $p_0=0$ and $s\mapsto p_s$ is continuous, there is $s_j\le s$ with $p^{(0)}_{s_j}(\{|y|>N_j\})=c_B/2$. Truncating to $R_j=\{N_j<|y|\le M_j\}$ with $M_j$ large keeps the mass in $[c_B/4,c_B/2]$, and then $4\sqrt{p(1-p)}\ge\sqrt{2c_B}$. Proposition~\ref{prop:comm} at spin time $t_j=s_j/2$ gives the bound. The times tend to infinity, because a fixed fraction of the packet cannot reach distance $N_j\to\infty$ in bounded time, by the ballistic bound. Finally $N_j\ge(kT_j)^a\ge t_j^a$ for large $j$, since $1-h/b_{n_j}>a$ eventually.

For the last statement, $\Gamma_x(\theta)=\Gamma_0(\theta+x\alpha)$ and $\theta+x\alpha\in\Ga^\Z$. Part~1 gives $\Gamma_0\le\gamma$ when $\beta<\infty$, and part~2 gives $\Gamma_0\ge\gamma$ when $\beta>L$; when $\beta\le L$, $\gamma=0$. When $\beta=\infty$, part~2 applies with any $a<1$.
\end{proof}

\begin{remark}\label{rem:LRprior}
Damanik, Lemm, Lukic and Yessen identified anomalous Lieb--Robinson exponents of the Fibonacci XY chain with upper one-body transport exponents and treated uniformity over the hull~\cite{DLLY2016}. Kachkovskiy proved velocity lower bounds for isotropic quasi-periodic XY chains in an absolutely continuous regime~\cite{Kachkovskiy2016}. Theorem~\ref{thm:LRfixed} is the corresponding statement for the Aubry--Andr\'e chain with Liouville frequencies, with its exact arithmetic exponent.
\end{remark}

The next lemma turns a polynomially small one-body tail into a commutator bound.

\begin{lemma}\label{lem:tailcomm}
Let $\varphi_n$ be phases, $T_n\to\infty$ and $N_n\ge T_n^{a_*-o(1)}$ with $a_*>0$, and suppose $P_{\varphi_n}(N_n,T_n)\ge T_n^{-C}$. Then for every $a<a_*$ there are $C'$, times $t_n\to\infty$ and finite sets $R_n\subset\{y:|y|>t_n^a\}$, each a union of two intervals, such that the chain with phase $\varphi_n$ satisfies $\|[\tau_{t_n}(Z_0),P_{R_n}]\|\ge t_n^{-C'}$.
\end{lemma}

The phases may vary with $n$, and a constant sequence is allowed; the proof treats each $n$ separately, and its only other input, the ballistic bound, holds uniformly in the phase because $\|v\|_\infty$ bounds the field at every phase.

\begin{proof}
By~\eqref{eq:abelsplit} with $e^{-2k}=T_n^{-C}/2$, some one-body time $s\le CT_n\log T_n$ has $p^{(0)}_s(\{|y|>N_n\})\ge T_n^{-C}/2$. Take the first time $s_n$ at which this mass equals $T_n^{-C}/4$, and truncate to a finite annulus $R_n=\{N_n<|y|\le M\}$ with mass $p\in[T_n^{-C}/8,T_n^{-C}/4]$. Proposition~\ref{prop:comm} gives $\|[\tau_{s_n/2}(Z_0),P_{R_n}]\|\ge4\sqrt{p/2}\ge T_n^{-C/2}$. The ballistic bound forces $s_n\ge cN_n$, so $T_n^{-C}$ is bounded below by a negative power of $t_n=s_n/2$; and $t_n\le CT_n\log T_n$, so $N_n\ge t_n^a$ for large $n$.
\end{proof}

\begin{theorem}[The uniform light cone]\label{thm:LRunif}
\begin{enumerate}
\item Let $v$ be real analytic with $H_{v,\alpha}$ supercritical, $\beta(\alpha)<\infty$, and $h_n=v(\theta+n\alpha)$. Then $\Gamma_{\mathrm{unif}}(\theta)\le\FU(\beta(\alpha))$ for every $\theta$.
\item Let $h_n=2\lambda\cos2\pi(\theta+n\alpha)$ with $\lambda>1$. Then $\Gamma_{\mathrm{unif}}(\theta)\ge\max\{0,1-L/(2\beta_e(\alpha))\}$ for every $\theta$, read as $1$ if $\beta_e=\infty$.
\item If moreover $\beta_e(\alpha)=\beta(\alpha)\in[L,\infty)$, then for every $\theta$
\[
\Gamma_{\mathrm{unif}}(\theta)=1-\frac L{2\beta(\alpha)},
\]
while $\Gamma_x(\theta)=1-L/\beta(\alpha)$ at every site $x$ when $\theta\in\Ga^\Z$.
\end{enumerate}
\end{theorem}

\begin{proof}
1. Let $a>a'>\FU(\beta)$. By Theorem~\ref{thm:U}, $\sup_\varphi P_\varphi((2T)^{a'},2T)=O(T^{-m})$ for every $m$. Lemma~\ref{lem:abel} turns this into $\sup_z\sup_{|s|\le2T}p^{(z)}_s(\{y:|y-z|>(2T)^{a'}\})=O(T^{-m'})$, uniformly in the site $z$, since its inequality is pointwise. For $A\in\mathcal A_l(x)$ and $B\in\mathcal B(x,T^a)$, the relevant sets lie in $\{y:|y-z|>T^a-l\}$ for $|z-x|\le l$, and Lemma~\ref{lem:locality} gives the bound uniformly in $x$.

2. Let $\beta_e>L/2$ and $h>L$ with $h/2+(h-L)<\beta_e$. For an even scale $n$ with $b_n\ge h/2+(h-L)$, the $q_{n+1}$ points $\theta+x\alpha$, $0\le x<q_{n+1}$, have gaps at most $\|q_n\alpha\|+\|q_{n+1}\alpha\|\le2/q_{n+1}$, and $1/q_{n+1}<2|\delta_n|$; so some $x_n$ with $0\le x_n<q_{n+1}<T_n=|\delta_n|^{-1}$ has $\|\theta+x_n\alpha-1/4\|\le2|\delta_n|$. Theorem~\ref{thm:quarter} at the phase $\theta+x_n\alpha$ gives $P_{\theta+x_n\alpha}(N_n,T_n)\ge T_n^{-2}e^{-2hq_n}$ with $N_n=\lfloor c_\eta T_ne^{-hq_n/2}\rfloor$. Lemma~\ref{lem:tailcomm}, applied to the chain at phase $\theta+x_n\alpha$, which is the chain at phase $\theta$ seen from the site $x_n$, gives commutators $\|[\tau_{t_n}(Z_{x_n}),P_{R_n}]\|\ge t_n^{-C'}$ with $R_n\subset\{y:|y-x_n|>t_n^a\}$ for every $a<1-h/(2\beta_e)$, along even scales with $b_n\to\beta_e$. So $\Gamma_{\mathrm{unif}}(\theta)\ge1-h/(2\beta_e)$; let $h\downarrow L$. If $\beta_e=\infty$ the same argument gives $1$.

3. Parts~1 and~2, with $\FU(\beta)=1-L/(2\beta)$ for the almost Mathieu operator when $\beta\ge L$ (Corollary~\ref{cor:typeI}), and Theorem~\ref{thm:LRfixed}.
\end{proof}

So uniformity over sites changes the light-cone exponent: when $\beta_e=\beta\in[L,\infty)$, almost every phase has fixed-site exponent $1-L/\beta$ at every site, and uniform exponent $1-L/(2\beta)$. At $\beta=\infty$ both equal $1$.

\subsection{Uncomputable light cones}

A real $x$ is \emph{right-c.e.\ relative to $\mathbf 0'$} if $\{r\in\Q:r>x\}$ is computably enumerable relative to the halting problem. Let $K_2$ be the halting set relative to $\mathbf 0'$ and $s_*=2-\sum_{e\in K_2}4^{-(e+2)}\in[23/12,2]$, a real of Turing degree $\mathbf 0''$~\cite[proof of Theorem~D]{DMamo}.

\begin{corollary}[A computable chain with a light cone of degree $\mathbf 0''$]\label{cor:LRdeg}
Let $\alpha_*$ be the explicit computable irrational of~\cite[Theorem~D(3)]{DMamo}, whose partial quotients after the first are even, so that all its denominators are odd, and $\beta(\alpha_*)=s_*\log2$. Consider the chain with $h_n=4\cos2\pi(1/4+n\alpha_*)=-4\sin(2\pi n\alpha_*)$. At every site $x$, the light-cone exponent $\Gamma_x(1/4)=1-1/s_*\in[11/23,1/2]$ has Turing degree $\mathbf 0''$. As $\alpha$ ranges over the computable irrationals with eventually odd denominators, the exponents $\Gamma_0(1/4)$ of the chains $h_n=4\cos2\pi(1/4+n\alpha)$ range over exactly the reals in $[0,1]$ that are right-c.e.\ relative to $\mathbf 0'$. For every computable $\lambda'$ with $1<\lambda'<2^{23/12}$, the chain $h_n=2\lambda'\cos2\pi(1/4+n\alpha_*)$ has $\Gamma_x(1/4)=1-\log\lambda'/\beta(\alpha_*)$ at every site, again of degree $\mathbf 0''$.
\end{corollary}

\begin{proof}
For odd $q_n$ and every fixed $x$, $|\sin(2\pi q_n(1/4+x\alpha))|=|\cos(2\pi x\delta_n)|\ge1/2$ for large $n$. So $1/4\in\Ga^\Z$ whenever the denominators are eventually odd, and Theorem~\ref{thm:LRfixed} identifies $\Gamma_x(1/4)$ with the one-body exponent of~\cite[Theorem~D]{DMamo}, including the value $1$ when $\beta=\infty$. The degree and value-set statements are~\cite[Theorem~D(2),(3)]{DMamo}. For $\lambda'$: since $\beta(\alpha_*)\ge(23/12)\log2>\log\lambda'$, the exponent $x'=1-\log\lambda'/\beta$ is positive, and $\beta=\log\lambda'/(1-x')$; both maps are computable for computable $\lambda'$, so $x'$ and $\beta$ have the same degree.
\end{proof}

The value-set statement concerns this family, with $\lambda=2$, phase $1/4$ and eventually odd denominators; it is not a classification over all computable fields and phases.

\begin{corollary}[Two exponents of degree $\mathbf 0''$ in one computable chain]\label{cor:twoexp}
Let $(v_m)$ be a computable sequence of rationals in $[1,2]$ with $\limsup v_m=s_*$~\cite[Lemma~10.1]{DMamo}. Define $\alpha=[0;a_1,a_2,\dots]$ by $a_1=2$ and, with $n_1=1$ and $n_{m+1}=n_m+3$,
\[
a_{n_m+1}=2^{\lceil v_mq_{n_m}\rceil},\qquad a_{n_m+2}=a_{n_m+3}=1 .
\]
Then $\alpha$ is computable, $\beta_e(\alpha)=\beta(\alpha)=s_*\log2$ and $\beta_o(\alpha)=0$. For the chains $h_n=4\cos2\pi(\theta+n\alpha)$:
\begin{enumerate}
\item for every $\theta$, $\Gamma_{\mathrm{unif}}(\theta)=1-1/(2s_*)\in[17/23,3/4]$;
\item for every $\theta\in\Ga^\Z$, in particular for almost every $\theta$, $\Gamma_x(\theta)=1-1/s_*\in[11/23,1/2]$ at every site $x$;
\item at $\theta=1/4$, $\Gamma_x(1/4)=1-1/(2s_*)$ at every site $x$;
\item there is a computable phase $\theta_*\in\Ga^\Z$.
\end{enumerate}
Both exponent values, $1-1/(2s_*)$ and $1-1/s_*$, have Turing degree $\mathbf 0''$. In particular the computable chain with phase $\theta_*$ has uniform exponent $1-1/(2s_*)$ and fixed-site exponent $1-1/s_*$ at every site: two different exponents of degree $\mathbf 0''$.
\end{corollary}

\begin{proof}
\emph{Arithmetic.} $q_{n_1}=q_1=2$ is even. If $q_{n_m}$ is even, then $q_{n_m-1}$ is odd, $q_{n_m+1}$ and $q_{n_m+2}$ are odd, and $q_{n_m+3}=q_{n_m+2}+q_{n_m+1}$ is even. The large partial quotient gives $b_{n_m}=v_m\log2+O(\log q_{n_m}/q_{n_m})$, while $b_n\to0$ along all other $n$. So $\beta=\beta_e=s_*\log2\ge\log2=L$ and $\beta_o=0$. The partial quotients are computable, so $\alpha$ is computable, although $\limsup v_m$ is not.

Part~1 is Theorem~\ref{thm:LRunif}(3), and part~2 is Theorem~\ref{thm:LRfixed}.

Part~3. Fix $x$. The first block to the right of $x$ in the chain with phase $1/4$ has phases $1/4+(x+1+k)\alpha$, and the rational block with phases $1/4+(x+1+k)p_n/q_n$ is a turning block, since $\cos(2\pi q_n(1/4+(x+1)p_n/q_n))=(-1)^{q_n/2}$ for even $q_n$. The site discrepancies are $(x+1+k)\delta_n/q_n$, at most $2|\delta_n|$ once $q_n>|x|$, so the proof of Theorem~\ref{thm:quarter} applies with the same exponents and gives $\au(1/4+x\alpha)\ge1-L/(2\beta_e)$. Lemma~\ref{lem:tailcomm} gives $\Gamma_x(1/4)\ge1-1/(2s_*)$, and part~1 gives the reverse inequality.

Part~4. Let $E_n=\{\theta:|\sin2\pi q_n\theta|<q_n^{-1/2}\}$, a finite union of intervals with computable endpoints and $|E_n|\le q_n^{-1/2}$. Since $q_n$ grows at least like the Fibonacci numbers, there is an explicit $N$ with $\sum_{n\ge N}|E_n|<1/2$, with an effective tail bound. Hence the measure of $(\bigcup_{n\ge N}E_n)\cap I$ is computable for every dyadic interval $I$, by finite unions and the tail bound. Start from $I_0=\T$; given $I_k$ whose intersection with the closed set $F=\T\setminus\bigcup_{n\ge N}E_n$ has positive measure, compute the measures for the two dyadic children until one is certified positive, which must happen, and let $I_{k+1}$ be that child. The nested intervals define a computable $\theta_*$. If $\theta_*$ were in the open set $\bigcup_{n\ge N}E_n$, a neighbourhood of $\theta_*$, hence some $I_k$, would lie in it, contradicting $|I_k\cap F|>0$. So $|\sin2\pi q_n\theta_*|\ge q_n^{-1/2}$ for $n\ge N$. For each fixed $x$, $q_n(\theta_*+x\alpha)=q_n\theta_*+xp_n+x\delta_n$, and the sine changes by at most $2\pi|x||\delta_n|\le2\pi|x|/q_n$, which is eventually below $q_n^{-1/2}/2$. So $\theta_*+x\alpha\in\Ga$ for every $x$.

\emph{Degrees.} $1-1/s_*$ and $1-1/(2s_*)$ are computable functions of $s_*$ on $[23/12,2]$ with computable inverses, so both have the degree of $s_*$, which is $\mathbf 0''$.
\end{proof}

Counting lattice sites is not averaging over phases: at the phase $1/4$ every site sees the exceptional exponent, while at almost every phase every site sees the typical one.

\section{Computability}\label{sec:computability}

The companion paper A showed that the transport exponent of one explicit computable almost Mathieu operator has Turing degree $\mathbf 0''$, and that this survives changes of the coupling in an interval~\cite[Theorem~D]{DMamo}. Here we show that, at the frequency $\alpha_*$ and the phase $1/4$, the degree survives every computable even change of the potential in an explicit strip ball (Theorem~\ref{thm:robust}), and we locate the source of the uncomputability: for a computable analytic potential the Lyapunov exponent is computable at every computable irrational frequency, so in the transport laws identified here the uncomputability sits in the arithmetic of the frequency.

\subsection{Conventions}

A real number is \emph{computable from an oracle} if the oracle computes rationals $r_k$ with $|r_k-x|\le2^{-k}$. A name of a real $x$ is any such sequence. A function on $\R$, or on $(\R\setminus\Q)\times\R$, is \emph{computable} if an algorithm turns every name of the argument into a name of the value; equivalently, if the sets $\{f>r\}$ and $\{f<r\}$, for rational $r$, are enumerated uniformly by rational boxes, that is, if $f$ is both lower and upper semicomputable. A real-analytic $v$ is \emph{computable} if its Fourier coefficients $\hat v(k)$, with $\hat v(-k)=\overline{\hat v(k)}$, form a computable sequence and there are computable rationals $C_v,\delta'_v>0$ with $|\hat v(k)|\le C_ve^{-2\pi\delta'_v|k|}$. Then $v$ is uniformly computable on every closed strip $|\operatorname{Im}z|\le y<\delta'_v$. Trigonometric polynomials with computable coefficients are the main examples. Two reals have the same Turing degree if each is computable from the other.

We write $A_{E,y}(x)=T(E,x+iy)$ for the complexified cocycle and $L_n(E,y)=\frac1n\int_\T\log\|A_n(E,x+iy)\|\,dx$, so that $L(E,y)=\inf_nL_n(E,y)$ by subadditivity. We use the following facts from Avila's global theory~\cite{Avila2015}, with the numbering of the published version: convexity of $y\mapsto L(E,y)$ (Section~2.1); quantization (Theorem~5); regularity, which for real cocycles holds if and only if the acceleration is zero (Definition~3 and Remark~4); regular with positive Lyapunov exponent if and only if uniformly hyperbolic (Theorem~6); and Corollary~7: for every analytic $SL(2,\C)$ cocycle at an irrational frequency there is $\eps_0>0$ such that either $L(A_\eps)=0$ for all $0<\eps<\eps_0$, or $A_\eps$ is uniformly hyperbolic for all $0<\eps<\eps_0$, and in the latter case $\eps\mapsto L(A_\eps)$ is affine on $(0,\eps_0)$. Uniformly hyperbolic cocycles form an open set, have continuous invariant line fields, and their Lyapunov exponent is continuous~\cite[Section~2.1]{Avila2015}.

\subsection{The Lyapunov exponent is computable}

\begin{lemma}[Chord bound]\label{lem:chord}
For real $E$ and $0<y<y'<\delta_v$,
\[
\ell(E)\ge\frac{y'L(E,y)-yL(E,y')}{y'-y},
\]
with equality if and only if $L(E,\cdot)$ is affine on $[0,y']$.
\end{lemma}

\begin{proof}
Write $y=(1-\sigma)\cdot0+\sigma y'$ with $\sigma=y/y'$. Convexity gives $L(E,y)\le(1-\sigma)\ell(E)+\sigma L(E,y')$, which rearranges to the claim; equality at an interior point of a chord forces affineness on the whole chord.
\end{proof}

\begin{lemma}[The minimum lies on the spectrum]\label{lem:minonsigma}
For every irrational $\alpha$ and real-analytic $v$, $\inf_{E\in\R}\ell(E)=\min_\Sigma\ell$. In particular $\ell_*=\min_{|E|\le K}\ell(E)$ for any $K\ge\|v\|_\infty+2$.
\end{lemma}

\begin{proof}
$\ell$ is continuous~\cite{BourgainJitomirskaya2002}. On a bounded gap $(a,b)$ of $\Sigma$, the Thouless formula $\ell(E)=\int\log|E-E'|\,dk(E')$~\cite{CraigSimon1983} can be differentiated twice, since $\supp dk=\Sigma$ is at positive distance from compact subintervals of the gap, and $\ell''(E)=-\int(E-E')^{-2}\,dk(E')<0$. So $\ell$ is concave on the gap, and its minimum over $[a,b]$ is attained at an endpoint. Outside the convex hull of $\Sigma$, $\ell'(E)=\int(E-E')^{-1}\,dk(E')$ is negative for $E<\min\Sigma$ and positive for $E>\max\Sigma$, so $\ell$ increases away from $\Sigma$.
\end{proof}

\begin{lemma}[Certificates]\label{lem:cert}
Let $B:\T\to SL(2,\C)$ be continuous and $\alpha$ irrational. A \emph{certificate} is a pair $(m,Q)$ with $m\ge1$ and $Q$ a trigonometric-polynomial matrix with Gaussian-rational coefficients such that
\begin{itemize}
\item[(V1)] $\det Q(x)\ne0$ for all $x\in\T$;
\item[(V2)] the entries of $\mathcal C(x)=Q(x+m\alpha)^{-1}B^{(m)}(x)Q(x)$ satisfy $|c_{11}|-|c_{12}|>|c_{21}|+|c_{22}|$ for all $x\in\T$,
\end{itemize}
where $B^{(m)}(x)=B(x+(m-1)\alpha)\cdots B(x)$. Its value is $\mathrm{val}(m,Q)=\frac1m\int_\T\log(|c_{11}|-|c_{12}|)\,dx$.
\begin{enumerate}
\item (Soundness) $L(\alpha,B)\ge\mathrm{val}(m,Q)$.
\item (Effectiveness) From names of $B$ and $\alpha$, the strict conditions (V1)--(V2) are semi-decidable and $\mathrm{val}(m,Q)$ is computable. A certificate verified for $(B,\alpha)$ remains one, with value changed by less than any prescribed $\eta>0$, for all $(B',\alpha')$ in a neighbourhood.
\item (Completeness) If $(\alpha,B)$ is uniformly hyperbolic, then for every $\eta>0$ there is a certificate with $\mathrm{val}>L(\alpha,B)-\eta$.
\end{enumerate}
\end{lemma}

\begin{proof}
1. Let $\phi=|c_{11}|-|c_{12}|>0$. By (V2) the cone $\{(u,w)\in\C^2:|w|\le|u|\}$ is mapped into itself by $\mathcal C(x)$, and the first coordinate grows at least by the factor $\phi(x)$. Iterating along $x,x+m\alpha,\dots$ gives $\|\mathcal C^{(k)}(x)\|\ge\prod_{j<k}\phi(x+jm\alpha)$, so $\log\|B^{(km)}(x)\|\ge\sum_{j<k}\log\phi(x+jm\alpha)-\log(\|Q\|_\infty\|Q^{-1}\|_\infty)$. Integrate over $\T$, divide by $km$ and let $k\to\infty$.

2. All quantities are continuous functions of $x$ on the compact circle, computable from the names, and a strict inequality between such functions is certified by finitely many evaluations with error bounds. Once (V1) and (V2) are certified, $\phi$ is effectively bounded away from zero, so its logarithm and its integral are computable. The strict margins persist under small perturbations.

3. A uniformly hyperbolic cocycle has continuous invariant line fields $u(x)$, $s(x)$, transverse to each other. Every complex line bundle over the circle is trivial, so they are spanned by continuous nonvanishing periodic vector fields, and conjugation by $P=[u,s]$ makes the cocycle diagonal, $\operatorname{diag}(\lambda_u,\lambda_s)$, with $\int\log|\lambda_u|=L(\alpha,B)$ and $\int\log|\lambda_s|=-L(\alpha,B)$. By unique ergodicity, $\frac1m\log|\Lambda_m|\to L$ and $\frac1m\log|\Sigma_m|\to-L$ uniformly, where $\Lambda_m,\Sigma_m$ are the $m$-step products of $\lambda_u,\lambda_s$. Approximate $P$ by a trigonometric polynomial matrix $Q$ with rational coefficients within $\eps'$ (Fej\'er). Then $\mathcal C=(I+Y)\operatorname{diag}(\Lambda_m,\Sigma'_m)(I+X)$ with $\|X\|,\|Y\|\le C\eps'$ and $|\Sigma'_m|\le C|\Sigma_m|$, and for $C\eps'<1/8$ and large $m$ the conditions (V1)--(V2) hold with $\phi\ge|\Lambda_m|(1-C\eps')$. So $\mathrm{val}\ge L+\frac1m\log(1-C\eps')$.
\end{proof}

Adding the condition $\min_\T\phi>\mu$ for a rational $\mu>1$ gives a certificate of uniform hyperbolicity: then $\|B^{(km)}(x)\|\ge c\mu^k$ uniformly. Completeness holds for these certificates as well, since $|\Lambda_m|\to\infty$ uniformly.

\begin{theorem}[The Lyapunov exponent is computable]\label{thm:LEcomp}
Let $v$ be a computable real-analytic function. The map $(\alpha,E)\mapsto\ell(E)$, the Lyapunov exponent of $H_{v,\alpha}$ at $E$, is computable on $(\R\setminus\Q)\times\R$, uniformly in a name of $v$. Consequently, for every computable irrational $\alpha$, $E\mapsto\ell(E)$ is a computable function, $\ell(E)$ is a computable real at every computable $E$, and $\ell_*$ is a computable real, uniformly in $v$ and $\alpha$. No hypothesis of supercriticality or type is needed.
\end{theorem}

\begin{proof}
\emph{Upper semicomputability.} $\ell(E)<r$ if and only if there are $n$ and a rational interval $I\ni E$ with $\max_{E'\in\bar I}L_n(E',0)<r$, because $\ell=\inf_nL_n$ and each $L_n$ is continuous in $E$ and in $\alpha$. The maxima are computable.

\emph{Lower semicomputability.} Let $W$ be the computably enumerable set of pairs $(I,r)$, with $I$ a rational interval and $r$ rational, such that either $r<0$, or there are rationals $0<y<y'<\delta'_v$, a certificate $(m,Q)$ for the cocycles $A_{E,y}$ valid uniformly for $E\in\bar I$, with $c_{\mathrm{low}}=\min_{E\in\bar I}\mathrm{val}_E(m,Q)$, and an $n$ with $c_{\mathrm{up}}=\max_{E\in\bar I}L_n(E,y')+2^{-n}$, such that $(y'c_{\mathrm{low}}-yc_{\mathrm{up}})/(y'-y)>r$.

\emph{Soundness}: if $(I,r)\in W$ and $E\in I$, then $\ell(E)>r$, by Lemma~\ref{lem:chord} with $L(E,y)\ge c_{\mathrm{low}}$ (Lemma~\ref{lem:cert}(1)) and $L(E,y')\le L_n(E,y')<c_{\mathrm{up}}$.

\emph{Completeness}: let $r<\ell(E_0)$ with $r\ge0$. Then $\ell(E_0)>0$, so $L(E_0,y)\ge\ell(E_0)>0$ for all $y$, by convexity and evenness. Corollary~7 of~\cite{Avila2015} gives $\eps_0>0$ such that $A_{E_0,y}$ is uniformly hyperbolic and $y\mapsto L(E_0,y)$ is affine for $0<y<\eps_0$; by continuity of the convex function at $0$, it is affine on $[0,\eps_0)$. Take rationals $y<y'=2y<\min\{\eps_0,\delta'_v\}$. Then Lemma~\ref{lem:chord} holds with equality at $E_0$: $2L(E_0,y)-L(E_0,2y)=\ell(E_0)$. Given $\eta>0$, Lemma~\ref{lem:cert}(3) and~(2) give a certificate at $(E_0,y)$ with value above $L(E_0,y)-2\eta$ on a neighbourhood of $E_0$. Choose $n$ with $L_n(E_0,2y)+2^{-n}<L(E_0,2y)+\eta$ and shrink the neighbourhood so that $L_n(\cdot,2y)+2^{-n}<L(E_0,2y)+2\eta$ on it. On a rational interval $I\ni E_0$ in that neighbourhood, $2c_{\mathrm{low}}-c_{\mathrm{up}}\ge\ell(E_0)-6\eta>r$ for small $\eta$. So $(I,r)\in W$.

Hence $\ell(E)>r$ if and only if some $(I,r)\in W$ has $E\in I$, which is semi-decidable from names of $E$ and $\alpha$; all verifications are continuous in $\alpha$, with margins. Irrationality of $\alpha$ is used only for completeness. The two halves give a computable function. Finally $\ell_*=\min_{|E|\le K}\ell(E)$ by Lemma~\ref{lem:minonsigma}, with a computable integer $K$ larger than a computable bound for $\|v\|_\infty+2$, and the minimum of a computable function over a computable compact interval is computable.
\end{proof}

The algorithm is: where the Lyapunov exponent is positive, quantization makes the cocycle regular, hence uniformly hyperbolic, at small positive heights; there a certificate gives a lower bound; convexity carries it to height zero exactly. Where it is zero, nonnegativity and finite-scale upper bounds suffice. No modulus of continuity and no large-deviation estimate is used, so Liouville frequencies are allowed. The proof gives no bound on the running time, since the height at which the Lyapunov exponent becomes affine is found by search.

\begin{corollary}[Density of states and spectrum]\label{cor:IDScomp}
For computable $v$ and computable irrational $\alpha$: the integrated density of states $k$ is a computable function; $\Sigma$ is a computable compact set; $\max_\Sigma\ell$ and $\min\Sigma$ are computable reals, and $\min\Sigma$ is a computable point of $\Sigma$.
\end{corollary}

\begin{proof}
Let $c_N(\theta,E)$ count the eigenvalues at most $E$ of $H_\theta$ restricted to $\{0,\dots,N-1\}$. Decoupling two adjacent sites is a rank-two perturbation with one positive and one negative eigenvalue, so $|c_{N+M}(\theta)-c_N(\theta)-c_M(\theta+N\alpha)|\le1$. With $a_N(E)=\int_\T c_N(\theta,E)\,d\theta$, this gives $|a_{N+M}-a_N-a_M|\le1$ and hence $|k(E)-a_N(E)/N|\le1/N$. The function $a_N$ is upper semicomputable, as a sum of measures of effectively closed sets cut out by computable eigenvalue functions, and the analogous count with strict inequalities is lower semicomputable. The density of states has no atoms: an atom at $E$ would make $\ell(E)=\int\log|E-E'|\,dk(E')=-\infty$. So both bounds converge to $k(E)$, and $k$ is computable. The complement of $\Sigma$ is enumerated by the certificates of uniform hyperbolicity, uniformly on rational intervals; and an open interval $(a,b)$ meets $\Sigma$ if and only if $k(b)>k(a)$, which is semi-decidable. Together with a bounding interval this makes $\Sigma$ a computable compact set, and the extrema follow.
\end{proof}

\begin{remark}[Related work]\label{rem:compprior}
Computing the spectrum is not new: the rational approximation of~\cite[Lemma~2.7]{DMuniv}, with finite Floquet matrices, already gives an effective route, and spectral computations for broad operator classes are studied through the solvability complexity index~\cite{BACHNS2015,ColbrookHansen2019}; Colbrook, Embree and Fillman study the computation of spectral size, including for the almost Mathieu operator~\cite{CEF2026}. Theorem~\ref{thm:LEcomp} is an algorithmic consequence of Avila's global theory. We did not find this formulation, uniform over all irrational frequencies, in the literature we checked, but we make no priority claim for it.
\end{remark}

\subsection{An explicit type I neighbourhood}

Let $v_0(x)=4\cos2\pi x$, the almost Mathieu potential with $\lambda=2$, and $S_Y=\{z:|\operatorname{Im}z|\le Y\}$.

\begin{theorem}[An explicit neighbourhood]\label{thm:nbhd}
Let $Y,\eps>0$, $t=e^{-2\pi Y}$ and $E_{\max}=6+\eps$, and suppose
\[
M:=\frac2t-2t-E_{\max}-\eps>2,\qquad p_0:=\frac{E_{\max}t+2t^2}2<1,\qquad\kappa:=\frac{t(\eps+\rho)}{2(1-p_0)}<\frac12,
\]
where $\rho=(M-\sqrt{M^2-4})/2$. Let $v$ be real analytic with $w=v-v_0$ holomorphic on a neighbourhood of $S_Y$ and $\sup_{S_Y}|w|\le\eps$. Then for every irrational $\alpha$:
\begin{enumerate}
\item $\ell(E)\ge\log2+\log(1-\kappa)>0$ for every real $E$ with $|E|\le E_{\max}$;
\item $H_{v,\alpha}$ is supercritical and type~I, with $L(E,y)=\ell(E)+2\pi|y|$ for $E\in\Sigma$ and $|y|\le Y$;
\item $\log2+\log(1-\kappa)\le\ell(E)\le\log2+\log(1+\kappa)$ for $E\in\Sigma$.
\end{enumerate}
For $Y=\eps=1/4$ one has $M>2.7051$, $\rho<0.4419$, $p_0<0.6929$ and $0.2341<\kappa<0.23411$, so $\ell\in[0.4264,0.9035]$ on $\Sigma$.
\end{theorem}

\begin{proof}
Fix a real $E$ with $|E|\le E_{\max}$. The three strict inequalities persist, with constants $M(y),\rho(y),p_0(y),\kappa(y)$ depending continuously on the height, for $y$ in an interval $[Y-\tau,Y]$. Write $z=x+iy$ and $\mathsf q=e^{2\pi iz}$, so $|\mathsf q|=e^{-2\pi y}$. Then $a(z):=E-v(z)=-2\mathsf q^{-1}+E-2\mathsf q-w(z)$ and $|a(z)|\ge M(y)>2$ on the line $\operatorname{Im}z=y$.

\emph{An invariant section.} The maps $\sigma\mapsto1/(a(\zeta)-\sigma)$ send the disk $|\sigma|\le\rho(y)$ into itself, since $|a-\sigma|\ge M(y)-\rho(y)=1/\rho(y)$, and contract it. Choosing one common disk on a slightly smaller substrip, using its lower endpoint, the limit $r(z)=\lim_n f_{z-\alpha}\circ\dots\circ f_{z-n\alpha}(0)$ exists uniformly, is holomorphic in $z$, satisfies $r(z+\alpha)=1/(a(z)-r(z))$, and on each line $\operatorname{Im}z=y$ satisfies $|r|\le\rho(y)$. With $u=(1,r)^{\mathsf T}$ and $g=a-r$, $T(E,z)u(z)=g(z)u(z+\alpha)$, and conjugation by $P=[u,e_2]$, of determinant one, gives the triangular cocycle $\bigl(\begin{smallmatrix}g&-1\\0&1/g\end{smallmatrix}\bigr)$ with $|g|\ge1/\rho>1$. Its $n$-step products are comparable to their diagonal, so $L(E,y)=\int_0^1\log|g(x+iy)|\,dx$.

\emph{Slope $2\pi$.} Write $g=-2\mathsf q^{-1}(1-\Phi)$ with $\Phi=\mathsf q(E-2\mathsf q-w-r)/2=\Phi_0+\Phi_1$, where $\Phi_0=\mathsf q(E-2\mathsf q)/2$ satisfies $|\Phi_0|\le p_0(y)$ on the disk $|\mathsf q|\le e^{-2\pi y}$ by the maximum principle, and $|\Phi_1|\le e^{-2\pi y}(\eps+\rho(y))/2$. So $|\Phi|<1$, the winding number of $1-\Phi$ along a horizontal line is zero, and that of $g$ is $-1$. Since $\log|g|$ is harmonic, $y\mapsto\int_0^1\log|g(x+iy)|\,dx$ is affine, with slope $-2\pi$ times the winding number, that is $2\pi$.

\emph{The value at height $Y$.} The mean of $\log|1-\Phi_0|$ over the circle $|\mathsf q|=t$ is $\log|1-\Phi_0(0)|=0$, and $|\Phi_1/(1-\Phi_0)|\le\kappa$. So $L(E,Y)=\log2+2\pi Y+c$ with $\log(1-\kappa)\le c\le\log(1+\kappa)$.

\emph{Conclusions.} By convexity every slope of $L(E,\cdot)$ on $[0,Y]$ is at most $2\pi$, so $\ell(E)\ge L(E,Y)-2\pi Y\ge\log2+\log(1-\kappa)>0$. The spectrum lies in $[-2-\sup_\R|v|,2+\sup_\R|v|]\subset[-E_{\max},E_{\max}]$, so $\ell>0$ on $\Sigma$. At $E\in\Sigma$ the acceleration is an integer at least one (Section~\ref{sec:prelim}) and at most one, so the slope is $2\pi$ on all of $(0,Y]$, which gives part~2 and $\ell(E)=L(E,Y)-2\pi Y$, which gives part~3. The numerical values are elementary; they were also certified independently with exact rational interval arithmetic, which gives $\kappa\in[0.2341020899,0.2341020900]$, $\log2+\log(1-\kappa)>0.42644$ and $\log2+\log(1+\kappa)<0.90350$.
\end{proof}

The estimates are pointwise in $x$, so they hold for every irrational frequency at once, for even and non-even $v$.

\begin{proposition}[Openness]\label{prop:open}
Fix an irrational $\alpha$, and let $v_0$ be any real-analytic potential with $H_{v_0,\alpha}$ supercritical and type~I, with common strip $h$~\cite[Lemma~5.2]{DMuniv}. There is $\eps_0>0$ such that every real-analytic $v$ with $\|v-v_0\|_h:=\sup_{|\operatorname{Im}z|\le h}|v-v_0|<\eps_0$ is supercritical and type~I at $\alpha$, and $\ell_*(v)\to\ell_*(v_0)$ as $\|v-v_0\|_h\to0$.
\end{proposition}

\begin{proof}
For $E\in\Sigma_0=\Sigma(v_0)$ and $y\in[h/4,3h/4]$, $L_{v_0}(E,\cdot)$ is affine near $y$ with positive value, so the cocycle is uniformly hyperbolic. Uniform hyperbolicity is open, and $(E,y,v)\mapsto T_v(E,\cdot+iy)$ is continuous into $C^0$ for $v$ in the norm $\|\cdot\|_h$. A finite cover of the compact set $K\times[h/4,3h/4]$, with $K=\{E:\dist(E,\Sigma_0)\le e_0\}$, gives $e_0,\eps_1>0$ such that these cocycles are uniformly hyperbolic for $E\in K$, $y\in[h/4,3h/4]$ and $\|v-v_0\|_h\le\eps_1$. There $y\mapsto L_v(E,y)$ is affine on $[h/4,3h/4]$, with a slope in $2\pi\Z$ that depends continuously on $(E,v)$, hence is constant on each connected component of $K\times\{\|v-v_0\|_h\le\eps_1\}$; each component meets $\Sigma_0\times\{v_0\}$, where the slope is $2\pi$. Let $\Psi(E,v)=L_v(E,h/4)-\pi h/2$. At $v_0$ and $E\in\Sigma_0$ it equals $\ell_{v_0}(E)$; shrinking $e_0$ and $\eps_1$, $\Psi\ge\ell_*(v_0)/2$ on the region, and, by joint continuity near $K\times\{v_0\}$ and a finite cover of $K$, $\sup_{E\in K}|\Psi(E,v)-\Psi(E,v_0)|\to0$ as $v\to v_0$. For $E\in\Sigma_v$, which lies within $\|v-v_0\|_h$ of $\Sigma_0$, convexity gives $\ell_v(E)\ge\Psi(E,v)>0$, and the argument of Theorem~\ref{thm:nbhd}(2) gives $\om_v(E)=1$ and $\ell_v=\Psi(\cdot,v)$ on $\Sigma_v$. Since $\Psi(\cdot,v_0)$ is uniformly continuous on $K$ and $\Sigma_v\to\Sigma_0$ in Hausdorff distance, $\ell_*(v)=\min_{\Sigma_v}\Psi(\cdot,v)\to\min_{\Sigma_0}\Psi(\cdot,v_0)=\ell_*(v_0)$.
\end{proof}

The topology is that of a fixed strip. This is not openness in the operator norm among all bounded potentials.

\subsection{A robust exponent of degree \texorpdfstring{$\mathbf 0''$}{0''}}

Let
\[
\mathcal U=\Bigl\{v\text{ real analytic}:\ v-v_0\text{ holomorphic near }S_{1/4},\ \sup_{S_{1/4}}|v-v_0|<\frac14\Bigr\},
\]
an open ball in the norm of the strip $S_{1/4}$, and its even part $\mathcal U_{\mathrm{even}}=\{v\in\mathcal U:v(-x)=v(x)\}$. Let $\alpha_*$ be the explicit computable frequency of~\cite[Theorem~D(3)]{DMamo}: all its denominators are odd, $a_{n+1}\ge2^{q_n}$, and $\beta(\alpha_*)=s_*\log2$ with $s_*\in[23/12,2]$ of Turing degree $\mathbf 0''$. So $\beta(\alpha_*)\ge(23/12)\log2>1.3285$. For $v\in\mathcal U$ put
\[
X(v)=\au(H_{v,\alpha_*,1/4}),\qquad X_{\mathrm{typ}}(v)=1-\frac{\ell_*(v,\alpha_*)}{\beta(\alpha_*)} .
\]

\begin{theorem}[A robust uncomputable exponent]\label{thm:robust}
\begin{enumerate}
\item For every $v\in\mathcal U_{\mathrm{even}}$, $X(v)=X_{\mathrm{typ}}(v)\in[0.3199,0.6924]$, and this is also the value of every moment exponent $\beta^+(p)$, of the fastest exponent for the packet started at any site, and of the upper packing dimension of the spectral measure of $\delta_0$.
\item For every $v\in\mathcal U_{\mathrm{even}}$ with computable coefficients, $X(v)$ has Turing degree exactly $\mathbf 0''$.
\item On $\mathcal U_{\mathrm{even}}$, $v\mapsto X(v)$ is continuous in the norm of $S_{1/4}$ and not constant: along $v=2\lambda\cos2\pi x$ with $|\lambda-2|<0.0498$, $X=1-\log\lambda/\beta(\alpha_*)$. On all of $\mathcal U$, $X_{\mathrm{typ}}$ is continuous, and it is the exponent at almost every phase.
\end{enumerate}
\end{theorem}

\begin{proof}
By Theorem~\ref{thm:nbhd} with $Y=\eps=1/4$, every $v\in\mathcal U$ is supercritical and type~I at every irrational frequency, with $0.4264\le\ell_*\le0.9035<1.3285\le\beta(\alpha_*)<\infty$. Every denominator of $\alpha_*$ is odd, so $|\sin(2\pi q_n/4)|=1$ and $1/4\in G_{\alpha_*}$, which is contained in $G^v_{\alpha_*}$ for even $v$~\cite[Lemma~5.22]{DMuniv}. The transport law~\cite[Theorem~A]{DMuniv} gives $X(v)=1-\ell_*/\beta$, without clipping; \cite[Theorem~B(1)]{DMuniv} gives the same value for every $\beta^+(p)$ and shows that $\ell$ attains $\ell_*$ on the support of the spectral measure, and~\cite[Theorem~B(3)]{DMuniv} then gives the packing dimension. The packet started at $\delta_j$ is a translate of the packet started at $\delta_0$ for the phase $1/4+j\alpha_*$, and $|\sin(2\pi q_n(1/4+j\alpha_*))|=|\cos(2\pi j\delta_n)|\ge1/2$ for large $n$. This proves part~1.

For part~2, $\ell_*$ is computable by Theorem~\ref{thm:LEcomp}, and positive. So $X=1-\ell_*/\beta$ is computable from $\beta$, and $\beta=\ell_*/(1-X)$ is computable from $X$. Hence $X$ has the degree of $\beta(\alpha_*)$, which is that of $s_*$, namely $\mathbf 0''$.

Part~3 follows from Proposition~\ref{prop:open} at the fixed frequency $\alpha_*$, and from $|2\lambda\cos2\pi z-4\cos2\pi z|\le2|\lambda-2|\cosh(\pi/2)$ on $S_{1/4}$. The almost-everywhere statement is~\cite[Theorem~A]{DMuniv}.
\end{proof}

\begin{remark}[Computable coefficients are needed]\label{rem:noncomp}
Part~2 fails without them. Let $\lambda=\exp(\frac{51}{100}\beta(\alpha_*))$, so $1.969<\lambda<2.028$, and $v=2\lambda\cos2\pi x$. Then $2|\lambda-2|\cosh(\pi/2)<1/4$, so $v\in\mathcal U_{\mathrm{even}}$, and $X(v)=1-\log\lambda/\beta(\alpha_*)=49/100$ is computable; the coupling $\lambda$ is not. The fixed phase $1/4$ is also needed in part~3: for non-even $v\in\mathcal U$ we do not identify the exponent at the phase $1/4$, and use instead the phase of Proposition~\ref{prop:goodphase}. Not every direction changes the value: adding a constant to $v$ shifts the energy and leaves $X$ unchanged.
\end{remark}

\begin{theorem}[Energy-resolved exponents]\label{thm:energyres}
Let $v\in\mathcal U_{\mathrm{even}}$ have computable coefficients, and let $E_0\in\Sigma$ be computable; for instance $E_0=\min\Sigma$ (Corollary~\ref{cor:IDScomp}). Then $F_\beta(\ell(E_0))=1-\ell(E_0)/\beta(\alpha_*)$ has Turing degree exactly $\mathbf 0''$. If moreover $E_0$ lies in the support of the spectral measure $\mu$ of $\delta_0$ for $H_{v,\alpha_*,1/4}$, then $F_\beta(\ell(E_0))$ is the energy-resolved exponent at $E_0$: for every $\gamma\in(0,\ell_*/2)$ there are real $g_\gamma\in C_c^\infty(\R)$ supported in $\{E:|\ell(E)-\ell(E_0)|<\gamma\}$ with $g_\gamma(H)\delta_0\ne0$, and for every such family and every $p>0$, $\beta^+_{g_\gamma(H)\delta_0}(p)\to F_\beta(\ell(E_0))$ as $\gamma\to0$.
\end{theorem}

\begin{proof}
$\ell(E_0)$ is computable (Theorem~\ref{thm:LEcomp}) and lies in $(0,\beta)$ (Theorem~\ref{thm:nbhd}), so the degree argument of Theorem~\ref{thm:robust} applies. If $E_0\in\supp\mu$, a smooth $g\ge0$ equal to $1$ near $E_0$ and supported in the open set $\{|\ell-\ell(E_0)|<\gamma\}$ has $\|g(H)\delta_0\|^2=\int g^2\,d\mu>0$. For every admissible $g_\gamma$, \cite[Theorem~B(2)]{DMuniv} gives $F_\beta(\ell(E_0)+\gamma)\le\beta^+_{g_\gamma(H)\delta_0}(p)\le F_\beta(\ell(E_0)-\gamma)$.
\end{proof}

By~\cite[Theorem~B(3)]{DMuniv}, the upper local dimension of $\mu$ equals $F_\beta(\ell(E))$ at $\mu$-almost every $E$. This is an almost-everywhere statement; it does not identify the local dimension at a prescribed computable energy.

\begin{theorem}[The set of values]\label{thm:values}
For every $v\in\mathcal U_{\mathrm{even}}$ with computable coefficients, as $\alpha$ ranges over the computable irrationals with eventually odd denominators, $\au(H_{v,\alpha,1/4})$ ranges over exactly the reals in $[0,1]$ that are right-c.e.\ relative to $\mathbf 0'$.
\end{theorem}

\begin{proof}
By Theorem~\ref{thm:nbhd}, $H_{v,\alpha}$ is supercritical and type~I for every irrational $\alpha$, and $1/4\in\Ga$ when the denominators are eventually odd, so the transport law~\cite[Theorem~A]{DMuniv} applies at the phase $1/4$, including the endpoints $\beta\in\{0,\infty\}$.

\emph{Every value is right-c.e.\ relative to $\mathbf 0'$.} For computable $\alpha$, $\ell=\ell_*(v,\alpha)$ is computable and positive, and $\beta(\alpha)$ is right-c.e.\ relative to $\mathbf 0'$ or infinite~\cite[Lemma~10.1]{DMamo}. For rational $r\in(0,1)$, $F_\beta(\ell)<r$ if and only if some rational $s'$ satisfies $\beta<s'<\ell/(1-r)$, a $\mathbf 0'$-c.e.\ condition.

\emph{Every such value occurs.} The values $0$ and $1$ occur for $\alpha=[0;1,2,2,\dots]$ and for $a_1=1$, $a_{n+1}=2^{q_n^2}$. Let $x\in(0,1)$ be right-c.e.\ relative to $\mathbf 0'$ and $\sigma=1/(1-x)>1$. By~\cite[Lemma~10.1]{DMamo} there are computable rationals $u_n$ with $\limsup u_n=\sigma/\log2$; clamp them to $[1,B]$ for an integer $B>\sigma/\log2$. Put $a_1=1$. Given $a_1,\dots,a_n$, let $\alpha^{(n)}=[0;a_1,\dots,a_n,2,2,\dots]$, a computable irrational, uniformly in $n$; compute a rational $c_n$ with $|c_n-\ell_*(v,\alpha^{(n)})|<2^{-n}$ (Theorem~\ref{thm:LEcomp}); and put $a_{n+1}=2^{\max\{1,\lceil u_nc_nq_n\rceil\}}$. All $a_{n+1}$ are even, so all $q_n$ are odd, and $\alpha=\lim\alpha^{(n)}$ is computable, with $|\alpha^{(n)}-\alpha|\le2/q_n^2$. By Lemma~\ref{lem:minonsigma}, $\ell_*(v,a)=\min_{|E|\le K}\ell(a,E)$ with $K$ independent of $a$, and the Lyapunov exponent is jointly continuous in $(a,E)$ at irrational $a$~\cite[Theorem~1]{BourgainJitomirskaya2002}; so $c_n\to\ell:=\ell_*(v,\alpha)>0$. Since $a_{n+1}q_n\le q_{n+1}\le2a_{n+1}q_n$, $\log q_{n+1}/q_n=(\log2)u_nc_n+O(\log q_n/q_n)$, so $\beta(\alpha)=\ell\sigma>\ell$ and $\au=1-\ell/\beta=x$.
\end{proof}

\begin{proposition}[A computable good phase]\label{prop:goodphase}
Let $v$ be a computable real-analytic potential with $H_{v,\alpha_*}$ supercritical and type~I. There is a computable phase $\theta_v$ such that $\theta_v+j\alpha_*\in G^v_{\alpha_*}$ for every $j\in\Z$. Hence $\au(H_{v,\alpha_*,\theta_v})=F_{\beta(\alpha_*)}(\ell_*(v))$ for the packet started at any site; if $\ell_*(v)<\beta(\alpha_*)$, this number has Turing degree exactly $\mathbf 0''$.
\end{proposition}

\begin{proof}
By~\cite[Lemma~5.20(1)]{DMuniv} there is $q_*$, depending on $v$ and $\alpha_*$, such that for convergent denominators $q\ge q_*$ each band centre has exactly one maximum and one minimum per period. Band centres are the ordered simple eigenvalues of a computable Hermitian Floquet matrix~\cite[Lemma~5.3]{DMuniv}, hence computable functions of the phase, and the unique maximizer of a computable continuous function on the circle is computable. So the set $\mathrm{Turn}_q$ of these $2q$ points is computable uniformly in $q\ge q_*$. Let $n_*$ be the first index with $q_{n_*}\ge q_*$; it is a fixed integer, which we do not compute but hard-code. Let $B_n=\{\theta:\dist(q_n\theta,\mathrm{Turn}_{q_n})<3q_n^{-3}\}$, a union of at most $2q_n^2$ arcs with computable endpoints and total measure at most $12q_n^{-2}$. Choose $N\ge n_*$ with $\sum_{n\ge N}|B_n|<1/2$, using the effective tail given by Fibonacci growth, and run the nested dyadic construction of Corollary~\ref{cor:twoexp}(4) with the sets $B_n$ in place of $E_n$. It gives a computable $\theta_v$ outside every $B_n$, $n\ge N$. For a site $j$, $q_n(\theta_v+j\alpha_*)\equiv q_n\theta_v+j\delta_n$ and $|\delta_n|\le1/q_{n+1}\le2^{-q_n}$, so $\dist(q_n(\theta_v+j\alpha_*),\mathrm{Turn}_{q_n})\ge3q_n^{-3}-|j|2^{-q_n}\ge2q_n^{-3}$ for large $n$, which is goodness in the sense of~\cite[Definition~5.21]{DMuniv}. The value follows from~\cite[Theorem~A]{DMuniv}, and the degree as in Theorem~\ref{thm:robust}.
\end{proof}

The program for $\theta_v$ exists, but we do not exhibit it, because $n_*$ is not computed. For even $v$ the phase $1/4$ is explicit.

\subsection{What is fragile}

\begin{theorem}[Fragility]\label{thm:fragile}
Let $v\in\mathcal U_{\mathrm{even}}$ have computable coefficients, and write $X(v,\alpha,\theta)=\au(H_{v,\alpha,\theta})$.
\begin{enumerate}
\item The value $X(v,\alpha_*,\theta)=X(v)$ is the same for every $\theta$ in the full-measure set $G^v_{\alpha_*}$, which is not a neighbourhood of any phase. At every other phase, $X(v)\le X(v,\alpha_*,\theta)\le1-\ell_*(v)/(2\beta(\alpha_*))$, by~\cite[Theorem~A]{DMuniv} and Corollary~\ref{cor:typeI}; we do not know the exact value there.
\item For every irrational $\alpha_0$ and every open interval $J\ni\alpha_0$, $J$ contains computable frequencies $\alpha$ with $X(v,\alpha,1/4)=0$, with $X(v,\alpha,1/4)=1$, and with $X(v,\alpha,1/4)=x$ for any prescribed $x\in(0,1)$ that is right-c.e.\ relative to $\mathbf 0'$. In particular $\alpha\mapsto X(v,\alpha,1/4)$ is discontinuous at every irrational. Moreover $X(v,\alpha,\theta)=0$ for almost every $\alpha$ and every $\theta$, and $X(v,\alpha,\theta)=1$ for a dense $G_\delta$ set of $\alpha$ and every $\theta$.
\end{enumerate}
\end{theorem}

\begin{proof}
1. \cite[Theorem~A]{DMuniv} and Corollary~\ref{cor:typeI}, since $\ell_*\le\beta(\alpha_*)$. The good set has empty interior. For even $v$ the turning points are $0$ and $1/2$~\cite[Lemma~5.22]{DMuniv}, and the bad sets $\{\theta:\dist(q_n\theta,\{0,1/2\})<2q_n^{-3}\}$ are open and $1/q_n$-periodic. So each tail union $\bigcup_{n\ge n_0}$ of them is open and dense, their intersection over $n_0$ is residual, and it is disjoint from $G^v_{\alpha_*}$. A set disjoint from a dense $G_\delta$ set has empty interior.

2. Choose a continued-fraction cylinder $[0;a_1,\dots,a_k]$ contained in $J$ and containing $\alpha_0$. Two consecutive denominators are coprime, so not both even; appending at most two partial quotients equal to $1$ makes the last two denominators odd. Now rerun the construction of Theorem~\ref{thm:values} from this prefix, using the actual denominators of the prefixed frequencies: even partial quotients keep every later denominator odd, and the continuity argument gives $\beta=\ell\sigma$ as before. A finite prefix does change the frequency, and, through it, the later denominators; this is why the construction is rerun rather than a fixed tail being appended, which would in general rescale $\beta$. The endpoints $0$ and $1$ are obtained with tails $2,2,\dots$ and $a_{n+1}=2^{q_n^2}$. Finally $\beta(\alpha)=0$ for almost every $\alpha$ and $\beta(\alpha)=\infty$ for a dense $G_\delta$ set of $\alpha$; at these frequencies the exponent is $0$, respectively $1$, at every phase~\cite[Theorem~A]{DMuniv}.
\end{proof}

So, for the transport laws identified here at the prescribed good phases, the frequency is the only fragile datum, and it is exactly the datum that carries the uncomputability; we have proved no robustness in the phase beyond those good phases: $X=1-\ell_*/\beta$ with $\ell_*$ computable and continuous in the potential, and $\beta(\alpha)$ noncomputable as a function of $\alpha$ and discontinuous everywhere.

\begin{theorem}[A dichotomy]\label{thm:dichotomy}
Let $v$ be a computable even real-analytic potential with $H_{v,\alpha_*}$ supercritical and type~I. Then either $\ell_*(v)>\beta(\alpha_*)$ and $X(v)=0$, or $\ell_*(v)<\beta(\alpha_*)$ and $X(v)$ has Turing degree exactly $\mathbf 0''$. For non-even $v$ in this class the same holds at the phase $\theta_v$ of Proposition~\ref{prop:goodphase}.
\end{theorem}

\begin{proof}
$\ell_*(v)$ is computable and $\beta(\alpha_*)$ is not, so they differ. If $\ell_*>\beta$ the law gives $0$; if $\ell_*<\beta$ the exponent computes $\beta$, as in Theorem~\ref{thm:robust}.
\end{proof}

For instance the coupling $\lambda=8$ gives $\ell_*=\log8>2\log2\ge\beta(\alpha_*)$ and exponent $0$. The dichotomy concerns the supercritical type~I class; for critical potentials and for supercritical potentials that are not of type~I we do not know the exponent at $\alpha_*$.

\begin{theorem}[Exponents that exist]\label{thm:exist}
\begin{enumerate}
\item The limit of a computable sequence of rationals is computable from $\mathbf 0'$. So a quantity that is the limit of finite-time, finite-volume observables computable uniformly along computable times has degree at most $\mathbf 0'$; degree $\mathbf 0''$ requires a $\limsup$, and the $\limsup$s of computable sequences of rationals are exactly the reals right-c.e.\ relative to $\mathbf 0'$~\cite[Lemma~10.1]{DMamo}.
\item For real-analytic $v$ with $H_{v,\alpha}$ supercritical, every irrational $\alpha$, every phase and every $p>0$, the lower transport exponent $\beta^-(p)=\liminf_{T\to\infty}\log M_p(T)/(p\log T)$ is $0$.
\end{enumerate}
\end{theorem}

\begin{proof}
1. This is the limit lemma.

2. Hopping has range one, so $|\langle\delta_j,e^{-itH}\delta_0\rangle|\le(\|H\|t)^{|j|}/|j|!$, and the Abel average gives $\sum_{|j|>T^2}|j|^pa(j,T)\to0$. By~\cite[Theorem~5.9]{DMuniv}, with $N_0=C'q_n$, $C'\ge C_D$, $c_DC'\ge\kappa$, and $T_n=e^{\kappa q_n/(2p+4)}$, $P(N_0,T_n)\le Ce^{-c_DC'q_n}+CT_n^3e^{-\kappa q_n}$. Hence $M_p(T_n)\le N_0^p+T_n^{2p}P(N_0,T_n)+o(1)=O((\log T_n)^p)$, so the liminf is at most $0$. It is at least $0$, because the Abel return probability $a(0,T)$ tends to the sum of the squared atomic spectral weights of $\delta_0$, which is less than one since $\delta_0$ is not an eigenvector.
\end{proof}

Part~2 is known. It follows directly from Jitomirskaya and Mavi~\cite[Theorem~1.7(1)]{JitomirskayaMavi2017}, for piecewise H\"older potentials with continuous positive Lyapunov exponent, and it is in substance in the complex-energy method of Damanik and Tcheremchantsev~\cite{DT2007}; we include the short proof for completeness. Consequently, in the supercritical analytic class, a moment exponent that exists as a limit is $0$, and the other natural existing quantities, namely $\ell(E)$ at computable $E$, $\ell_*$, the density of states, the spectrum and $\max_\Sigma\ell$, are computable at computable data (Theorem~\ref{thm:LEcomp}, Corollary~\ref{cor:IDScomp}). In the examples of this section, uncomputability is carried by $\limsup$-type quantities: upper transport exponents and upper packing dimensions. We do not claim a classification of all observables.

\begin{remark}[Related work on uncomputability]\label{rem:uncompprior}
Undecidability and uncomputability of physical quantities are known in several forms: the spectral gap of Cubitt, Perez-Garcia and Wolf~\cite{CPGW2015}; phase diagrams, including continuous parameters and undecidability on sets of positive measure, by Bausch, Cubitt and Watson~\cite{BCW2021}; and a single uncomputable transition point in a continuous Hamiltonian family, by Purcell, Li and Cubitt~\cite{PLC2024}. We do not claim that robust or continuous-parameter uncomputability is new. Paper A already gave an exponent of degree $\mathbf 0''$ that is robust along a coupling interval. The new point here is that, at the frequency $\alpha_*$ and the phase $1/4$, the degree is stable for the computable members of an explicit neighbourhood in the even analytic potentials, open relative to that class, with the analogous statement for non-even computable potentials at the good phase of Proposition~\ref{prop:goodphase} (Theorem~\ref{thm:dichotomy}); and that the Lyapunov exponent of a computable analytic potential is itself computable, so that in these laws the frequency carries the uncomputability.
\end{remark}

\section{Discussion and open problems}\label{sec:discussion}

\paragraph{The picture.} At a resonant scale a packet moves with the speed allowed by the widest band near its energy and stops where the drift of the local phase stops it. For type~I potentials, at almost every phase the widest band has width about $e^{-\ell q}$; this is the heuristic behind the typical law. At the worst phase it can be as wide as $e^{-\ell q/(2\om)}$, because $2\om$ zeros of the discriminant can form an exponentially tight cluster at one local phase, and such clusters do form when the potential has a hidden period; the zeros of a cluster need not coincide. For the almost Mathieu operator the cluster sits at the closed central gap of even-denominator approximants, and it is visited by a residual set of phases; when $\beta_e=\beta$ is finite and at least $L$, this gives the residual law of Theorem~\ref{thm:C}. For the Aubry--Andr\'e XY chain and the observable classes of Definition~\ref{def:LR}, in the same regime, almost every phase has fixed-site exponent $1-L/\beta$, while the exponent uniform over sites is $1-L/(2\beta)$ at every phase; exceptional phases can have the larger exponent at every fixed site as well, as the phase $1/4$ does in Corollary~\ref{cor:twoexp}. All exponents here are infinite-time quantities of an infinite system; at finite times the dynamics is a sequence of bursts at the resonant times.

\subsection{Open problems}

\paragraph{1. The factor for primitive potentials.} Theorem~\ref{thm:A} charges the price of a cluster of $2\om$ zeros, and Proposition~\ref{prop:NP} shows that the price is paid when the potential has a hidden period $1/s$. Write $v(x)=w(sx)$ with $w$ primitive, as in~\cite[Theorem~5.29]{DMuniv}, so that $H_{v,\alpha,\theta}=H_{w,s\alpha,s\theta}$.

\begin{conjecture}\label{conj:M2}
Let $v=w(s\,\cdot)$ be real analytic with $L>0$ on $\Sigma$ and $w$ primitive. If $\max_\Sigma\ell\le\beta(s\alpha)<\infty$, then $\aunif=\sup_\theta\au(\theta)=1-\ell_*/(2\beta(s\alpha))$.
\end{conjecture}

For the two-harmonic family of Proposition~\ref{prop:twoharm} with $\lambda\ne0$, which is then primitive and has acceleration two exactly where $\ell=\log|\mu|$, assume that $\om=2$ somewhere on $\Sigma$, so that $\ell_*=\log|\mu|$, and that $\ell^+\le\beta<\infty$. Then Theorem~\ref{thm:A} gives the upper bound $1-\log|\mu|/(4\beta)$, while Conjecture~\ref{conj:M2} predicts $1-\log|\mu|/(2\beta)$. We do not know a mechanism that would force clusters of more than two zeros for a primitive potential, nor a proof that they cannot occur at the scales that matter.

\begin{remark}[Numerical observations; not used in any proof]\label{rem:numerics}
We computed a linearized bandwidth proxy at single rationals $p/q$ near the golden mean, $q=8,13,21,34,55$: the quantity $e_q/\ell(E_*)$, where $e_q=-q^{-1}\log\max_{x,b}4/|\partial_Et_q(x,E_b(x))|$ over phases and band centres, and $\ell(E_*)$ is a Birkhoff estimate of the Lyapunov exponent at the energy $E_*$ of the widest band. Ranges registered before the runs were $[0.40,0.60]$ for primitive potentials and for the almost Mathieu control, and $[0.18,0.32]$ for the non-primitive control $4\cos4\pi x$ at even $q$. Three primitive potentials with estimated acceleration two or three at $E_*$ gave values near $1/2$ rather than $1/(2\om)$: $\cos2\pi x+5\cos4\pi x$ gave $0.50$--$0.58$, a three-harmonic potential without symmetry gave $0.52$--$0.55$, and the odd-harmonic $\cos2\pi x+5\cos6\pi x$ gave $0.50$--$0.59$ for $q\le34$ and $0.630$ at $q=55$, a miss. There were further misses: the almost Mathieu control gave $0.629$, $0.709$, $0.659$ at $q=8,13,21$; the non-primitive control gave $0.683$ and $0.646$ at odd $q=13,21$, and $0.388$ at even $q=34$, against $0.268$ at $q=8$. The proxy is not the exact band width, the maximization over phases is a finite search, and the acceleration estimates are not certificates. These are finite-size observations, not evidence of an asymptotic law.
\end{remark}

\paragraph{2. The matching lower bound for general potentials.} Theorem~\ref{thm:Wv} reduces the lower bound $1-\ell_0/(2\beta)$ to the existence of a wide band, $(\mathrm W_v)$, near an energy with $\ell\approx\ell_0$. A sufficient condition is a small gap between band centres: if $E_b(x)$ and $E_{b+1}(x)$ are roots of $t_q(x,\cdot)$ near $E'$ with $E_{b+1}(x)-E_b(x)\le e^{-(\ell(E')/2-4\xi)q}$, then the maximum principle for $t_q(x,z)/((z-E_b)(z-E_{b+1}))$ on a disk of radius $e^{-\xi q}$, with~\cite[Corollary~5.6]{DMuniv}, bounds $|\partial_Et_q(x,E_b)|$ by $e^{(\ell(E')/2+7\xi+o(1))q}$. We do not know whether this condition is also necessary. It is proved here only at $E=0$ for even $q$ and odd-harmonic potentials, through exactly closed gaps; for general type~I potentials, and at the minimizer of $\ell$ when $\ell(0)>\ell_*$, it is open.

\paragraph{3. The regime $0<\beta<\ell$.} For supercritical type~I operators and $0<\beta<\ell_*$, the exponent is $0$ at almost every phase~\cite[Theorem~A]{DMuniv}, while Corollary~\ref{cor:typeI} only gives the positive uniform bound $\beta/(\beta+\ell_*)$. For general supercritical analytic operators, Theorem~\ref{thm:A} gives the uniform bound $\max_{E\in\Sigma}f(\beta;\ell(E),\om(E))$, which involves $(2\om-1)\beta/((2\om-1)\beta+\ell)$ at energies with $\beta<\ell$; the typical value is not $0$ in general there, since a hidden period replaces $\beta(\alpha)$ by $\beta(s\alpha)$~\cite[Theorem~C]{DMuniv}, which can be larger. For the almost Mathieu operator with $L/2<\beta_e=\beta<L$ the worst exponent lies between $1-L/(2\beta)$ and $\beta/(\beta+L)$. Whether $\au=0$ at every phase for $0<\beta\le\ell_*/(2\om^+)$ is open. The obstruction is the error term $e^{(\ell-b)q}$ of the $q$-step identity, which raises the sublevel threshold in Proposition~\ref{prop:Comega} when $b<\ell$.

\paragraph{4. Odd scales for the almost Mathieu operator.} Theorem~\ref{thm:quarter} uses the closed gaps at even denominators. For odd $q$ every gap of the rational approximants is open~\cite{vanMouche1989}, so no closure argument is available, and a quantitative bound on one gap is needed. Fix an odd $q\ge3$ and $p$ coprime to $q$, and write $P^\lambda$ for the Chambers polynomial of coupling $\lambda$. Its critical points $c_1<\dots<c_{q-1}$ are real and simple, and $\mathrm{crit}_k(\lambda)=|P^\lambda(c_k)|-2\lambda^q=\min_\varphi|\Delta_\varphi(c_k)|\ge2$, because every critical value of a periodic discriminant has modulus at least two. Put $\mathrm{ex}_k(\lambda)=\mathrm{crit}_k(\lambda)-2$, and let $\mathrm{ex}_{\mathrm{ext}}$ refer to the two extreme gap labels $m=\pm(q-1)/2$, where $k\equiv mp\pmod q$.

\begin{lemma}[Duality of the excess]\label{lem:X}
$\mathrm{ex}_k(1/\lambda)=\lambda^{-q}\,\mathrm{ex}_k(\lambda)$ for every $k$ and every $\lambda>0$.
\end{lemma}

\begin{proof}
On the cyclic $q$-site space let $H_\lambda(a,b)=e^{ia}S+e^{-ia}S^*+2\lambda\operatorname{diag}\{\cos(2\pi np/q+b)\}$. The discrete Fourier transform with entries $q^{-1/2}e^{2\pi ipns/q}$ exchanges shifts and multiplications, and conjugates $H_\lambda(a,b)$ to $\lambda H_{1/\lambda}(b,a)$, up to orientation signs in $a,b$. So $\det(E-H_\lambda(a,b))=\lambda^q\det(E/\lambda-H_{1/\lambda}(b,a))$. Averaging over both phases removes the phase-dependent harmonics of Chambers' formula and gives $P^\lambda(E)=\lambda^qP^{1/\lambda}(E/\lambda)$. So the critical points scale by $\lambda$, preserving their order, and $\mathrm{crit}_k(\lambda)-2=\lambda^q(\mathrm{crit}_k(1/\lambda)-2)$.
\end{proof}

\begin{proposition}[The strong-coupling coefficient]\label{prop:X2}
For odd $q\ge3$ and $p$ coprime to $q$,
\[
\lim_{\lambda\to\infty}\frac{\mathrm{ex}_{\mathrm{ext}}(\lambda)}{\lambda}=\lim_{\mu\to0}\frac{\mathrm{ex}_{\mathrm{ext}}(\mu)}{\mu^{q-1}}=C_q:=4\sin^2\Bigl(\frac{\pi mp}q\Bigr),\qquad m=\frac{q-1}2 .
\]
In particular $0<C_q\le4$.
\end{proposition}

\begin{proof}
The first equality is Lemma~\ref{lem:X}. At weak coupling $\mu$, the Bloch fibre of the $q$-periodic operator in the plane-wave basis indexed by $s=0,\dots,q-1$ is a $q$-cycle with diagonal $\epsilon_s=2\cos(k+2\pi sp/q)$ and hopping of size $\mu$. The gap with label $m$ sits at the crossing $\epsilon_0=\epsilon_m$, at $k\equiv-\pi mp/q\pmod\pi$ and $E_0=2\cos k$. For fixed $p,q$ the two crossing states are isolated from the others, and a Schur complement onto them is finite-dimensional analytic perturbation theory. By reflection the two diagonal self-energies agree, and the off-diagonal coupling comes from the unique path of length $m$, the other path having length $m+1$:
\[
|t_{\mathrm{eff}}|=\frac{\mu^m}{\prod_{s=1}^{m-1}|E_0-\epsilon_s|}\bigl(1+O_{p,q}(\mu)\bigr).
\]
Using $|E_0-\epsilon_s|=4|\sin(k+\pi sp/q)\sin(\pi sp/q)|$ and the identity $\prod_{s=1}^{m}|\sin(\pi sp/q)|=\sqrt q\,2^{-m}$, which follows from $\prod_{s=1}^{q-1}|\sin(\pi sp/q)|=q2^{-(q-1)}$ and the symmetry $s\mapsto q-s$, one gets $\prod_{s=1}^{m-1}|E_0-\epsilon_s|=q/(4\sin^2(\pi mp/q))$. The half-gap is $|t_{\mathrm{eff}}|$, and the free discriminant $2T_q(E/2)$ has second derivative of modulus $q^2/(2\sin^2(\pi mp/q))$ at the corresponding critical point. A Taylor expansion gives $\mathrm{ex}_{\mathrm{ext}}(\mu)=\frac12|P''|\,|t_{\mathrm{eff}}|^2(1+O_{p,q}(\mu))=4\sin^2(\pi mp/q)\mu^{q-1}(1+O_{p,q}(\mu))$; the phase-union term $2\mu^q$ is of higher order.
\end{proof}

The remainders depend on $p$ and $q$, so Proposition~\ref{prop:X2} gives no control uniform in $q$. Sharp exponential asymptotics of spectral gaps are known for Diophantine frequencies in the subcritical regime~\cite{LYZZ2017}, and upper bounds on gaps at Liouville frequencies in a regime of weak dual coupling~\cite[Theorem~1.3]{LiuShi2019}; as stated, neither gives the bound on the extreme gaps of the rational approximants that is needed here.

\begin{conjecture}[Mon]\label{conj:mon}
For every odd $q$ and every $p$ coprime to $q$, the function $\lambda\mapsto\mathrm{ex}_{\mathrm{ext}}(\lambda)/\lambda$ is nondecreasing on $(0,\infty)$.
\end{conjecture}

\begin{proposition}[Conditional exact law]\label{prop:X4}
If Conjecture~\ref{conj:mon} holds, then $\aunif=1-\log\lambda/(2\beta(\alpha))$ for the almost Mathieu operator with $\lambda>1$ and every irrational $\alpha$ with $\log\lambda\le\beta(\alpha)<\infty$, without any parity condition.
\end{proposition}

\begin{proof}
Under the conjecture, Proposition~\ref{prop:X2} gives $\mathrm{ex}_{\mathrm{ext}}(\lambda)\le C_q\lambda\le4\lambda$, so the critical value $D=\mathrm{crit}_{\mathrm{ext}}(\lambda)\le4\lambda+2$ is bounded uniformly in $p,q$. Let $\varphi$ be a phase at which $|\Delta_\varphi(c)|=D$ for the extreme critical point $c$, and let $e$ be a point of an adjacent band with $|\Delta_\varphi(e)|\le1$; it lies within $e^{-\xi q}$ of $\Sigma$ when $b_n$ is bounded below. On a $1/q$-neighbourhood of $e$, \cite[Corollary~5.6]{DMuniv} gives $|\Delta_\varphi'|\le2qe^{(L+\eps)q}$ and $|\Delta_\varphi''|\le M=8q^2e^{(L+\eps)q}$. Moving from $e$ towards $c$ by $t=|\Delta_\varphi'(e)|/M\le1/(4q)$, which does not pass $c$ because $\Delta'_\varphi(c)=0$, Taylor's formula on a segment where the signed discriminant is at most $D$ gives $|\Delta_\varphi'(e)|^2\le2M(D+1)$. So $(\mathrm W_v)$ holds at every large odd scale, with $\eta=1$ and $\ell_0=L$, and Theorem~\ref{thm:Wv}, together with Theorem~\ref{thm:quarter} at even scales, gives $\aunif\ge1-L/(2\beta)$. Corollary~\ref{cor:typeI} gives the upper bound.
\end{proof}

\begin{remark}[Numerical observations; not used in any proof]\label{rem:monnum}
For $p/q=8/13$ and $13/21$, the discrete second differences of $t\mapsto\log\mathrm{ex}_{\mathrm{ext}}(e^t)$ on a grid of step $1/4$ in $[-3,4]$ were all negative. For twelve rationals with $5\le q\le21$, the ratio $\mathrm{ex}_{\mathrm{ext}}(\lambda)/(C_q\lambda)$ was below one at $\lambda=1.2,2,4$, and $\mathrm{ex}_{\mathrm{ext}}(10^5)/10^5$ agreed with $C_q$ to about $2\cdot10^{-5}$, the difference being of the size of the next term. Global concavity of $t\mapsto\log\mathrm{ex}_{\mathrm{ext}}(e^t)$, together with Proposition~\ref{prop:X2}, would imply Conjecture~\ref{conj:mon}. A grid does not prove concavity.
\end{remark}

\paragraph{5. Weak separation.} The typical exponent beyond type~I is known only conditionally on the weak separation hypothesis WXS for primitive potentials with acceleration at least two~\cite[Definition~7.1 and Theorem~D]{DMuniv}. It remains open, and so does the typical exponent of the two-harmonic family.

\paragraph{6. How much of the packet moves at the worst phase.} Theorem~\ref{thm:quarter} produces a polynomially small tail at the exceptional radius, whose exponent is $1-L/(2\beta)$ when $\beta_e=\beta$ is finite and at least $L$. For the almost Mathieu operator with $h>L$, \cite[Theorem~B]{DMamo} gives at every phase a fixed Abel-averaged fraction of the packet beyond $Te^{-hq_n}$, for $e^{hq_n}\le T\le|\delta_n|^{-1}$ on the resonant scales with $b_n\ge c>h$; at $T=|\delta_n|^{-1}$ this is the smaller, typical radius. Whether a fixed fraction reaches the larger exceptional radius of Theorem~\ref{thm:quarter}, and what the moment exponents are at the exceptional phases, is open.

\subsection*{How this work was done}

The authors developed the results and proofs with assistance from AI models from more than one developer. They take responsibility for every claim. The authors chose the question, directed the work and decided which lines to pursue; the models did most of the exploration, derivation and drafting.

The proof was developed over several long working sessions, each building on the previous ones and recording its proofs, failed routes and self-checks. Each step was then checked by a separate session that had not seen the material, with instructions to refute it and to read the primary sources it relies on. The repairs found in these checks are incorporated here. Priority and literature searches were AI-assisted. The assembled paper was then audited as a whole by another session that had not seen it, also instructed to refute; the repairs made in answer to that audit and to each later check were checked again by a fresh session, including one session that compared every summary sentence with the result it summarizes, and the final small repair uses the wording that the last check proposed. No computation is used in any proof, except that the elementary numerical constants of Theorem~\ref{thm:nbhd} were also certified with exact rational arithmetic; the numerical remarks are not used in any proof.


\bibliographystyle{plainurl}
\bibliography{references}
\end{document}
